% International Economics — Economics Upper Sixth — Cours signé OnBuch+
% Official MINESEC syllabus. Compiler avec : tectonic ECO06-international-economics.tex
\newcommand{\DOCMATIERE}{Economics}
\newcommand{\DOCNIVEAU}{Upper Sixth}
\newcommand{\DOCMODULE}{Upper Sixth — Economics}
\newcommand{\DOCLECON}{6}
\newcommand{\DOCTITRE}{International Economics}
% OnBuch+ — préambule « cours signés » ENRICHI (compilé avec tectonic / XeLaTeX)
% Système visuel « Pop » d'OnBuch+ : fond crème (jamais blanc), bordures encre
% épaisses, ombres « dures » décalées, titres Archivo Black en capitales.
% Version MATHÉMATIQUES : graphes (pgfplots/tikz), boîtes pédagogiques
% (définition, propriété, méthode, attention, le savais-tu, exercice résolu,
% exercice, TP, bilan), page de garde et signature OnBuch+ ancrée en bas.
%
% Macros attendues dans main.tex AVANT \input{preamble} :
%   \DOCMATIERE  \DOCNIVEAU  \DOCMODULE  \DOCLECON (numéro)  \DOCTITRE
\documentclass[11pt]{article}

\usepackage{fontspec}
\usepackage[english]{babel}
\usepackage[a4paper,margin=2cm,top=3.4cm,bottom=2.8cm,headheight=1.4cm,footskip=1.3cm]{geometry}
\usepackage{amsmath,amssymb,amsfonts}
\usepackage{mathtools} % \xmapsto, \coloneqq... souvent utilisés par les modèles
\usepackage{cancel} % simplifications barrées (bilans de forces)
% \abs{z} (module, valeur absolue) et \norm{u} (norme), utilisés par les modèles
\providecommand{\abs}{}\let\abs\relax\DeclarePairedDelimiter{\abs}{\lvert}{\rvert}
\providecommand{\norm}{}\let\norm\relax\DeclarePairedDelimiter{\norm}{\lVert}{\rVert}
\usepackage[table]{xcolor} % [table] : fournit aussi \rowcolors (alternance de lignes)
\usepackage{pagecolor}
\usepackage{tikz}
\usetikzlibrary{shapes.symbols,circuits.logic.US,circuits.logic.IEC,arrows.meta,positioning,calc,shapes.geometric,shapes.misc,shapes.arrows,decorations.pathmorphing,decorations.pathreplacing,decorations.markings,patterns,shadows,mindmap,trees,fit,backgrounds,matrix,intersections,angles,quotes}
\usepackage{pgfplots}
\usepackage[version=4]{mhchem} % équations nucléaires \ce{^{235}_{92}U}
\usepackage[european,siunitx]{circuitikz} % schémas électriques
\pgfplotsset{compat=1.18}
\usepgfplotslibrary{fillbetween,groupplots,polar,statistics,dateplot} % aires entre courbes (calcul intégral)
\usepackage{enumitem}
\usepackage{fancyhdr}
% [most] charge toutes les bibliothèques tcolorbox (listings, minted, raster,
% documentation...) et gonfle inutilement la mémoire TeX sur un document long
% (~300 boîtes) : on ne charge que ce qui est réellement utilisé.
\usepackage[skins,breakable]{tcolorbox}
\usepackage{graphicx}
\usepackage{colortbl}
\usepackage{booktabs}
\usepackage{tabularx}
\usepackage{diagbox} % case d'en-tête diagonale des tableaux à double entrée
% Colonnes extensibles centrée / gauche / droite (C, L, R), souvent utilisées par les modèles
\newcolumntype{C}{>{\centering\arraybackslash}X}
\newcolumntype{L}{>{\raggedright\arraybackslash}X}
\newcolumntype{R}{>{\raggedleft\arraybackslash}X}
\usepackage{array}
\usepackage{multicol}
\usepackage{siunitx}
% parse-numbers=false : accepte aussi \SI{2,5 \times 10^{-3}}{...}
\sisetup{locale=UK,output-decimal-marker={.},group-minimum-digits=4,parse-numbers=false}
\DeclareSIUnit\ohm{\ensuremath{\symup{\Omega}}} % U+2126 absent de la police
\DeclareSIUnit\division{div} % réglages d'oscilloscope (ms/div, V/div)
\DeclareSIUnit\year{yr}\DeclareSIUnit\annum{yr}\DeclareSIUnit\Ma{Ma}\DeclareSIUnit\tr{tr} % tours (vitesse de rotation, ex. \SI{1500}{\tr\per\minute})
\usepackage{qrcode}
% \degree : commande fréquemment utilisée par les modèles pour un angle ou une
% température en dehors de \SI{}{} (non fournie par les paquets chargés).
\providecommand{\degree}{^{\circ}}
% \qed (fin de démonstration), utilisé par les modèles sans amsthm
\providecommand{\qed}{\hfill$\square$}
% \barre : double barre des valeurs interdites dans un tableau de variations (tkz-tab)
\providecommand{\barre}{\ensuremath{\|}}
% environnement proof (amsthm non chargé) utilisé par les modèles
\ifdefined\proof\else\newenvironment{proof}[1][Proof]{\par\noindent\textit{#1.}\ }{\qed\par}\fi
\usepackage{lastpage}
\usepackage{hyperref}
\hypersetup{hidelinks}

% ============================================================
% Palette « Pop » OnBuch+ — mêmes hex que la classe P (plus_ui.dart)
% ============================================================
\definecolor{poporange}{HTML}{F59321}
\definecolor{popdark}{HTML}{E07A0C}
\definecolor{popcream}{HTML}{FFF6E8}
\definecolor{popink}{HTML}{1C1714}
\definecolor{poppurple}{HTML}{7A5AE0}
\definecolor{popgreen}{HTML}{2E9E6A}
\definecolor{poppink}{HTML}{E0457B}
\definecolor{popblue}{HTML}{2F7DE1}
\definecolor{popgold}{HTML}{C98A12}
\definecolor{popmuted}{HTML}{8A7F76}
\definecolor{popline}{HTML}{EADFD2}
\definecolor{popgray}{HTML}{8A7F76}
\colorlet{popgrey}{popgray}
% Alias tolérants (le modèle nomme parfois les couleurs différemment)
\colorlet{popwhite}{white}
\colorlet{popblack}{popink}
\colorlet{popred}{poppink}
\colorlet{popyellow}{popgold}
% Teintes claires pour fonds de tableaux / zones
\colorlet{popblueL}{popblue!10!white}
\colorlet{poporangeL}{poporange!14!white}
\colorlet{popgreenL}{popgreen!12!white}
\colorlet{poppurpleL}{poppurple!12!white}
\colorlet{poppinkL}{poppink!12!white}
\colorlet{popgoldL}{popgold!14!white}

\pagecolor{popcream}

% ============================================================
% Typographie
% ============================================================
\defaultfontfeatures{Ligatures=TeX}
\setmainfont{PlusJakartaSans-Medium.ttf}[Path=../../../fonts/,BoldFont=PlusJakartaSans-Bold.ttf]
% Maths en sans-serif (Fira Math) avec chiffres et lettres droites de Plus
% Jakarta, pour que les formules se fondent dans le texte.
\usepackage[math-style=ISO,bold-style=ISO]{unicode-math}
\setmathfont{FiraMath-Regular.otf}[Path=../../../fonts/]
\setmathfont{PlusJakartaSans-Medium.ttf}[Path=../../../fonts/,range={up/num,up/{Latin,latin},\mathup}]
% (icomma retiré : décimales avec point en anglais)
% Caractères mathématiques Unicode absents des polices de texte (Plus Jakarta,
% Archivo Black) : rendus par leur équivalent mathématique, en texte comme en
% formule — sinon ils s'affichent en carré vide (ex. « Arithmétique dans ℤ »).
\usepackage{newunicodechar}
\newunicodechar{ℝ}{\ensuremath{\mathbb{R}}}
\newunicodechar{ℤ}{\ensuremath{\mathbb{Z}}}
\newunicodechar{ℂ}{\ensuremath{\mathbb{C}}}
\newunicodechar{ℕ}{\ensuremath{\mathbb{N}}}
\newunicodechar{ℚ}{\ensuremath{\mathbb{Q}}}
\newunicodechar{𝔻}{\ensuremath{\mathbb{D}}}
\newunicodechar{α}{\ensuremath{\alpha}}
\newunicodechar{β}{\ensuremath{\beta}}
\newunicodechar{γ}{\ensuremath{\gamma}}
\newunicodechar{δ}{\ensuremath{\delta}}
\newunicodechar{ε}{\ensuremath{\varepsilon}}
\newunicodechar{θ}{\ensuremath{\theta}}
\newunicodechar{λ}{\ensuremath{\lambda}}
\newunicodechar{μ}{\ensuremath{\mu}}
\newunicodechar{σ}{\ensuremath{\sigma}}
\newunicodechar{τ}{\ensuremath{\tau}}
\newunicodechar{φ}{\ensuremath{\varphi}}
\newunicodechar{ψ}{\ensuremath{\psi}}
\newunicodechar{ω}{\ensuremath{\omega}}
\newunicodechar{Σ}{\ensuremath{\Sigma}}
% majuscules produites par \MakeUppercase dans les titres (α-aminés -> Α)
\newunicodechar{Α}{\ensuremath{\alpha}}
\newunicodechar{Β}{\ensuremath{\beta}}
\newunicodechar{Γ}{\ensuremath{\Gamma}}
\newunicodechar{Δ}{\ensuremath{\Delta}}
\newunicodechar{Ε}{\ensuremath{\varepsilon}}
\newunicodechar{Θ}{\ensuremath{\Theta}}
\newunicodechar{Λ}{\ensuremath{\Lambda}}
\newunicodechar{Μ}{\ensuremath{\mu}}
\newunicodechar{Τ}{\ensuremath{\tau}}
\newunicodechar{Φ}{\ensuremath{\Phi}}
\newunicodechar{Ψ}{\ensuremath{\Psi}}
\newunicodechar{Ω}{\ensuremath{\Omega}}
\newunicodechar{↦}{\ensuremath{\mapsto}}
\newunicodechar{∈}{\ensuremath{\in}}
\newunicodechar{∉}{\ensuremath{\notin}}
\newunicodechar{∧}{\ensuremath{\wedge}}
\newunicodechar{⇔}{\ensuremath{\Leftrightarrow}}
\newunicodechar{⟺}{\ensuremath{\Longleftrightarrow}}
\newunicodechar{⇒}{\ensuremath{\Rightarrow}}
\newunicodechar{⟹}{\ensuremath{\Longrightarrow}}
\newunicodechar{∘}{\ensuremath{\circ}}
\newunicodechar{∪}{\ensuremath{\cup}}
\newunicodechar{∩}{\ensuremath{\cap}}
\newunicodechar{≡}{\ensuremath{\equiv}}
\newunicodechar{⊂}{\ensuremath{\subset}}
\newunicodechar{⊥}{\ensuremath{\perp}}
\newunicodechar{∃}{\ensuremath{\exists}}
\newunicodechar{∀}{\ensuremath{\forall}}
\newunicodechar{∛}{\ensuremath{\sqrt[3]{\,}}}
\newunicodechar{∜}{\ensuremath{\sqrt[4]{\,}}}
\newunicodechar{⃗}{\ensuremath{{}^{\to}}}
\newunicodechar{ⁿ}{\ifmmode^{n}\else\textsuperscript{\ensuremath{n}}\fi}
\newunicodechar{ˣ}{\ifmmode^{x}\else\textsuperscript{\ensuremath{x}}\fi}
\newunicodechar{ᵇ}{\ifmmode^{b}\else\textsuperscript{\ensuremath{b}}\fi}
\newunicodechar{ʳ}{\ifmmode^{r}\else\textsuperscript{\ensuremath{r}}\fi}
\newunicodechar{ᵘ}{\ifmmode^{u}\else\textsuperscript{\ensuremath{u}}\fi}
\newunicodechar{ʸ}{\ifmmode^{y}\else\textsuperscript{\ensuremath{y}}\fi}
\newunicodechar{ᵃ}{\ifmmode^{a}\else\textsuperscript{\ensuremath{a}}\fi}
\newunicodechar{ᵉ}{\ifmmode^{e}\else\textsuperscript{\ensuremath{e}}\fi}
\newunicodechar{ᵏ}{\ifmmode^{k}\else\textsuperscript{\ensuremath{k}}\fi}
\newunicodechar{ᵖ}{\ifmmode^{p}\else\textsuperscript{\ensuremath{p}}\fi}
\newunicodechar{ᴺ}{\ifmmode^{N}\else\textsuperscript{\ensuremath{N}}\fi}
\newunicodechar{ᶜ}{\ifmmode^{c}\else\textsuperscript{\ensuremath{c}}\fi}
\newunicodechar{ʰ}{\ifmmode^{h}\else\textsuperscript{\ensuremath{h}}\fi}
\newunicodechar{ᵠ}{\ifmmode^{q}\else\textsuperscript{\ensuremath{q}}\fi}
\newunicodechar{ₙ}{\ifmmode_{n}\else\textsubscript{\ensuremath{n}}\fi}
\newunicodechar{ₐ}{\ifmmode_{a}\else\textsubscript{\ensuremath{a}}\fi}
\newunicodechar{ᵢ}{\ifmmode_{i}\else\textsubscript{\ensuremath{i}}\fi}
\newunicodechar{ᵣ}{\ifmmode_{r}\else\textsubscript{\ensuremath{r}}\fi}
\newunicodechar{ₖ}{\ifmmode_{k}\else\textsubscript{\ensuremath{k}}\fi}
\newunicodechar{ᵦ}{\ifmmode_{\beta}\else\textsubscript{\ensuremath{\beta}}\fi}
\newunicodechar{ₓ}{\ifmmode_{x}\else\textsubscript{\ensuremath{x}}\fi}
\newfontfamily\popdisplay{ArchivoBlack-Regular.ttf}[Path=../../../fonts/]
\newfontfamily\poplogo{SpaceGrotesk-Bold.ttf}[Path=../../../fonts/]
\newfontfamily\popsemi{PlusJakartaSans-SemiBold.ttf}[Path=../../../fonts/]
\newfontfamily\popxbold{PlusJakartaSans-ExtraBold.ttf}[Path=../../../fonts/]
\newfontfamily\popxboldls{PlusJakartaSans-ExtraBold.ttf}[Path=../../../fonts/,LetterSpace=3.0]

\setlist[enumerate]{leftmargin=1.7em,itemsep=5pt,topsep=4pt}
\setlist[itemize]{leftmargin=1.7em,itemsep=4pt,topsep=4pt,label={\color{poporange}\textbullet}}
\setlength{\parindent}{0pt}
\setlength{\parskip}{5pt}
\renewcommand{\arraystretch}{1.25}

% pgfplots : style Pop par défaut
\pgfplotsset{
  every axis/.append style={axis line style={popink,line width=1pt},tick style={popink},
    grid style={popline,line width=0.5pt},label style={font=\small},tick label style={font=\footnotesize}},
  popcurve/.style={poporange,line width=1.8pt,no markers,smooth},
}
% Même style disponible dans un \draw TikZ simple (hors pgfplots)
\tikzset{popcurve/.style={poporange,line width=1.8pt}}
\tikzset{popgrid/.style={popline,very thin}}
\colorlet{popcurve}{poporange} % le nom du style est parfois utilisé comme couleur

% ============================================================
% Pill / squiggle
% ============================================================
\newcommand{\poppill}[3]{%
  \tikz[baseline=-0.55ex]\node[draw=popink,line width=1.2pt,rounded corners=2.6pt,%
    fill=#2,inner xsep=6.5pt,inner ysep=3pt]{%
    \popxboldls\fontsize{7}{7}\selectfont\color{#3}\MakeUppercase{#1}};%
}
\newcommand{\popsquiggle}[1]{%
  \tikz[baseline=-0.4ex]{
    \draw[#1,line width=2.6pt,line cap=round]
      plot [smooth] coordinates {(0,0.09) (0.34,0.01) (0.68,0.21) (1.02,0.01) (1.36,0.10)};
  }%
}
% Mot-clé surligné (marqueur orange sous le texte)
\newcommand{\cle}[1]{{\popxbold\color{popdark}#1}}

% ============================================================
% En-tête / pied
% ============================================================
\pagestyle{fancy}
\fancyhf{}
\renewcommand{\headrulewidth}{0pt}
\renewcommand{\footrulewidth}{0pt}
\lhead{\raisebox{-0.15ex}{\poplogo\fontsize{13}{13}\selectfont\color{popink}OnBuch\color{poporange}+}}
\rhead{\poppill{\DOCMATIERE\ \textbullet\ \DOCNIVEAU\ \textbullet\ Chapter \DOCLECON}{white}{popink}}
\lfoot{\popxbold\fontsize{7.6}{7.6}\selectfont\color{popmuted}\MakeUppercase{Signed course OnBuch+}\ \textbullet\ {\popsemi\DOCTITRE}}
\rfoot{\tikz[baseline=-0.6ex]\node[rounded corners=8.5pt,fill=popink,minimum height=17pt,minimum width=17pt,inner xsep=5pt,inner sep=0]{\popxbold\fontsize{8}{8}\selectfont\color{popcream}\thepage\,/\,\pageref*{LastPage}};}

% ============================================================
% Titres
% ============================================================
\newcommand{\chaptitre}[2]{%
  \begin{tcolorbox}[enhanced,colback=poporange,colframe=popink,boxrule=2.6pt,arc=11pt,
      drop shadow=popink,left=15pt,right=15pt,top=12pt,bottom=11pt,before skip=3pt,after skip=18pt]
    \poppill{#2}{popink}{popcream}\par\vspace{9pt}
    {\raggedright\hyphenpenalty=10000\exhyphenpenalty=10000\popdisplay\fontsize{22}{24}\selectfont\color{popcream}\MakeUppercase{#1}\par}
  \end{tcolorbox}
}
% Section numérotée : pastille orange avec le numéro + titre Archivo
\newcounter{popsec}
\newcommand{\coursec}[1]{%
  \stepcounter{popsec}\setcounter{popsub}{0}%
  \par\vspace{16pt}\noindent
  \begin{minipage}{\linewidth}
  \tikz[baseline=-0.7ex]\node[draw=popink,line width=1.6pt,fill=poporange,rounded corners=4pt,minimum size=22pt,inner sep=0]{\popdisplay\fontsize{12}{12}\selectfont\color{popcream}\thepopsec};%
  \hspace{8pt}{\popdisplay\fontsize{15}{17}\selectfont\color{popink}\MakeUppercase{#1}}\par
  \vspace{4pt}\noindent{\color{poporange}\rule{36pt}{3.6pt}}
  \end{minipage}\par\nopagebreak\vspace{8pt}
}
\newcounter{popsub}
\newcommand{\courssub}[1]{%
  \stepcounter{popsub}%
  \par\vspace{9pt}\noindent{\popxbold\fontsize{11.5}{13}\selectfont\color{popink}{\color{poporange}\thepopsec.\thepopsub}\hspace{6pt}#1}\par\nopagebreak\vspace{4pt}
}

% ============================================================
% Boîtes pédagogiques — même recette : fond blanc, bordure épaisse colorée,
% ombre dure encre, badge plein. Titre optionnel [#1].
% ============================================================
\tcbset{popbase/.style={breakable,enhanced,colback=white,boxrule=2.3pt,arc=9pt,
  shadow={3pt}{-3pt}{0pt}{popink},left=14pt,right=14pt,top=11pt,bottom=11pt,before skip=11pt,after skip=15pt,
  fontupper=\popsemi\fontsize{10.6}{14.5}\selectfont\color{popink},
  fontlower=\popsemi\fontsize{10.6}{14.5}\selectfont\color{popink}}}
% Coche dessinée (le glyphe ✓ n'existe pas dans les polices du document)
\newcommand{\popcheck}[1]{\tikz[baseline=-0.5ex]\draw[#1,line width=1.6pt,line cap=round,line join=round](-0.13,0.0)--(-0.04,-0.09)--(0.13,0.1);}
\providecommand{\coche}{\popcheck{popgreen}}
\providecommand{\resultat}[1]{\par\smallskip\noindent\fcolorbox{popgreen}{popgreenL}{\parbox{\dimexpr\linewidth-2\fboxsep-2\fboxrule\relax}{\textbf{Result.} #1}}\par\smallskip}
\newcommand{\popboxhead}[3]{\poppill{#1}{#2}{white}\ifx\relax#3\relax\else\hspace{6pt}{\popxbold\fontsize{10.6}{12}\selectfont\color{popink}#3}\fi\par\vspace{7pt}}

\newtcolorbox{aretenir}[1][]{popbase,colframe=popgreen,before upper={\popboxhead{Key points}{popgreen}{#1}}}
\newtcolorbox{exemplebox}[1][]{popbase,colframe=poporange,before upper={\popboxhead{Exemple}{poporange}{#1}}}
\newtcolorbox{definition}[1][]{popbase,colframe=popblue,before upper={\popboxhead{Definition}{popblue}{#1}}}
\newtcolorbox{propriete}[1][]{popbase,colframe=poppurple,before upper={\popboxhead{Property}{poppurple}{#1}}}
\newtcolorbox{methode}[1][]{popbase,colframe=popgold,before upper={\popboxhead{Method}{popgold}{#1}}}
\newtcolorbox{attention}[1][]{popbase,colframe=poppink,before upper={\popboxhead{Warning}{poppink}{#1}}}
\newtcolorbox{experience}[1][]{popbase,colframe=popblue,colback=popblueL,before upper={\popboxhead{Experiment}{popblue}{#1}}}
% « Le savais-tu ? » : carte violette pleine (contexte, culture, Cameroun)
\newtcolorbox{savaistu}[1][]{popbase,colback=poppurple,colframe=popink,
  fontupper=\popsemi\fontsize{10.4}{14.2}\selectfont\color{white},
  before upper={\poppill{Did you know?}{white}{poppurple}\ifx\relax#1\relax\else\hspace{6pt}{\popxbold\color{white}#1}\fi\par\vspace{7pt}}}
% Exercice résolu : énoncé en haut, \tcblower puis la solution sur fond crème
\newcounter{popexo}\newcounter{popexr}
\newtcolorbox{exoresolu}[1][]{popbase,colframe=poporange,colbacklower=popcream,
  segmentation style={popink,line width=1.2pt,dashed},
  before upper={\stepcounter{popexr}\popboxhead{Worked example \thepopexr}{poporange}{#1}},
  before lower={\poppill{Solution}{popink}{popcream}\par\vspace{6pt}}}
% Exercice (non résolu en place) — la correction est dans la partie Corrigés
\newtcolorbox{exercice}[1][]{popbase,colframe=popink,
  before upper={\stepcounter{popexo}\popboxhead{Exercise \thepopexo}{popink}{#1}}}
% Corrigé
\newtcolorbox{corrige}[1][]{popbase,colframe=popgreen,colback=popgreenL,
  before upper={\popboxhead{Solution}{popgreen}{#1}}}
% Objectifs / prérequis / bilan
\newtcolorbox{objectifs}{popbase,colframe=popink,colback=poporangeL,
  before upper={\popboxhead{Chapter objectives}{popink}{}}}
\newtcolorbox{prerequis}{popbase,colframe=popink,colback=popblueL,
  before upper={\popboxhead{Prerequisites}{popblue}{}}}
\newtcolorbox{bilan}{popbase,colframe=popink,colback=white,boxrule=2.8pt,
  before upper={\popboxhead{Fiche bilan}{poporange}{L'essentiel à connaître pour l'examen}}}
% Figure encadrée avec légende
\newtcolorbox{popfigure}[1][]{enhanced,breakable=false,colback=white,colframe=popline,boxrule=1.4pt,arc=8pt,
  left=10pt,right=10pt,top=10pt,bottom=8pt,before skip=10pt,after skip=12pt,halign=center,
  lower separated=false,halign lower=center,
  fontlower=\popsemi\fontsize{9}{11.5}\selectfont\color{popmuted},
  #1}
\newcommand{\legende}[1]{\par\vspace{6pt}{\popsemi\fontsize{9}{11.5}\selectfont\color{popmuted}#1}}

% ============================================================
% Signature OnBuch+
% ============================================================
\newcommand{\poplogoinline}[1]{{\poplogo\fontsize{#1}{#1}\selectfont\color{popink}OnBuch\color{poporange}+}}
\newcommand{\signatureonbuch}{%
  \begin{tcolorbox}[enhanced,colback=popink,colframe=popink,boxrule=2.2pt,arc=10pt,
      drop shadow=poporange,left=14pt,right=14pt,top=10pt,bottom=10pt,before skip=0pt,after skip=0pt]
    \tikz[baseline=-0.7ex]\node[circle,fill=popgreen,draw=popcream,line width=1.3pt,minimum size=22pt,inner sep=0]{\popcheck{white}};%
    \hspace{9pt}\begin{minipage}[c]{0.62\linewidth}
      {\popxbold\fontsize{12}{13}\selectfont\color{popcream}Signed\ \textbullet\ {\poplogo OnBuch\color{poporange}+}}\par\vspace{2pt}
      {\raggedright\popsemi\fontsize{8}{10.5}\selectfont\color{popline}\MakeUppercase{OnBuch+ teaching team}\\\MakeUppercase{Follows the official MINESEC syllabus}\par}
    \end{minipage}\hfill
    {\poplogo\fontsize{22}{22}\selectfont\color{popcream}OnBuch\color{poporange}+}
  \end{tcolorbox}%
}

% ============================================================
% Page de garde
% #1 titre · #2 matière · #3 niveau · #4 module · #5 n° leçon · #6 durée ·
% #7 description / objectifs (texte ou itemize)
% ============================================================
\newcommand{\pagegarde}[7]{%
  \thispagestyle{empty}
  \newgeometry{a4paper,margin=1.7cm,top=1.6cm,bottom=1.4cm}%
  \begin{tikzpicture}[remember picture,overlay]
    \fill[poporange!18] ([xshift=-3.2cm,yshift=-3.6cm]current page.north east) circle (4.6cm);
    \fill[poppurple!14] ([xshift=2.4cm,yshift=7.6cm]current page.south west) circle (3.8cm);
  \end{tikzpicture}%
  {\poplogo\fontsize{30}{30}\selectfont\color{popink}OnBuch\color{poporange}+}\hfill
  \raisebox{0.6ex}{\poppill{Signed course \textbullet\ Premium}{popink}{popcream}}\par
  \vspace{1pt}\popsquiggle{poppurple}\par
  \vspace{8pt}
  \begin{tcolorbox}[enhanced,colback=poporange,colframe=popink,boxrule=2.8pt,arc=13pt,
      drop shadow=popink,left=18pt,right=18pt,top=12pt,bottom=13pt,after skip=10pt]
    \poppill{#2 \textbullet\ #3}{popink}{popcream}\hspace{5pt}\poppill{#4}{white}{popink}\par\vspace{8pt}
    {\popdisplay\fontsize{14}{14}\selectfont\color{popink}CHAPTER #5}\par\vspace{4pt}
    {\raggedright\hyphenpenalty=10000\exhyphenpenalty=10000\popdisplay\fontsize{25}{27}\selectfont\color{popcream}\MakeUppercase{#1}\par}
    \vspace{8pt}
    {\popxbold\fontsize{9.5}{9.5}\selectfont\color{popink}\MakeUppercase{Suggested time: #6}}
  \end{tcolorbox}
  \begin{tcolorbox}[enhanced,colback=white,colframe=popink,boxrule=2.2pt,arc=10pt,
      drop shadow=popink,left=16pt,right=16pt,top=10pt,bottom=10pt,after skip=8pt]
    \poppill{In this chapter}{poppurple}{white}\par\vspace{5pt}
    \popsemi\fontsize{9.3}{12.6}\selectfont\color{popink}#7
  \end{tcolorbox}
  \noindent\begin{minipage}[t]{0.487\linewidth}
    \begin{tcolorbox}[enhanced,colback=popgreen,colframe=popink,boxrule=2.2pt,arc=10pt,
        drop shadow=popink,left=11pt,right=11pt,top=10pt,bottom=10pt,height=3.1cm,valign=center]
      \tcbox[colback=white,colframe=popink,boxrule=1.3pt,arc=5pt,boxsep=0pt,nobeforeafter,
        left=3pt,right=3pt,top=3pt,bottom=3pt,valign=center]{\qrcode[height=1.9cm]{https://play.google.com/store/apps/details?id=com.luvvix.onbuch&hl=en}}%
      \hspace{8pt}\begin{minipage}[c]{0.5\linewidth}
        {\popdisplay\fontsize{11}{12.5}\selectfont\color{white}\MakeUppercase{Download OnBuch}}\par\vspace{4pt}
        {\raggedright\popsemi\fontsize{8.4}{11}\selectfont\color{white}Free \textbullet\ Google Play\\AI tutor, past papers, results\par}
      \end{minipage}
    \end{tcolorbox}
  \end{minipage}\hfill
  \begin{minipage}[t]{0.487\linewidth}
    \begin{tcolorbox}[enhanced,colback=poppurple,colframe=popink,boxrule=2.2pt,arc=10pt,
        drop shadow=popink,left=11pt,right=11pt,top=10pt,bottom=10pt,height=3.1cm,valign=center]
      \tikz[baseline=-0.7ex]\node[circle,fill=white,draw=popink,line width=1.3pt,minimum size=1.6cm,inner sep=0]{\popdisplay\fontsize{24}{24}\selectfont\color{poppurple}+};%
      \hspace{8pt}\begin{minipage}[c]{0.6\linewidth}
        {\popdisplay\fontsize{11}{12.5}\selectfont\color{white}\MakeUppercase{Go OnBuch+}}\par\vspace{4pt}
        {\raggedright\popsemi\fontsize{8.4}{11}\selectfont\color{white}All signed courses, unlimited AI tutor, PDF sheets — OnBuch+ tab in the app\par}
      \end{minipage}
    \end{tcolorbox}
  \end{minipage}
  \vspace{8pt}
  \popcontenu
  \vfill
  \signatureonbuch
  \restoregeometry
  \newpage
}

% Rangée « Ce cours contient » (page de garde)
\newcommand{\poptile}[3]{% #1 couleur #2 gros texte #3 légende
  \begin{tcolorbox}[enhanced,colback=#1,colframe=popink,boxrule=2pt,arc=9pt,shadow={2.5pt}{-2.5pt}{0pt}{popink},
      left=6pt,right=6pt,top=9pt,bottom=9pt,halign=center,valign=center,height=1.75cm,width=\linewidth,nobeforeafter]
    {\popdisplay\fontsize{12.5}{13}\selectfont\color{white}\MakeUppercase{#2}}\par\vspace{3pt}
    {\centering\popsemi\fontsize{7.8}{9.5}\selectfont\color{white}#3\par}
  \end{tcolorbox}}
\newcommand{\popcontenu}{%
  \par\noindent{\popxboldls\fontsize{7.5}{7.5}\selectfont\color{popmuted}THIS SIGNED COURSE CONTAINS}\par\vspace{7pt}
  \noindent\begin{minipage}[t]{0.235\linewidth}\poptile{popblue}{Course}{complete, illustrated, step by step}\end{minipage}\hfill
  \begin{minipage}[t]{0.235\linewidth}\poptile{poporange}{Methods}{and worked examples}\end{minipage}\hfill
  \begin{minipage}[t]{0.235\linewidth}\poptile{poppink}{Exercises}{exam style + solutions}\end{minipage}\hfill
  \begin{minipage}[t]{0.235\linewidth}\poptile{popgreen}{Summary sheet}{the essentials to remember}\end{minipage}\par}

% Sommaire
\newtcolorbox{sommaire}{enhanced,colback=white,colframe=popink,boxrule=2.2pt,arc=10pt,
  drop shadow=popink,left=18pt,right=18pt,top=13pt,bottom=13pt,before skip=6pt,after skip=20pt,
  before upper={\poppill{Contents}{poppurple}{white}\par\vspace{9pt}\popsemi\fontsize{10.6}{14.5}\selectfont\color{popink}}}

% Signature de fin de cours — poussée tout en bas de la dernière page
\newcommand{\findecours}{%
  \par\vfill\nopagebreak
  \begin{center}\popsquiggle{poporange}\end{center}\vspace{4pt}
  \signatureonbuch
}
\AtBeginDocument{\def\llbracket{\mathopen{[\mkern-3.5mu[}}\def\rrbracket{\mathclose{]\mkern-3.5mu]}}} % intervalles d'entiers (glyphe ⟦ absent de la police)
% Glyphes absents de Fira Math : redéfinis à partir de glyphes disponibles
\AtBeginDocument{%
  \def\setminus{\mathbin{\backslash}}\let\smallsetminus\setminus
  \def\vdots{\mathord{\vbox{\baselineskip3.2pt\lineskiplimit0pt\kern4pt\hbox{.}\hbox{.}\hbox{.}}}}%
  \def\ddots{\mathinner{\mkern1mu\raise6.5pt\hbox{.}\mkern2mu\raise3.6pt\hbox{.}\mkern2mu\raise0.7pt\hbox{.}\mkern1mu}}%
  \def\triangle{\mathord{\scalebox{0.9}{$\symup{\Delta}$}}}%
  \def\nabla{\mathord{\raisebox{\depth}{\scalebox{1}[-1]{$\symup{\Delta}$}}}}%
  \def\checkmark{\popcheck{popdark}}%
  \def\star{\mathbin{\ast}}\def\bigstar{\mathord{\ast}}%
  \def\diamond{\mathbin{\scalebox{0.6}{$\Diamond$}}}%
  \def\bigcup{\mathop{\mathchoice{\raisebox{-0.3em}{\scalebox{1.7}{$\cup$}}}{\scalebox{1.25}{$\cup$}}{\cup}{\cup}}\displaylimits}%
  \def\bigcap{\mathop{\mathchoice{\raisebox{-0.3em}{\scalebox{1.7}{$\cap$}}}{\scalebox{1.25}{$\cap$}}{\cap}{\cap}}\displaylimits}%
  \def\lesseqqgtr{\mathrel{\lessgtr}}\def\bigoplus{\mathop{\oplus}\displaylimits}%
}
\providecommand{\ssi}{if and only if} % abréviation fréquente
\AtBeginDocument{\providecommand{\wideparen}{\overparen}} % arc de cercle

% Fonctions usuelles que les modèles utilisent sans que amsmath les fournisse
\providecommand{\arsinh}{\operatorname{arsinh}}\providecommand{\arcosh}{\operatorname{arcosh}}
\providecommand{\artanh}{\operatorname{artanh}}\providecommand{\arsech}{\operatorname{arsech}}
\providecommand{\arcsch}{\operatorname{arcsch}}\providecommand{\arcoth}{\operatorname{arcoth}}
\providecommand{\arcsinh}{\operatorname{arsinh}}\providecommand{\arccosh}{\operatorname{arcosh}}
\providecommand{\arctanh}{\operatorname{artanh}}\providecommand{\sech}{\operatorname{sech}}
\providecommand{\csch}{\operatorname{csch}}\providecommand{\arcsec}{\operatorname{arcsec}}
\providecommand{\arccsc}{\operatorname{arccsc}}\providecommand{\arccot}{\operatorname{arccot}}
\providecommand{\sgn}{\operatorname{sgn}}\providecommand{\lcm}{\operatorname{lcm}}
\providecommand{\Var}{\operatorname{Var}}\providecommand{\Cov}{\operatorname{Cov}}
\providecommand{\permil}{\ensuremath{{}^{0}\!/_{00}}}\providecommand{\textperthousand}{\ensuremath{{}^{0}\!/_{00}}}

%% FIGFIX-BEGIN v5 — correctifs globaux des figures (docs/onbuch-generation/tools/figfix)
\usepackage{adjustbox}\usepackage{etoolbox}\tcbuselibrary{raster}
% toute figure TikZ/pgfplots plus large que la ligne est réduite à la largeur disponible
\BeforeBeginEnvironment{tikzpicture}{\begin{adjustbox}{max width=\linewidth}}
\AfterEndEnvironment{tikzpicture}{\end{adjustbox}}
% légendes pgfplots lisibles par-dessus les courbes
\pgfplotsset{every axis/.append style={legend style={fill=white,fill opacity=0.92,text opacity=1,draw=black!15,inner sep=3pt}}}
% étiquettes d'arêtes (syntaxe edge["..."]) avec fond blanc
\tikzset{every edge quotes/.append style={fill=white,inner sep=1.2pt}}
%% FIGFIX-END
\begin{document}

\pagegarde{International Economics}{Economics}{Upper Sixth}{Upper Sixth — Economics}{6}{21 periods}{This chapter introduces the theory and practice of international trade: why countries specialise and exchange, how trade is measured and restricted, how the balance of payments records a nation's transactions with the rest of the world, and how foreign exchange markets determine the value of currencies. It equips learners with the analytical tools needed for the GCE Advanced Level paper and for understanding Cameroon's place in the global economy.
\vspace{4pt}
\begin{multicols}{2}\small
\begin{itemize}
\item By the end of this chapter I can define international trade and distinguish it from internal trade
\item By the end of this chapter I can explain and apply the theories of absolute and comparative advantage to numerical examples
\item By the end of this chapter I can calculate the gains from specialisation and the exchange ratio between two countries
\item By the end of this chapter I can describe the main instruments of trade restriction and evaluate their effects
\item By the end of this chapter I can compute and interpret the terms of trade index
\item By the end of this chapter I can present the structure of the balance of payments and identify deficits and surpluses
\end{itemize}\end{multicols}}

\begin{sommaire}
\begin{multicols}{2}
\begin{enumerate}
\item Meaning and Importance of International Trade
\item Theories of International Trade
\item Specialisation and the Gains from Trade
\item International Trade Restrictions
\item Terms of Trade
\item The Balance of Payments: Structure and Recording
\item Balance of Payments Disequilibrium: Deficit, Surplus and Correction
\item International Economic Institutions and Regional Groupings
\item Foreign Exchange Market, Exchange Rates and Adjustments
\item Integration activity
\item Methods and worked examples
\item Exercises
\item Solutions to the exercises
\item Summary sheet
\end{enumerate}
\end{multicols}
\end{sommaire}

\begin{objectifs}
\begin{itemize}
\item By the end of this chapter I can define international trade and distinguish it from internal trade
\item By the end of this chapter I can explain and apply the theories of absolute and comparative advantage to numerical examples
\item By the end of this chapter I can calculate the gains from specialisation and the exchange ratio between two countries
\item By the end of this chapter I can describe the main instruments of trade restriction and evaluate their effects
\item By the end of this chapter I can compute and interpret the terms of trade index
\item By the end of this chapter I can present the structure of the balance of payments and identify deficits and surpluses
\item By the end of this chapter I can propose measures to correct a balance of payments disequilibrium, with reference to Cameroon
\item By the end of this chapter I can describe the roles of major international economic institutions and regional groupings
\item By the end of this chapter I can explain how the foreign exchange market works and how exchange rates are determined and adjusted
\end{itemize}
\end{objectifs}

\begin{prerequis}
\begin{itemize}
\item Demand, supply and market equilibrium (Lower Sixth)
\item Opportunity cost and the production possibility frontier
\item Price elasticity of demand and supply
\item National income and the circular flow of income
\item Basic index number calculation (price indices)
\end{itemize}
\end{prerequis}

\newpage

\coursec{Meaning and Importance of International Trade}

Every morning in Douala, a trader at Marché Mboppi sells imported rice from Thailand, a biker fuels his machine with petrol refined from imported crude, and a student in Buea pays school fees with a phone assembled in China --- yet Cameroon exports cocoa, oil and timber to pay for all this. Why does Cameroon not produce everything itself? Why does the government worry when imports exceed exports? And why does the value of the FCFA against the dollar matter to the price of bread in Bamenda? This chapter answers these questions by exploring how nations trade, how they pay each other, and how exchange rates shape our daily lives.

\courssub{The Meaning of International Trade}

\textbf{Intuition first.}\par
No household produces everything it consumes. A family in Ngaoundéré buys salt, soap and school books from others, and sells its own maize or labour to earn the money to pay for them. Exactly the same logic applies between nations. Just as specialisation makes a household richer, exchange between countries allows each nation to obtain goods it cannot produce --- or cannot produce cheaply --- in return for the goods it produces best.

\begin{definition}[International Trade]
\cle{International trade} (also called \emph{foreign} or \emph{external trade}) is the exchange of goods and services across national borders, that is, between residents of different countries. It consists of \cle{exports} (goods and services sold to foreigners) and \cle{imports} (goods and services bought from foreigners).
\end{definition}

It is useful to contrast this with what happens inside a single country.

\begin{definition}[Internal (Domestic) Trade]
\cle{Internal trade} is the buying and selling of goods and services \emph{within} the borders of one country, between its own residents, using the national currency --- for example, a Douala wholesaler selling garri to a retailer in Yaoundé.
\end{definition}

\begin{center}
\begin{tabularx}{\linewidth}{|l|X|X|}
\hline
\rowcolor{popblueL} \textbf{Feature} & \textbf{Internal trade} & \textbf{International trade} \\
\hline
Geographical scope & Within one country & Across national borders \\
\hline
Currency used & One national currency & Foreign currencies (dollars, euros) needed \\
\hline
Factors of production & Labour and capital move freely & Factors move with difficulty between nations \\
\hline
Government policy & Same laws and taxes everywhere & Tariffs, quotas, exchange controls may apply \\
\hline
Transport \& risk & Shorter distances, lower risk & Long distances, higher costs and risks \\
\hline
\end{tabularx}
\end{center}

\begin{attention}[Why the distinction matters]
Because international trade crosses borders, it faces obstacles that internal trade does not: different currencies (requiring foreign exchange), different languages and legal systems, trade barriers (tariffs and quotas), and greater transport and insurance costs. These obstacles explain why governments need special policies --- and special institutions --- to manage foreign trade, topics we develop later in this chapter.
\end{attention}

\courssub{Visible and Invisible Trade}

International trade is not only about physical goods you can touch at the port of Douala. A Cameroonian student paying tuition to a British university, or a tourist spending money at Kribi beach, is also ``trading'' internationally.

\begin{definition}[Visible and Invisible Trade]
\cle{Visible trade} is the export and import of \emph{physical, tangible goods} (merchandise): cocoa, crude oil, timber, cars, rice. \cle{Invisible trade} is the export and import of \emph{services and transfers}: tourism, banking, insurance, shipping, education, dividends, and remittances from Cameroonians living abroad.
\end{definition}

\begin{exemplebox}[Visible and invisible items for Cameroon]
\begin{itemize}
\item \textbf{Visible exports:} crude oil, cocoa beans, timber, coffee, cotton, bananas.
\item \textbf{Visible imports:} machinery, vehicles, refined fuel, medicines, rice, manufactured goods.
\item \textbf{Invisible exports:} spending by foreign tourists visiting Mount Cameroon; freight and port services provided to foreign ships.
\item \textbf{Invisible imports:} insurance premiums paid to foreign companies; school fees paid abroad; interest on foreign debt.
\end{itemize}
\end{exemplebox}

\begin{attention}[Do not confuse international trade with international finance]
A very common examination mistake is to treat \emph{international trade} and \emph{international finance} as the same thing. \textbf{International trade} concerns the flow of \emph{goods and services} (exports and imports). \textbf{International finance} concerns the flow of \emph{money and capital} (loans, investment, foreign exchange). The two are linked --- trade must be paid for --- but they are not identical. When a question asks about trade, discuss goods and services, not loans and capital flows.
\end{attention}

\courssub{Why Do Countries Trade?}

Why does Cameroon not simply produce everything at home? The fundamental reason is that countries are \emph{different}, and these differences make exchange mutually beneficial.

\begin{enumerate}
\item \textbf{Differences in natural resource endowments.} Some countries possess oil, minerals or fertile land; others do not. Cameroon has crude oil and tropical hardwoods; Japan has almost none and must import them.
\item \textbf{Differences in climate.} Cocoa and bananas grow in Cameroon's humid south; wheat and grapes grow better in temperate climates. No amount of effort would let Cameroon grow wheat as cheaply as Canada.
\item \textbf{Differences in technology and capital.} Industrialised countries produce aircraft, phones and medicines efficiently because of advanced technology; developing countries import these goods.
\item \textbf{Differences in labour skills.} Some nations have highly trained engineers; others have abundant cheap labour suited to agriculture or textiles.
\item \textbf{Differences in tastes and demand.} Even when a country \emph{could} produce a good, consumers may prefer the foreign variety --- Cameroonians buy Thai rice although rice is grown in the Ndop plain.
\item \textbf{Differences in costs of production.} Because of all the above, the \emph{opportunity cost} of producing each good differs from country to country. This is the deep economic basis of trade, developed formally in the theory of comparative advantage in the next section.
\end{enumerate}

\courssub{The Open Economy: Circular Flow}

In a \emph{closed economy}, income circulates only between households and firms. In an \emph{open economy} like Cameroon's, two new flows appear: \textbf{imports} are a \cle{leakage} (money leaves the national circuit to pay foreigners), while \textbf{exports} are an \cle{injection} (foreign money enters the circuit to pay our producers).

\begin{popfigure}
\begin{tikzpicture}[font=\small, >=stealth, node distance=1.6cm]
\node[draw, rounded corners, fill=popblueL, text width=2.6cm, align=center] (hh) at (0,0) {\textbf{Households}\\ (factor services)};
\node[draw, rounded corners, fill=popgreenL, text width=2.6cm, align=center] (fm) at (7,0) {\textbf{Firms}\\ (goods \& services)};
\node[draw, rounded corners, fill=popblueL, text width=2.8cm, align=center] (fw) at (3.5,-3.2) {\textbf{Foreign sector}\\ (rest of the world)};
\draw[->, thick, popblue] (hh.north east) to[bend left=15] node[above, text width=3cm, align=center]{Factors of production} (fm.north west);
\draw[->, thick, popblue] (fm.south west) to[bend left=15] node[below, text width=3cm, align=center]{Incomes (wages, profit)} (hh.south east);
\draw[->, thick, poporange] (hh.south) to[bend right=10] node[left, text width=2.2cm, align=center]{Imports $M$ (leakage)} (fw.west);
\draw[->, thick, popgreen] (fw.east) to[bend right=10] node[right, text width=2.2cm, align=center]{Exports $X$ (injection)} (fm.south);
\end{tikzpicture}
\legende{The circular flow in an open economy: imports leak income abroad; exports inject foreign income into the domestic circuit.}
\end{popfigure}

This diagram explains why the government ``worries when imports exceed exports'': if the leakage ($M$) persistently exceeds the injection ($X$), money drains out of the economy faster than it flows back in, creating a trade deficit and pressure on foreign currency reserves. The difference between export earnings and import spending on goods is called the \cle{balance of trade}; when services and transfers are added, we obtain the \cle{balance of payments}, studied later in this chapter.

\courssub{Advantages of International Trade}

\begin{enumerate}
\item \textbf{Wider markets.} Trade gives producers access to millions of consumers abroad. Cameroon's cocoa farmers sell to Europe and America, far beyond the small domestic market.
\item \textbf{Economies of scale.} Larger markets allow mass production at lower average cost, making firms more efficient and competitive.
\item \textbf{Access to goods not produced locally.} Cameroon obtains aircraft, medicines, machinery and phones that it cannot (yet) manufacture.
\item \textbf{Competition and lower prices.} Foreign competition forces local firms to improve quality and keep prices down, benefiting consumers.
\item \textbf{Foreign exchange earnings.} Exports earn dollars and euros, which the country needs to pay for imports and service its debts.
\item \textbf{Employment and growth.} Export industries (cocoa, timber, oil) create jobs and raise national income through the export multiplier.
\item \textbf{International cooperation.} Trade builds political and economic links between nations, encouraging peace and integration.
\end{enumerate}

\courssub{Disadvantages and Dangers of International Trade}

\begin{enumerate}
\item \textbf{Over-dependence.} A country relying on one or two exports (e.g.\ oil and cocoa) suffers badly when world prices fall.
\item \textbf{Dumping.} Foreign firms may sell goods below cost to capture the market, destroying local producers, then raise prices later.
\item \textbf{Exhaustion of natural resources.} Heavy export of timber and minerals can deplete Cameroon's forests and reserves, harming future generations.
\item \textbf{Imported inflation.} If the FCFA weakens or world prices rise, the cost of imported fuel and food rises, pushing up domestic prices --- the ``price of bread in Bamenda'' problem.
\item \textbf{Threat to infant industries.} Young local industries cannot compete with established foreign giants and may die before maturing, unless protected.
\item \textbf{Balance of payments problems.} Persistent excess of imports over exports creates deficits and external debt.
\end{enumerate}

\begin{attention}[Exam tip --- always localise your answer]
When an essay asks for the advantages (or disadvantages) of international trade, \textbf{always illustrate each point with a Cameroonian example}: cocoa exports for foreign exchange, imported crude for resource gaps, imported rice for consumer choice, timber exports for resource exhaustion. Examiners reward applied, localised answers far more than textbook lists.
\end{attention}

\courssub{Cameroon's International Trade}

Cameroon is a typical open developing economy. Its \textbf{main exports} are primary products: crude oil, cocoa, timber, coffee, cotton and bananas. Its \textbf{main imports} are manufactured and capital goods: machinery, vehicles, refined petroleum products, pharmaceuticals, food (especially rice and wheat) and other manufactured consumer goods. This structure --- exporting raw materials, importing manufactures --- is characteristic of many African economies and explains both their vulnerability to world price swings and the push for local transformation (e.g.\ grinding cocoa locally instead of exporting raw beans).

\begin{popfigure}
\begin{tikzpicture}
\begin{axis}[
    ybar, bar width=0.55cm,
    width=0.85\linewidth, height=6.2cm,
    ymin=0, ymax=35,
    ylabel={Share of export earnings (\%)},
    symbolic x coords={Crude oil, Cocoa, Timber, Cotton, Coffee, Bananas},
    xtick=data,
    x tick label style={rotate=25, anchor=east, font=\small},
    nodes near coords, nodes near coords style={font=\small},
    ymajorgrids, grid style={popline},
    every axis plot/.append style={fill=popblue, draw=popdark}
]
\addplot coordinates {(Crude oil,30) (Cocoa,15) (Timber,12) (Cotton,8) (Coffee,6) (Bananas,5)};
\end{axis}
\end{tikzpicture}
\legende{Illustrative structure of Cameroon's main export products by approximate share of export earnings (orders of magnitude for teaching purposes).}
\end{popfigure}

The bar chart shows the danger of concentration: crude oil alone dominates export earnings, so a fall in world oil prices immediately reduces government revenue and foreign exchange --- a powerful illustration of the risk of over-dependence discussed above.

\begin{exoresolu}[Classifying trade flows]
A GCE question gives the following transactions of Cameroon and asks you to classify each as (i) visible export, (ii) visible import, (iii) invisible export, (iv) invisible import:
\begin{enumerate}
\item Sale of cocoa beans to a chocolate maker in the Netherlands.
\item Purchase of rice from Thailand by a Douala importer.
\item School fees paid by a Cameroonian student at a university in France.
\item Spending by a German tourist visiting the Waza National Park.
\item Purchase of a generator from China by an electricity company in Yaoundé.
\end{enumerate}
\tcblower
\textbf{Solution.}
\begin{enumerate}
\item \textbf{Visible export} --- a physical good (cocoa) sold abroad; it earns foreign exchange for Cameroon.
\item \textbf{Visible import} --- a physical good (rice) bought from abroad; it uses up foreign exchange.
\item \textbf{Invisible import} --- Cameroon is \emph{buying} an education service from France; money flows out for a service, not a good.
\item \textbf{Invisible export} --- Cameroon is \emph{selling} a tourism service to the German visitor; foreign money flows in.
\item \textbf{Visible import} --- a generator is a tangible capital good purchased from abroad.
\end{enumerate}
\textbf{Key test:} ask two questions --- (a) is it a \emph{good} (visible) or a \emph{service/transfer} (invisible)? (b) does money flow \emph{in} (export) or \emph{out} (import)?
\end{exoresolu}

\begin{aretenir}[Key points]
\begin{itemize}
\item International trade is the exchange of goods and services across national borders; internal trade occurs within one country.
\item Visible trade = goods; invisible trade = services and transfers (tourism, insurance, fees, remittances).
\item Countries trade because of differences in resources, climate, technology, skills, tastes and costs.
\item Imports are a leakage, exports an injection, in the circular flow of an open economy.
\item Trade brings wider markets, scale economies, cheaper goods and foreign exchange --- but also dependence, dumping, imported inflation and threats to infant industries.
\item Cameroon exports mainly primary products (oil, cocoa, timber) and imports manufactures --- a source of vulnerability.
\end{itemize}
\end{aretenir}

\begin{exercice}[Practice]
\begin{enumerate}
\item Define international trade and distinguish it from internal trade, giving one Cameroonian example of each.
\item Classify the following as visible or invisible, export or import: (a) export of timber logs to China; (b) insurance premium paid to a London firm; (c) remittances sent home by Cameroonians in Germany; (d) import of refined fuel.
\item Explain four reasons why Cameroon trades with the rest of the world.
\item ``International trade is an unmixed blessing for Cameroon.'' Discuss, using at least three advantages and three disadvantages with local examples.
\end{enumerate}
\end{exercice}

\begin{corrige}[Exercise 1]
\textbf{1.} International trade is the exchange of goods and services across national borders (e.g.\ Cameroon selling cocoa to the Netherlands); internal trade is exchange within one country (e.g.\ a Douala trader selling garri to a Yaoundé retailer).
\textbf{2.} (a) Visible export; (b) invisible import; (c) invisible export (a transfer inflow, recorded as an invisible receipt); (d) visible import.
\textbf{3.} Any four: differences in resource endowments (oil, timber), climate (cocoa cannot grow in Europe), technology (imports of machinery and phones), skills, tastes (preference for Thai rice), and cost differences making exchange cheaper than self-sufficiency.
\textbf{4.} Structure: introduction (define trade); advantages with examples (wider markets for cocoa, foreign exchange from oil, cheaper imported goods, scale economies); disadvantages with examples (dependence on oil prices, dumping of cheap imports killing local firms, imported inflation raising bread prices, exhaustion of forests); conclusion --- trade is beneficial \emph{if} managed with diversification and protection of infant industries.
\end{corrige}

\coursec{Theories of International Trade}

In the previous section we established that international trade is the exchange of goods and services across national borders, and we saw why countries trade: differences in climate, natural resources, technology and tastes. But a deeper question remains: \emph{which} goods should a country produce and export, and \emph{why} does trade benefit \emph{both} partners even when one of them seems better at producing everything? Economists have answered these questions over more than three centuries. This section presents the great theories of international trade, from the mercantilists to the Heckscher--Ohlin model, with the numerical tools you must master for the GCE Advanced Level examination.

\courssub{The Mercantilist View of Trade}

The earliest systematic thinking about trade came from the \cle{mercantilists}, a group of writers, merchants and state advisers active in Europe roughly between the sixteenth and eighteenth centuries (think of the era of the great colonial trading companies).

Their central ideas were:

\begin{itemize}
\item A nation's wealth is measured by its stock of \textbf{precious metals} (gold and silver).
\item A country should therefore \textbf{export more than it imports}, so that gold and silver flow in as payment.
\item Exports are ``good'' and imports are ``bad''; the state should actively promote exports and restrict imports through tariffs and quotas.
\item Trade is a \textbf{zero-sum game}: what one nation gains, another must lose.
\end{itemize}

\begin{attention}[The mercantilist fallacy]
The mercantilists confused \emph{money} with \emph{wealth}. Real wealth is the goods and services a nation can consume, not the coins in its treasury. Moreover, trade is \textbf{not} a zero-sum game: as Smith and Ricardo later proved, \emph{both} trading partners can gain simultaneously. Nevertheless, mercantilist instincts survive today whenever politicians speak of exports as victories and imports as defeats.
\end{attention}

\courssub{Adam Smith and Absolute Advantage (1776)}

In his famous book \emph{An Inquiry into the Nature and Causes of the Wealth of Nations} (1776), the Scottish economist \textbf{Adam Smith} attacked the mercantilists. His intuition was simple and powerful: just as it would be foolish for a tailor to make his own shoes when the shoemaker makes them better and cheaper, it is foolish for a country to produce at home what another country produces more efficiently.

\begin{definition}[Absolute advantage]
A country has an \cle{absolute advantage} in the production of a good if it can produce a unit of that good using \textbf{fewer resources} (for example, fewer labour hours) than another country --- or, equivalently, if it can produce \textbf{more output} from the same quantity of resources.
\end{definition}

\begin{propriete}[Smith's principle of trade]
Each country should \textbf{specialise} in the production of the good in which it has an absolute advantage, export part of its output, and import the good in which it has an absolute \emph{disadvantage}. Total world output then rises, and both countries can consume more than in self-sufficiency (autarky). (Statement examinable; no formal proof required.)
\end{propriete}

\textbf{Numerical illustration.} Suppose one hour of labour produces the following quantities in two countries:

\begin{center}
\begin{tabular}{|l|c|c|}
\hline
\rowcolor{popblueL} & \textbf{Cloth (metres per labour hour)} & \textbf{Wine (litres per labour hour)} \\
\hline
Cameroon & 10 & 2 \\
\hline
Portugal & 6 & 8 \\
\hline
\end{tabular}
\end{center}

Cameroon produces more cloth per hour ($10>6$): it has the \textbf{absolute advantage in cloth}. Portugal produces more wine per hour ($8>2$): it has the \textbf{absolute advantage in wine}. If each country shifts labour towards its speciality, world output of \emph{both} goods rises. For example, moving one hour of Cameroonian labour from wine to cloth loses 2 litres of wine but gains 10 metres of cloth, while moving one hour of Portuguese labour from cloth to wine loses 6 metres of cloth but gains 8 litres of wine: the world gains $10-6=4$ metres of cloth \emph{and} $8-2=6$ litres of wine from the same total labour. Exchange then allows both countries to consume beyond what they could produce alone. This is Smith's case for \textbf{free trade}.

\begin{attention}[The question Smith could not answer]
Smith's theory works only when each country has an absolute advantage in \emph{something}. But what if one country is more efficient at producing \textbf{both} goods? Should trade take place at all? Smith had no answer. It was David Ricardo who solved this puzzle in 1817.
\end{attention}

\courssub{David Ricardo and Comparative Advantage (1817)}

In \emph{On the Principles of Political Economy and Taxation} (1817), \textbf{David Ricardo} demonstrated the most famous result in all of international economics: even if a country has an absolute advantage in \emph{everything}, both countries still gain from trade, provided each specialises according to its \textbf{comparative advantage}.

\begin{definition}[Comparative advantage]
A country has a \cle{comparative advantage} in the production of a good if it can produce that good at a \textbf{lower opportunity cost} than the other country --- that is, it sacrifices \emph{less} of the other good to produce one more unit of it.
\end{definition}

\begin{propriete}[Ricardo's law of comparative advantage]
If each country specialises in the good for which its opportunity cost is lowest, total world output increases, and mutually beneficial trade is possible at any exchange ratio lying \textbf{between} the two countries' opportunity cost ratios. (Statement and numerical application examinable.)
\end{propriete}

\begin{aretenir}[Assumptions of Ricardo's model]
Ricardo's demonstration rests on strong simplifying assumptions:
\begin{itemize}
\item Only \textbf{two countries} and \textbf{two commodities} (a $2\times 2$ model).
\item \textbf{Labour is the only factor} of production (labour theory of value).
\item \textbf{Constant costs}: productivity does not change as output expands.
\item \textbf{Full employment} of resources in both countries.
\item Labour is \textbf{mobile between industries} within a country but \textbf{immobile between countries}.
\item \textbf{No transport costs} and no trade barriers (free trade).
\item Technology is fixed and differs between countries.
\end{itemize}
These assumptions are exactly the points attacked by critics --- see the limitations below.
\end{aretenir}

\textbf{Ricardo's famous example: England and Portugal.} Ricardo used the labour hours needed to produce one unit of cloth and one unit of wine:

\begin{popfigure}
\begin{center}
\begin{tabular}{|l|c|c|}
\hline
\rowcolor{popblueL} \textbf{Country} & \textbf{Cloth (labour hours per unit)} & \textbf{Wine (labour hours per unit)} \\
\hline
Portugal & 90 & 80 \\
\hline
England & 100 & 120 \\
\hline
\end{tabular}
\end{center}
\legende{Labour hours required to produce one unit of each good (Ricardo's example). Portugal is more efficient in both goods, yet trade still benefits both countries.}
\end{popfigure}

Observe carefully: Portugal needs \emph{fewer} hours for both goods, so Portugal has an \textbf{absolute advantage in both}. Does this kill trade? No. Compute the opportunity costs:

\begin{itemize}
\item \textbf{Portugal:} 1 unit of cloth costs $\dfrac{90}{80}=1.125$ units of wine; 1 unit of wine costs $\dfrac{80}{90}\approx 0.89$ units of cloth.
\item \textbf{England:} 1 unit of cloth costs $\dfrac{100}{120}\approx 0.83$ units of wine; 1 unit of wine costs $\dfrac{120}{100}=1.2$ units of cloth.
\end{itemize}

England gives up only $0.83$ wine for a cloth, against $1.125$ in Portugal: \textbf{England has the comparative advantage in cloth}. Portugal gives up only $0.89$ cloth for a wine, against $1.2$ in England: \textbf{Portugal has the comparative advantage in wine}. Each country should specialise accordingly, and any exchange ratio between $0.83$ and $1.125$ wine per cloth benefits both.

\begin{attention}[The classic examination trap]
A country can have an \textbf{absolute advantage in both goods} and yet still gain from trade. Never conclude ``no trade is possible'' because one country is better at everything. What matters is \emph{comparative} (opportunity-cost) advantage, never absolute advantage alone. This is the single most common error in GCE answers on this topic.
\end{attention}

\courssub{Computing Opportunity Costs: Method and Worked Examples}

\begin{methode}[How to identify comparative advantage from data]
\begin{enumerate}
\item \textbf{Read the table carefully.} Are the figures \emph{labour hours per unit of output} (input data) or \emph{output per labour hour} (output data)? The treatment differs.
\item \textbf{Input data (hours per unit):} the opportunity cost of good $X$ in terms of good $Y$ is $\dfrac{\text{hours for one } X}{\text{hours for one } Y}$.
\item \textbf{Output data (units per hour):} the opportunity cost of good $X$ in terms of good $Y$ is $\dfrac{\text{output of } Y \text{ per hour}}{\text{output of } X \text{ per hour}}$.
\item \textbf{Compare the ratios} between the two countries for each good.
\item \textbf{Assign specialisation:} each country specialises in the good for which its opportunity cost is \emph{lower}.
\item \textbf{Conclude on the gains:} a terms-of-trade ratio between the two opportunity costs makes both countries better off.
\end{enumerate}
\end{methode}

\begin{exoresolu}[Applying comparative advantage to Cameroon and Nigeria]
The table below shows the output per labour hour in two (hypothetical) economies:

\begin{center}
\begin{tabular}{|l|c|c|}
\hline
\rowcolor{popblueL} & \textbf{Cocoa (kg per hour)} & \textbf{Cassava (kg per hour)} \\
\hline
Cameroon & 12 & 6 \\
\hline
Nigeria & 4 & 8 \\
\hline
\end{tabular}
\end{center}

\textbf{Required:} (a) Identify absolute advantages. (b) Compute opportunity costs and assign specialisation. (c) Suggest a mutually beneficial exchange ratio.

\tcblower
\textbf{Solution.}

\textbf{(a) Absolute advantage.} Cameroon produces more cocoa per hour ($12>4$); Nigeria produces more cassava per hour ($8>6$). Cameroon has the absolute advantage in cocoa, Nigeria in cassava.

\textbf{(b) Opportunity costs} (output data, so we divide the \emph{other} good's output by the good's own output):

Cameroon: 1 kg of cocoa costs $\dfrac{6}{12}=0.5$ kg of cassava; 1 kg of cassava costs $\dfrac{12}{6}=2$ kg of cocoa.

Nigeria: 1 kg of cocoa costs $\dfrac{8}{4}=2$ kg of cassava; 1 kg of cassava costs $\dfrac{4}{8}=0.5$ kg of cocoa.

Cameroon sacrifices only $0.5$ kg of cassava per kg of cocoa (against $2$ in Nigeria): \textbf{Cameroon has the comparative advantage in cocoa}. Nigeria sacrifices only $0.5$ kg of cocoa per kg of cassava (against $2$ in Cameroon): \textbf{Nigeria has the comparative advantage in cassava}.

\textbf{(c) Exchange ratio.} Any ratio between $0.5$ and $2$ kg of cassava per kg of cocoa --- for example $1$ kg of cocoa for $1$ kg of cassava --- leaves \emph{both} countries better off than producing both goods themselves.
\end{exoresolu}

\begin{exoresolu}[Comparative advantage when one country is better at everything]
Labour hours per unit of output:

\begin{center}
\begin{tabular}{|l|c|c|}
\hline
\rowcolor{popblueL} & \textbf{Textiles} & \textbf{Timber} \\
\hline
Country A & 2 hours & 4 hours \\
\hline
Country B & 6 hours & 8 hours \\
\hline
\end{tabular}
\end{center}

\textbf{Required:} Show that Country B, although less efficient in both goods, still gains from trade.

\tcblower
\textbf{Solution.} Country A has the absolute advantage in both goods ($2<6$ and $4<8$). But compute opportunity costs (input data):

Country A: 1 textile costs $\dfrac{2}{4}=0.5$ timber; 1 timber costs $\dfrac{4}{2}=2$ textiles.

Country B: 1 textile costs $\dfrac{6}{8}=0.75$ timber; 1 timber costs $\dfrac{8}{6}\approx 1.33$ textiles.

Country A has the lower opportunity cost in \textbf{textiles} ($0.5<0.75$); Country B has the lower opportunity cost in \textbf{timber} ($1.33<2$). If A specialises in textiles and B in timber, and they exchange at, say, 1 textile for 0.6 timber (between 0.5 and 0.75), A receives more timber per textile than it could produce domestically, and B pays less timber per textile than its own cost of $0.75$. \textbf{Both gain}, confirming Ricardo's law.
\end{exoresolu}

\courssub{Illustrating the Gains: Production Possibility Frontiers}

The gains from specialisation can be shown graphically with \cle{production possibility frontiers} (PPFs). Under Ricardo's assumption of constant costs, each country's PPF is a straight line whose slope equals the opportunity cost ratio. Using Ricardo's own numbers, suppose England commands $12\,000$ labour hours and Portugal $9\,000$ labour hours. England can then produce at most $\frac{12000}{100}=120$ units of cloth or $\frac{12000}{120}=100$ units of wine; Portugal at most $\frac{9000}{90}=100$ units of cloth or $\frac{9000}{80}=112.5$ units of wine.

\begin{popfigure}
\begin{center}
\begin{tikzpicture}
\begin{axis}[
width=0.62\linewidth, height=7cm,
xlabel={Cloth (units)}, ylabel={Wine (units)},
xmin=0, xmax=135, ymin=0, ymax=135,
xtick={0,50,100,120}, ytick={0,50,100,112.5},
legend style={at={(0.98,0.98)}, anchor=north east, font=\small},
grid=both, grid style={popline!40},
axis lines=left, thick]
\addplot[popblue, very thick, domain=0:120] {100 - 0.833*x};
\addlegendentry{England PPF}
\addplot[poporange, very thick, domain=0:100] {112.5 - 1.125*x};
\addlegendentry{Portugal PPF}
\addplot[popgreen, dashed, very thick, domain=0:120] {120 - x};
\addlegendentry{England with trade (1 cloth : 1 wine)}
\end{axis}
\end{tikzpicture}
\end{center}
\legende{PPFs for England and Portugal under constant costs. England specialises fully in cloth (120 units) and trades at 1 cloth for 1 wine (a ratio between the two opportunity costs, $0.83$ and $1.125$): its consumption line (dashed) lies \emph{outside} its PPF, so England can consume more wine for any given amount of cloth than it could produce alone. Portugal gains symmetrically.}
\end{popfigure}

The key message of the figure: each country's PPF shows what it can \emph{produce}; with specialisation and trade at an intermediate exchange ratio, each country can \emph{consume} on a line lying \textbf{outside} its PPF. Trade is therefore equivalent to an outward shift of consumption possibilities --- a gain obtainable without any extra resources.

\courssub{Limitations of the Classical Theories}

Ricardo's model is logically flawless \emph{given its assumptions}, but the assumptions themselves are unrealistic. A good GCE essay must evaluate the theory:

\begin{itemize}
\item \textbf{Constant costs.} In reality, expanding production often meets \emph{increasing} opportunity costs (diminishing returns), so PPFs are concave, not straight. Complete specialisation is then rarely optimal.
\item \textbf{Full employment.} The model assumes all resources are employed. In economies with unemployment and underemployment (a real concern in Cameroon and the CEMAC region), the gains-from-trade argument must be qualified.
\item \textbf{Two countries, two goods.} The real world has many countries and thousands of products; extending the model is possible but the simple clarity is lost.
\item \textbf{Transport costs ignored.} Moving cocoa from Kumba to the port of Douala and on to Europe, or refined products from SONARA, is costly. High transport costs can wipe out a comparative advantage and make some goods non-tradable.
\item \textbf{Factor immobility.} Ricardo assumed labour moves freely between industries at home but not across borders. In fact, workers cannot instantly switch from farming to manufacturing (retraining, location), and labour \emph{does} migrate internationally.
\item \textbf{Labour as the only factor.} Capital, land and entrepreneurship also matter; ignoring them hides the true source of advantage.
\item \textbf{Static analysis.} Technology is taken as fixed, yet trade itself changes technology, skills and productivity over time.
\item \textbf{Distributional effects ignored.} Even when a nation gains overall, some groups (import-competing workers) may lose --- a point that explains resistance to free trade.
\end{itemize}

\courssub{The Heckscher--Ohlin Factor Endowment Theory (Outline)}

Two Swedish economists, \textbf{Eli Heckscher} and his student \textbf{Bertil Ohlin}, proposed (in the early twentieth century) a deeper explanation of \emph{why} comparative advantages differ.

\begin{definition}[Heckscher--Ohlin theorem]
A country has a comparative advantage in, and will \textbf{export}, the goods whose production uses intensively its \textbf{relatively abundant} (and therefore relatively cheap) factor of production; it will \textbf{import} goods intensive in its relatively \textbf{scarce} factor.
\end{definition}

The intuition: differences in factor endowments create differences in factor prices, which create differences in production costs, which create comparative advantage. A \textbf{labour-abundant} country (such as Cameroon, with a large and relatively low-cost workforce) should specialise in \textbf{labour-intensive} goods (agricultural products like cocoa, coffee, cotton); a \textbf{capital-abundant} country (such as Germany or Japan) should specialise in \textbf{capital-intensive} goods (machinery, vehicles, electronics).

\begin{savaistu}[The Leontief paradox]
When the economist Wassily Leontief tested the Heckscher--Ohlin theory on United States data (1953), he found that US \emph{exports} were more labour-intensive than US \emph{imports} --- the opposite of what the theory predicted for a capital-abundant country! This famous ``Leontief paradox'' showed that real-world trade also depends on skills (human capital), technology and natural resources, not just on capital and labour quantities.
\end{savaistu}

\courssub{Practical Work 23: Applying the Theories to Data}

\begin{experience}[Guided application: absolute and comparative advantage]
Work through the following data table (output per labour hour) for two hypothetical economies, ``Northland'' and ``Southland'':

\begin{center}
\begin{tabular}{|l|c|c|}
\hline
\rowcolor{popblueL} & \textbf{Steel (tonnes)} & \textbf{Rice (tonnes)} \\
\hline
Northland & 10 & 5 \\
\hline
Southland & 3 & 2 \\
\hline
\end{tabular}
\end{center}

\textbf{Step 1 --- Absolute advantage.} Northland produces more of both goods per hour: it has the absolute advantage in steel \emph{and} rice. A student reasoning like Adam Smith might wrongly stop here and predict no trade.

\textbf{Step 2 --- Opportunity costs.} Northland: 1 tonne of steel costs $\frac{5}{10}=0.5$ tonne of rice. Southland: 1 tonne of steel costs $\frac{2}{3}\approx 0.67$ tonne of rice. Northland's steel is cheaper in opportunity-cost terms.

\textbf{Step 3 --- Specialisation.} Northland $\rightarrow$ steel; Southland $\rightarrow$ rice (check: Southland's rice costs $\frac{3}{2}=1.5$ steel, Northland's rice costs $\frac{10}{5}=2$ steel; Southland is indeed the cheaper rice producer).

\textbf{Step 4 --- Interpretation.} Trade at any ratio between $0.5$ and $0.67$ rice per tonne of steel benefits both. Conclusion: Ricardo's law holds even when one country dominates in absolute terms.
\end{experience}

\begin{exercice}[Test yourself]
The table shows labour hours per unit of output:

\begin{center}
\begin{tabular}{|l|c|c|}
\hline
\rowcolor{popblueL} & \textbf{Shoes (pairs)} & \textbf{Baskets} \\
\hline
Country X & 5 & 10 \\
\hline
Country Y & 8 & 12 \\
\hline
\end{tabular}
\end{center}

\begin{enumerate}
\item Which country has the absolute advantage in each good?
\item Compute the opportunity cost of shoes in each country.
\item Assign specialisation and propose a mutually beneficial exchange ratio.
\item State two assumptions of Ricardo's model and explain why each may fail in the Cameroonian economy.
\end{enumerate}
\end{exercice}

\begin{corrige}[Exercise: shoes and baskets]
\textbf{1.} Country X needs fewer hours for both goods ($5<8$; $10<12$): absolute advantage in both.

\textbf{2.} Country X: 1 pair of shoes costs $\frac{5}{10}=0.5$ basket. Country Y: 1 pair of shoes costs $\frac{8}{12}\approx 0.67$ basket.

\textbf{3.} X has the comparative advantage in shoes ($0.5<0.67$); Y in baskets (Y: $\frac{12}{8}=1.5$ shoes per basket $<$ X: $\frac{10}{5}=2$). Any exchange between $0.5$ and $0.67$ basket per pair of shoes benefits both.

\textbf{4.} Any two of: constant costs (fails with diminishing returns in agriculture); full employment (fails given underemployment in Cameroon); no transport costs (fails given poor rural roads); labour immobility between countries (qualified by migration); labour as the only factor (ignores capital and land).
\end{corrige}

\begin{aretenir}[Key points of the section]
\begin{itemize}
\item Mercantilists saw trade as a zero-sum game aimed at accumulating gold --- a view refuted by classical economists.
\item \textbf{Smith:} specialise where you have an \emph{absolute} advantage; trade raises world output.
\item \textbf{Ricardo:} specialise where you have a \emph{comparative} (lower opportunity cost) advantage; even a country with absolute advantage in everything gains from trade.
\item Opportunity cost of good $X$ = amount of good $Y$ sacrificed per unit of $X$; beneficial exchange ratios lie between the two countries' opportunity costs.
\item The classical model rests on strong assumptions (constant costs, full employment, $2\times2$, no transport costs, factor immobility) which limit its realism.
\item \textbf{Heckscher--Ohlin:} export goods intensive in your abundant factor --- labour-abundant Cameroon should export labour-intensive products.
\end{itemize}
\end{aretenir}

Having established \emph{why} countries specialise and trade, the next section examines the practical question that follows directly: once countries have specialised, \emph{how is the exchange of goods actually determined}, and how are the gains shared between the trading partners?

\coursec{Specialisation and the Gains from Trade}

In the previous section we studied the \emph{theories} of international trade: Adam Smith's absolute advantage and Ricardo's comparative advantage both concluded that nations gain by specialising. But a theory is only convincing if we can \emph{measure} the gain. How much extra output does specialisation create? At what price will the two countries exchange their products? Who gets the larger share of the gain? This section answers these questions with simple but powerful numerical tools that the GCE examiner expects you to master perfectly.

\courssub{The Meaning of Specialisation}

\begin{definition}[Specialisation]
\cle{Specialisation} is the concentration of a country's (or a worker's, or a firm's) productive resources on the production of those goods and services in which it is most efficient, relying on trade to obtain the other goods it needs.
\end{definition}

We distinguish two forms:

\begin{definition}[Complete and Partial Specialisation]
\cle{Complete specialisation} occurs when a country devotes \emph{all} its resources to the production of a single commodity (the one in which it has a comparative advantage) and produces nothing of the other good.

\cle{Partial specialisation} occurs when a country shifts resources \emph{towards} the good of its comparative advantage but continues to produce some quantity of the other good domestically.
\end{definition}

\textbf{Why is specialisation often only partial in reality?}\par
\begin{itemize}
\item \textbf{Increasing opportunity costs.} The simple Ricardian model assumes constant costs. In reality, as more and more resources move into one industry, the resources transferred become less and less suitable, so the opportunity cost rises. Beyond a certain point, further specialisation is no longer profitable.
\item \textbf{Strategic and security reasons.} A country does not want to depend entirely on foreigners for food or defence goods. Cameroon, for example, maintains domestic food production even where imports might be cheaper.
\item \textbf{Risk diversification.} Depending on a single export (e.g.\ only cocoa or only oil) exposes the economy to world price shocks.
\item \textbf{Transport costs and tariffs} reduce the gains from full specialisation.
\end{itemize}

For the numerical analysis that follows, we retain the classical assumptions: two countries, two goods, one factor (labour), constant costs, no transport costs, free trade.

\courssub{World Output Before and After Specialisation}

The central claim of the theory of comparative advantage is that \emph{total world output rises} when each country specialises according to its comparative advantage, even if one country is more efficient at producing \emph{both} goods. Let us prove this numerically.

\begin{exemplebox}[Setting up the problem]
Suppose one unit of labour in each country can produce the following outputs per day:

\begin{center}
\begin{tabularx}{\linewidth}{|l|X|X|}
\hline
\rowcolor{popblueL} \textbf{Country} & \textbf{Cocoa (kg per labour unit)} & \textbf{Cloth (metres per labour unit)} \\
\hline
Cameroon & 10 & 5 \\
\hline
Nigeria & 6 & 12 \\
\hline
\end{tabularx}
\end{center}

Each country has 100 units of labour and, \emph{before specialisation}, devotes half of its labour (50 units) to each good.
\end{exemplebox}

\textbf{Step 1: identify the comparative advantage.}\par
Internal opportunity cost ratios:
\begin{itemize}
\item Cameroon: 10 kg of cocoa ``costs'' 5 m of cloth, so $1$ kg cocoa $= 0.5$ m cloth; equivalently $1$ m cloth $= 2$ kg cocoa.
\item Nigeria: 12 m of cloth ``costs'' 6 kg of cocoa, so $1$ m cloth $= 0.5$ kg cocoa; equivalently $1$ kg cocoa $= 2$ m cloth.
\end{itemize}
Cocoa is relatively cheap in Cameroon ($0.5$ m of cloth against $2$ m in Nigeria); cloth is relatively cheap in Nigeria. Cameroon should specialise in \cle{cocoa}, Nigeria in \cle{cloth} --- even though Nigeria is \emph{absolutely} more efficient in cloth and Cameroon in cocoa, the logic would be identical even if one country were better at both.

\textbf{Step 2: compute world output before specialisation.}\par
\begin{itemize}
\item Cameroon: $50 \times 10 = 500$ kg cocoa; $50 \times 5 = 250$ m cloth.
\item Nigeria: $50 \times 6 = 300$ kg cocoa; $50 \times 12 = 600$ m cloth.
\item \textbf{World total: 800 kg cocoa and 850 m cloth.}
\end{itemize}

\textbf{Step 3: compute world output after complete specialisation.}\par
\begin{itemize}
\item Cameroon (all 100 labour units on cocoa): $100 \times 10 = 1000$ kg cocoa; $0$ m cloth.
\item Nigeria (all 100 labour units on cloth): $0$ kg cocoa; $100 \times 12 = 1200$ m cloth.
\item \textbf{World total: 1000 kg cocoa and 1200 m cloth.}
\end{itemize}

\begin{center}
\begin{tabularx}{\linewidth}{|l|X|X|X|X|}
\hline
\rowcolor{popblueL} & \multicolumn{2}{c|}{\textbf{Before specialisation}} & \multicolumn{2}{c|}{\textbf{After specialisation}} \\
\hline
\rowcolor{popblueL} \textbf{Country} & \textbf{Cocoa (kg)} & \textbf{Cloth (m)} & \textbf{Cocoa (kg)} & \textbf{Cloth (m)} \\
\hline
Cameroon & 500 & 250 & 1000 & 0 \\
\hline
Nigeria & 300 & 600 & 0 & 1200 \\
\hline
\textbf{World total} & \textbf{800} & \textbf{850} & \textbf{1000} & \textbf{1200} \\
\hline
\end{tabularx}
\end{center}

\begin{propriete}[The Gain from Specialisation]
When each country specialises completely in the good of its comparative advantage, total world output of \emph{both} goods increases, with no additional resources used. In our example the world gains $+200$ kg of cocoa and $+350$ m of cloth from the same 200 units of labour. \emph{(The numerical demonstration, as above, is examinable and must be reproduced step by step.)}
\end{propriete}

This extra output is the \cle{gain from trade}: it is a ``free lunch'' created purely by reallocating labour according to comparative advantage.

\courssub{Determining the Terms of Exchange}

World output has risen, but the two countries must now \emph{exchange} cocoa against cloth. At what ratio? This ratio is called the \cle{terms of exchange} (or international exchange ratio).

\begin{definition}[Terms of Exchange]
The \cle{terms of exchange} is the rate at which one good is exchanged for another between two countries, e.g.\ 1 kg of cocoa against $x$ metres of cloth.
\end{definition}

\begin{propriete}[Limits of the Exchange Ratio]
For trade to be \emph{mutually beneficial}, the international exchange ratio must lie \textbf{strictly between} the two countries' internal (domestic) opportunity cost ratios:
\[
\text{Cameroon's ratio} \;<\; \text{terms of exchange} \;<\; \text{Nigeria's ratio}
\]
i.e.\ in our example, $1$ kg of cocoa must exchange for between $0.5$ m and $2$ m of cloth.
\end{propriete}

\textbf{Why must the ratio lie between the two domestic ratios?}\par
\begin{itemize}
\item Cameroon can produce 1 kg of cocoa at a domestic cost of 0.5 m of cloth. It will only \emph{sell} cocoa abroad if it receives \emph{more} than 0.5 m of cloth per kg; otherwise it would rather make its own cloth.
\item Nigeria can produce 1 m of cloth at a domestic cost of 0.5 kg of cocoa, i.e.\ it can obtain cocoa domestically at 2 m of cloth per kg. It will only \emph{buy} cocoa if it pays \emph{less} than 2 m of cloth per kg.
\item Hence any ratio with $0.5 < \text{cloth per kg of cocoa} < 2$ leaves \emph{both} countries better off than self-sufficiency.
\end{itemize}

\begin{popfigure}
\begin{center}
\begin{tikzpicture}
\begin{axis}[
    width=12cm, height=7.5cm,
    xlabel={Cocoa (kg)},
    ylabel={Cloth (m)},
    xmin=0, xmax=1100, ymin=0, ymax=2300,
    xtick={0,200,400,600,800,1000},
    ytick={0,500,1000,1500,2000},
    legend style={at={(0.02,0.98)}, anchor=north west, font=\small},
    grid=both, grid style={popline!40},
    axis lines=left,
]
\addplot[popblue, very thick, domain=0:1000, name path=A] {0.5*x};
\addlegendentry{Cameroon's ratio: 1 kg cocoa = 0.5 m cloth}
\addplot[poporange, very thick, domain=0:1000, name path=B] {2*x};
\addlegendentry{Nigeria's ratio: 1 kg cocoa = 2 m cloth}
\addplot[popgreen, very thick, dashed, domain=0:1000] {1*x};
\addlegendentry{Possible terms of exchange: 1 kg = 1 m}
\addplot[popgreenL, opacity=0.3] fill between[of=A and B];
\end{axis}
\end{tikzpicture}
\end{center}
\legende{The feasible range of the international exchange ratio: any exchange line lying between the two domestic cost-ratio lines (shaded region) makes trade beneficial to both countries.}
\end{popfigure}

The exact ratio \emph{within} this range is determined by the \textbf{relative strength of demand} in the two countries (reciprocal demand): if world demand for cocoa is strong relative to cloth, the ratio settles near Nigeria's limit (1 kg cocoa $\approx$ 2 m cloth), favouring Cameroon; if demand for cloth is stronger, it settles near Cameroon's limit, favouring Nigeria.

\courssub{Sharing the Gains from Trade}

Suppose the two countries agree on an exchange ratio of \textbf{1 kg of cocoa = 1 m of cloth} (midway between 0.5 and 2). Cameroon exports 400 kg of cocoa and receives 400 m of cloth in return.

\begin{center}
\begin{tabularx}{\linewidth}{|l|X|X|X|X|}
\hline
\rowcolor{popblueL} \textbf{Country} & \textbf{Cocoa kept (kg)} & \textbf{Cloth obtained (m)} & \textbf{Before trade} & \textbf{Gain} \\
\hline
Cameroon & $1000 - 400 = 600$ & 400 & 500 kg, 250 m & $+100$ kg, $+150$ m \\
\hline
Nigeria & 400 & $1200 - 400 = 800$ & 300 kg, 600 m & $+100$ kg, $+200$ m \\
\hline
\end{tabularx}
\end{center}

\emph{Both} countries end up with more of \emph{both} goods than before specialisation. This is the rigorous meaning of the statement that ``trade is mutually beneficial''.

\begin{aretenir}[Key Points]
\begin{itemize}
\item The closer the exchange ratio lies to a country's own domestic ratio, the \emph{smaller} that country's share of the gain.
\item The country whose export good is in stronger world demand captures the larger share of the gains.
\item The sum of the two countries' gains equals the increase in world output created by specialisation.
\end{itemize}
\end{aretenir}

\begin{methode}[Computing World Output Gains and the Exchange Ratio]
\begin{enumerate}
\item \textbf{Read the productivity table} and compute each country's internal opportunity cost ratio for each good.
\item \textbf{Assign specialisation}: each country specialises in the good for which its opportunity cost is \emph{lower} (comparative advantage).
\item \textbf{Compute world output before} specialisation (sum of both countries' outputs of each good).
\item \textbf{Compute world output after} complete specialisation (each country devotes all its labour to its speciality).
\item \textbf{Compute the gain}: world output after $-$ world output before, for each good.
\item \textbf{Set the limits of the exchange ratio}: it must lie strictly between the two domestic ratios.
\item \textbf{Choose a ratio within the limits}, fix a quantity exported, and compute each country's final consumption bundle.
\item \textbf{Verify} that each country ends up with at least as much of both goods as before trade.
\end{enumerate}
\end{methode}

\begin{attention}[Common Mistake: A Ratio Outside the Two Domestic Ratios]
Many candidates propose an exchange ratio that lies \emph{outside} the two domestic cost ratios (e.g.\ 1 kg cocoa $= 3$ m cloth in our example). At such a ratio, one country (here Nigeria, which can obtain cocoa domestically at only 2 m per kg) would \emph{lose} from trade and would refuse to exchange. Always check: \textbf{the ratio must lie strictly between the two internal ratios}. Another frequent error is to confuse the direction of specialisation: always compare \emph{opportunity costs}, never absolute outputs.
\end{attention}

\begin{exoresolu}[Practical Work 24: Determining the Exchange of Goods after Specialisation]
One unit of labour produces per day: in country A, 8 tonnes of yams or 4 tonnes of rice; in country B, 3 tonnes of yams or 6 tonnes of rice. Each country has 120 labour units and, before specialisation, uses 60 units on each crop. (a) Determine each country's comparative advantage. (b) Compute world output before and after complete specialisation. (c) State the limits of the exchange ratio. (d) If the agreed ratio is 1 tonne of yams $= 1$ tonne of rice and A exports 300 tonnes of yams, compute each country's final consumption and its gain.
\tcblower
\textbf{(a) Comparative advantage.} Country A: 8 yams cost 4 rice, so 1 tonne of yams $= 0.5$ tonne of rice. Country B: 6 rice cost 3 yams, so 1 tonne of yams $= 2$ tonnes of rice. Yams are cheaper in A ($0.5 < 2$): A specialises in \textbf{yams}; B specialises in \textbf{rice}.

\textbf{(b) World output.}
Before: A produces $60 \times 8 = 480$ yams and $60 \times 4 = 240$ rice; B produces $60 \times 3 = 180$ yams and $60 \times 6 = 360$ rice. World: \textbf{660 yams, 600 rice}.
After complete specialisation: A produces $120 \times 8 = 960$ yams, 0 rice; B produces 0 yams, $120 \times 6 = 720$ rice. World: \textbf{960 yams, 720 rice}.
Gain: $+300$ tonnes of yams and $+120$ tonnes of rice.

\textbf{(c) Limits of the exchange ratio.} $0.5 < \text{rice per tonne of yams} < 2$.

\textbf{(d) Exchange at 1 yam $= 1$ rice, A exports 300 tonnes of yams.}
A keeps $960 - 300 = 660$ yams and receives 300 rice. Before trade A had 480 yams and 240 rice: gain $= +180$ yams, $+60$ rice.
B keeps $720 - 300 = 420$ rice and receives 300 yams. Before trade B had 180 yams and 360 rice: gain $= +120$ yams, $+60$ rice.
Both countries gain in both goods: the exchange is mutually beneficial, as required.
\end{exoresolu}

\begin{methode}[Exam Tip: Method Marks]
In the GCE examination, marks are awarded for \emph{each step} of the calculation, not only for the final answer. Always: (1) write down the domestic cost ratios explicitly; (2) show the before/after output table; (3) state the two limits of the exchange ratio; (4) show the final consumption bundles. A correct final figure with no working earns only a fraction of the marks.
\end{methode}

\courssub{Real-World Qualifications}

The neat conclusion above rests on simplifying assumptions. In reality, three factors push the \emph{actual} exchange ratio and may shrink the gains:

\begin{itemize}
\item \textbf{Transport costs.} Shipping cocoa from Douala to Lagos or Rotterdam is costly. Transport costs narrow the feasible band between the two domestic ratios; if they are large enough, they can wipe out the gain entirely and make trade unprofitable.
\item \textbf{Tariffs and other trade restrictions.} An import duty raises the price paid by the importing country, moving the effective exchange ratio against the exporter and reducing the volume of mutually beneficial trade (studied in detail in the next section).
\item \textbf{Demand conditions.} The exact ratio within the feasible band depends on the intensity of each country's demand for the other's product. A small country trading with a large one often finds the ratio settling near its own domestic ratio, so the larger partner captures most of the gain.
\item \textbf{Increasing costs and partial specialisation.} With increasing opportunity costs, countries stop specialising before the ``complete'' point, so the actual gain in world output is smaller than the textbook maximum.
\end{itemize}

\begin{savaistu}[Cameroon and the Gains from Specialisation]
Cameroon's exports are heavily concentrated in a few primary products --- crude oil, cocoa, coffee, cotton, timber and bananas --- in which the country is presumed to hold a comparative advantage, while it imports manufactured goods. This pattern fits the theory of specialisation, but it also illustrates its risks: when world prices of cocoa or oil fall, Cameroon's export earnings drop sharply, which is why economists recommend \emph{diversifying} the economy rather than specialising completely.
\end{savaistu}

\begin{exercice}[Practice]
In country X, one labour unit produces 12 bags of maize or 4 bags of cement; in country Y, 5 bags of maize or 10 bags of cement. Each country has 200 labour units, split equally before specialisation. (a) Determine the comparative advantages. (b) Compute world output before and after complete specialisation and the gain. (c) State the limits of the exchange ratio between maize and cement. (d) At a ratio of 1 bag of maize $= 1$ bag of cement, with X exporting 600 bags of maize, compute each country's final consumption and gain. (e) Explain what would happen to trade if the ratio were fixed at 1 bag of maize $= 0.2$ bag of cement.
\end{exercice}

\begin{corrige}[Practice]
(a) X: 1 maize $= \tfrac{1}{3}$ cement; Y: 1 maize $= 2$ cement. Maize is cheaper in X: X specialises in maize, Y in cement.
(b) Before: X: $100 \times 12 = 1200$ maize, $100 \times 4 = 400$ cement; Y: $100 \times 5 = 500$ maize, $100 \times 10 = 1000$ cement. World: 1700 maize, 1400 cement. After: X: 2400 maize; Y: 2000 cement. Gain: $+700$ maize, $+600$ cement.
(c) $\tfrac{1}{3} < \text{cement per bag of maize} < 2$.
(d) X keeps $2400 - 600 = 1800$ maize and receives 600 cement: gain $+600$ maize, $+200$ cement. Y keeps $2000 - 600 = 1400$ cement and receives 600 maize: gain $+100$ maize, $+400$ cement.
(e) At 1 maize $= 0.2$ cement, the ratio lies \emph{below} X's domestic cost ($\tfrac{1}{3}$): X would receive less cement per bag of maize than it costs to produce cement itself, so X would refuse to trade. No mutually beneficial exchange occurs.
\end{corrige}

\begin{aretenir}[To Conclude]
Specialisation according to comparative advantage raises total world output; the terms of exchange, fixed between the two domestic opportunity cost ratios, determine how this gain is shared; and transport costs, tariffs and demand conditions explain why real-world gains are smaller and unevenly distributed. Mastering the numerical demonstration --- output table, cost ratios, exchange limits, final bundles --- is essential, as it is one of the most frequently examined calculations in the international trade paper.
\end{aretenir}

\coursec{International Trade Restrictions}

In the previous section we established a powerful result: when countries specialise according to comparative advantage and trade freely, world output rises and \emph{both} trading partners can gain. If free trade is so beneficial, a natural question arises: why do almost all governments in the world --- including Cameroon and the CEMAC zone --- still restrict trade in one way or another? This section answers that question. We define free trade and protectionism, study the main instruments of protection (tariffs, quotas and other barriers), analyse their effects rigorously with diagrams, and evaluate the arguments for and against protection.

\courssub{Free Trade and Protectionism: Meaning}

\begin{definition}[Free trade and protectionism]
\cle{Free trade} is a situation in which goods and services move between countries without any government-imposed barriers such as taxes, quantitative limits or administrative obstacles. Prices and quantities traded are determined purely by demand and supply on the world market.

\cle{Protectionism} is the deliberate policy of restricting international trade in order to shield domestic producers from foreign competition. It operates through instruments such as tariffs, quotas, embargoes, subsidies and administrative barriers.
\end{definition}

In practice, no country is completely free-trading and none is completely closed. Every economy chooses a \emph{degree} of openness. Cameroon, for example, trades freely within CEMAC but applies a common external tariff on imports from outside the zone.

\begin{savaistu}[Cameroon's trade openness]
Cameroon's trade-to-GDP ratio has historically hovered around 40--50\,\%, reflecting its integration into regional (CEMAC) and global markets. The country is a member of the WTO and has signed the African Continental Free Trade Area (AfCFTA) agreement, committing to progressive tariff liberalisation on intra-African trade while maintaining protective tariffs on sensitive products from outside the continent.
\end{savaistu}

\courssub{Tariffs}

\begin{definition}[Tariff]
A \cle{tariff} is a tax imposed on imported goods (and occasionally on exports). It raises the domestic price of the imported product above the world price, making foreign goods less competitive relative to domestic production. Two main types exist:
\begin{itemize}
\item A \cle{specific tariff} is a fixed amount of money per physical unit imported, e.g.\ 5\,000 FCFA per tonne of imported rice.
\item An \cle{ad valorem tariff} is a fixed percentage of the value of the import, e.g.\ 20\,\% of the customs value of an imported vehicle.
\end{itemize}
\end{definition}

\textbf{Effects of a tariff: graphical analysis.}\par
Consider the market for rice in Cameroon. Without trade, the domestic price would be set where domestic demand equals domestic supply. With free trade at the world price $P_w$ (below the autarky price), domestic consumers demand $Q_2$, domestic producers supply $Q_1$, and imports fill the gap $(Q_2 - Q_1)$. Now suppose the government imposes a specific tariff $t$ per unit. The domestic price rises to $P_t = P_w + t$. The consequences follow mechanically:

\begin{itemize}
\item \textbf{Price:} rises from $P_w$ to $P_t$.
\item \textbf{Consumption:} falls from $Q_2$ to $Q_4$ (movement up along the demand curve).
\item \textbf{Domestic production:} rises from $Q_1$ to $Q_3$ (movement up along the supply curve) --- this is the \emph{protective effect}.
\item \textbf{Imports:} shrink from $(Q_2 - Q_1)$ to $(Q_4 - Q_3)$.
\item \textbf{Government revenue:} equals the tariff per unit multiplied by the new volume of imports: $t \times (Q_4 - Q_3)$, shown as the shaded rectangle.
\item \textbf{Welfare:} consumers lose; producers gain; part of the consumer loss is a pure deadweight loss to society.
\end{itemize}

\begin{popfigure}
\begin{center}
\begin{tikzpicture}[scale=1.05]
\draw[->, thick] (0,0) -- (8.2,0) node[below left, font=\small] {Quantity of rice};
\draw[->, thick] (0,0) -- (0,5.6) node[above, font=\small] {Price};
% Demand and supply
\draw[very thick, popblue] (0.6,5.0) -- (7.4,0.6) node[pos=0.97, above, font=\small, popblue] {$D$};
\draw[very thick, poporange] (0.6,0.8) -- (7.4,5.2) node[pos=0.97, above, font=\small, poporange] {$S$};
% World price and tariff price
\draw[dashed, popmuted] (0,1.9) -- (7.4,1.9);
\draw[dashed, popmuted] (0,2.9) -- (7.4,2.9);
\node[left, font=\small] at (0,1.9) {$P_w$};
\node[left, font=\small] at (0,2.9) {$P_t$};
% Quantities: Q1=2.1, Q3=3.3, Q4=4.9, Q2=6.1
\draw[dashed, popmuted] (2.1,0) -- (2.1,1.9);
\draw[dashed, popmuted] (3.3,0) -- (3.3,2.9);
\draw[dashed, popmuted] (4.9,0) -- (4.9,2.9);
\draw[dashed, popmuted] (6.1,0) -- (6.1,1.9);
\node[below, font=\small] at (2.1,0) {$Q_1$};
\node[below, font=\small] at (3.3,0) {$Q_3$};
\node[below, font=\small] at (4.9,0) {$Q_4$};
\node[below, font=\small] at (6.1,0) {$Q_2$};
% Revenue rectangle
\fill[popgold, opacity=0.45] (3.3,1.9) rectangle (4.9,2.9);
\node[font=\footnotesize, align=center] at (4.1,2.4) {Govt Revenue};
% tariff bracket
\draw[<->, thick, popdark] (7.0,1.9) -- (7.0,2.9) node[midway, right, font=\small, popdark] {$t$};
% import braces
\draw[<->, thick, popgreen] (2.1,1.55) -- (6.1,1.55) node[midway, below, font=\footnotesize, popgreen] {Imports before tariff};
\draw[<->, thick, poppurple] (3.3,3.25) -- (4.9,3.25) node[midway, above, font=\footnotesize, poppurple] {Imports after tariff};
% Deadweight loss triangles (optional, for advanced students)
\draw[popred, thick, fill=popred, fill opacity=0.1] (2.1,1.9) -- (3.3,2.9) -- (2.1,2.9) -- cycle;
\draw[popred, thick, fill=popred, fill opacity=0.1] (4.9,1.9) -- (6.1,1.9) -- (4.9,2.9) -- cycle;
\node[font=\tiny, popred] at (1.5,2.4) {DWL};
\node[font=\tiny, popred] at (5.5,2.4) {DWL};
\end{tikzpicture}
\end{center}
\legende{Effect of a tariff on the domestic market: price rises from $P_w$ to $P_t$, domestic output expands from $Q_1$ to $Q_3$, consumption contracts from $Q_2$ to $Q_4$, imports fall, and the government collects the shaded revenue rectangle $t \times (Q_4 - Q_3)$. The two small red triangles represent the net welfare loss (deadweight loss) to society.}
\end{popfigure}

\begin{exemplebox}[Numerical illustration of a specific tariff]
Suppose the world price of rice is $P_w = 200$ FCFA per kg. At this price, Cameroonian consumers demand 100\,000 tonnes and domestic farmers supply 40\,000 tonnes, so imports are 60\,000 tonnes. The government imposes a specific tariff of 50 FCFA per kg, raising the domestic price to 250 FCFA. Demand falls to 90\,000 tonnes and domestic supply rises to 55\,000 tonnes. Imports fall to 35\,000 tonnes, and government revenue is $50 \times 35\,000\,000 = 1.75$ billion FCFA. Consumers pay more, farmers sell more, and the Treasury gains revenue --- but society as a whole loses efficiency.
\end{exemplebox}

\begin{exemplebox}[Numerical illustration of an ad valorem tariff]
Suppose the world price of a used vehicle is 5\,000\,000 FCFA. The government imposes a 30\,\% ad valorem tariff. The domestic price becomes $5\,000\,000 \times 1.30 = 6\,500\,000$ FCFA. If 2\,000 vehicles are imported at this price, government revenue is $0.30 \times 5\,000\,000 \times 2\,000 = 3$ billion FCFA. Note: for an ad valorem tariff, the tariff \emph{per unit} rises with the world price, providing automatic protection when world prices fall.
\end{exemplebox}

\begin{methode}[Step-by-step: analysing a tariff in an exam]
Whenever a question asks you to ``analyse the effects of a tariff'', follow these steps:
\begin{enumerate}
\item Draw the domestic demand ($D$) and supply ($S$) curves. Mark the world price $P_w$ as a horizontal line below the autarky equilibrium.
\item Identify free-trade consumption ($Q_2$ on $D$), domestic production ($Q_1$ on $S$) and imports ($Q_2 - Q_1$).
\item Draw the new domestic price line $P_t = P_w + t$ (specific tariff) or $P_t = P_w(1 + \text{rate})$ (ad valorem).
\item Read off the new consumption ($Q_4$ on $D$), domestic production ($Q_3$ on $S$) and imports ($Q_4 - Q_3$).
\item Shade and label the government revenue rectangle: height $= t$ (or $P_t - P_w$), width $= Q_4 - Q_3$.
\item Identify the two deadweight loss triangles (production inefficiency and consumption inefficiency).
\item Conclude with a written commentary: who gains (producers, government), who loses (consumers), and the net welfare loss.
\end{enumerate}
A fully labelled diagram with a clear written commentary typically earns the majority of the marks.
\end{methode}

\courssub{Quotas}

\begin{definition}[Quota]
A \cle{quota} is a physical limit on the quantity (or value) of a good that may be imported during a given period, e.g.\ no more than 20\,000 tonnes of frozen chicken per year. Once the quota is filled, no further imports are allowed, whatever the price.
\end{definition}

\textbf{Effects of a quota.}\par
A quota restricts supply. At the world price $P_w$, domestic supply is $Q_1$ and demand is $Q_2$; free-trade imports would be $(Q_2 - Q_1)$. If the government allows only a quota $\bar{M} < (Q_2 - Q_1)$, total supply on the domestic market becomes domestic supply plus the fixed quota. Scarcity pushes the domestic price up to $P_q$, where the gap between demand and domestic supply is exactly equal to the quota. Consumption falls, domestic production rises, and imports are compressed to $\bar{M}$.

\begin{popfigure}
\begin{center}
\begin{tikzpicture}[scale=1.05]
\draw[->, thick] (0,0) -- (8.2,0) node[below left, font=\small] {Quantity};
\draw[->, thick] (0,0) -- (0,5.6) node[above, font=\small] {Price};
\draw[very thick, popblue] (0.6,5.0) -- (7.4,0.6) node[pos=0.97, above, font=\small, popblue] {$D$};
\draw[very thick, poporange] (0.6,0.8) -- (7.4,5.2) node[pos=0.97, above, font=\small, poporange] {$S$};
% S + quota: supply shifted right by quota amount (1.6 units)
\draw[very thick, popgreen, dashed] (2.2,0.8) -- (8.0,4.6) node[pos=0.9, above, font=\small, popgreen] {$S + \text{quota}$};
% prices
\draw[dashed, popmuted] (0,1.9) -- (7.4,1.9);
\draw[dashed, popmuted] (0,2.7) -- (7.4,2.7);
\node[left, font=\small] at (0,1.9) {$P_w$};
\node[left, font=\small] at (0,2.7) {$P_q$};
% quantities: Q1=2.1 (S at Pw), Q3=3.0 (S at Pq), Q4=4.6 (S+quota at Pq and D at Pq), Q2=6.1 (D at Pw)
\draw[dashed, popmuted] (2.1,0) -- (2.1,1.9);
\draw[dashed, popmuted] (3.0,0) -- (3.0,2.7);
\draw[dashed, popmuted] (4.6,0) -- (4.6,2.7);
\draw[dashed, popmuted] (6.1,0) -- (6.1,1.9);
\node[below, font=\small] at (2.1,0) {$Q_1$};
\node[below, font=\small] at (3.0,0) {$Q_3$};
\node[below, font=\small] at (4.6,0) {$Q_4$};
\node[below, font=\small] at (6.1,0) {$Q_2$};
% quota brace
\draw[<->, thick, poppurple] (3.0,3.05) -- (4.6,3.05) node[midway, above, font=\footnotesize, poppurple] {Quota $\bar{M}$};
\draw[<->, thick, popgreen] (2.1,1.55) -- (6.1,1.55) node[midway, below, font=\footnotesize, popgreen] {Free-trade imports};
% Quota rent rectangle
\fill[poppink, opacity=0.4] (3.0,1.9) rectangle (4.6,2.7);
\node[font=\footnotesize, align=center] at (3.8,2.3) {Quota Rent};
\end{tikzpicture}
\end{center}
\legende{Effect of an import quota: limiting imports to $\bar{M}$ shifts total supply to $S + \text{quota}$, raising the domestic price from $P_w$ to $P_q$; consumption falls from $Q_2$ to $Q_4$ and domestic output rises from $Q_1$ to $Q_3$. The shaded rectangle represents the quota rent (extra profit per unit $\times$ quota volume).}
\end{popfigure}

\begin{attention}[Common mistake: tariff revenue versus quota revenue]
A \cle{tariff} always raises government revenue equal to the tariff per unit times the volume of imports that still enter. A \cle{quota} raises \textbf{no} government revenue: the extra profit created by the higher domestic price (the \emph{quota rent}) is captured by whoever holds the import licences --- usually foreign exporters or domestic importers. The government only obtains revenue from a quota if it \emph{sells or auctions} the import licences. Do not write in an examination that ``a quota increases government revenue'' --- this is a frequent and costly error.
\end{attention}

\textbf{Tariff versus quota: a comparison.}\par
\begin{center}
\begin{tabularx}{\linewidth}{|l|X|X|}
\hline
\rowcolor{popblueL} \textbf{Criterion} & \textbf{Tariff} & \textbf{Quota} \\
\hline
Instrument & Tax on imports & Physical limit on imports \\
\hline
Price effect & Raises price by the amount of the tariff & Raises price through artificial scarcity \\
\hline
Government revenue & Yes: $t \times$ imports & No (unless licences are sold/auctioned) \\
\hline
Certainty of import volume & Uncertain (depends on elasticities) & Certain: imports cannot exceed the quota \\
\hline
Effect of rising demand & Imports rise, domestic price unchanged & Domestic price rises, imports fixed at quota \\
\hline
Effect of falling world price & Domestic price falls, imports may rise & Domestic price unchanged, imports fixed \\
\hline
Administrative cost & Relatively low (customs collection) & Higher (licence allocation, monitoring) \\
\hline
\end{tabularx}
\end{center}

\courssub{Other Instruments of Protection}

\begin{definition}[Embargo]
An \cle{embargo} is a complete prohibition of trade in a particular good, or with a particular country, usually imposed for political or strategic reasons. It is the most extreme form of trade restriction --- a quota reduced to zero.
\end{definition}

Other widely used instruments include:
\begin{itemize}
\item \cle{Exchange controls}: the government limits access to foreign currency, so importers cannot obtain the FCFA--foreign currency conversion needed to pay for imports. This restricts imports indirectly.
\item \cle{Subsidies to domestic producers}: rather than taxing foreign goods, the government lowers the costs of home producers (e.g.\ subsidised fertiliser for Cameroonian farmers), allowing them to compete with cheaper imports without raising consumer prices.
\item \cle{Voluntary export restraints (VERs)}: the exporting country ``voluntarily'' agrees to limit its sales, usually under pressure from the importing country (a classic example is Japanese car exports to the USA in the 1980s).
\item \cle{Administrative and technical barriers}: complex customs procedures, strict health and safety standards, labelling rules or lengthy inspections that discourage imports without any explicit tax or quota. These are sometimes called \emph{non-tariff barriers} (NTBs).
\item \cle{Local content requirements}: laws requiring that a certain percentage of a final product be produced domestically.
\item \cle{Anti-dumping duties}: special tariffs imposed when foreign firms sell below cost or below home-market price (see definition of dumping below).
\end{itemize}

\courssub{Arguments For Protection}

Why do governments protect, despite the gains from free trade? The main arguments are:

\begin{enumerate}
\item \textbf{Infant industry argument.} New industries (e.g.\ Cameroonian agro-processing or cement production) cannot immediately match the prices of established foreign competitors. Temporary protection gives them time to grow, learn and reach economies of scale, after which protection can be removed.
\item \textbf{Employment protection.} Protecting domestic industries preserves jobs that would otherwise be lost to imports --- politically powerful in regions dependent on a single industry.
\item \textbf{Balance of payments correction.} Import restrictions reduce the import bill and can help correct a current account deficit (we return to this in the balance of payments lessons).
\item \textbf{Anti-dumping / unfair competition.}
\end{enumerate}

\begin{definition}[Dumping]
\cle{Dumping} is the practice of selling a good on a foreign market at a price below its cost of production, or below the price charged in the exporter's home market, usually to drive out local competitors and gain market dominance. Tariffs imposed to offset dumping are called \emph{anti-dumping duties} and are generally accepted as legitimate under WTO rules.
\end{definition}

\begin{enumerate}
\setcounter{enumi}{4}
\item \textbf{National security and strategic self-sufficiency.} A country may protect food production, energy or defence-related industries so as not to depend on foreigners in times of crisis.
\item \textbf{Government revenue.} For developing countries like Cameroon, where income tax collection is difficult, import duties are a relatively easy and important source of budget revenue.
\item \textbf{Protection of cultural identity / social standards.} Some countries restrict imports that conflict with local cultural values or that are produced under labour/environmental standards deemed unacceptable.
\end{enumerate}

\courssub{Arguments Against Protection (For Free Trade)}

\begin{enumerate}
\item \textbf{Higher prices and reduced consumer welfare.} Tariffs and quotas raise domestic prices, hurting consumers --- especially poor households that spend a large share of income on imported food and fuel.
\item \textbf{Retaliation and trade wars.} One country's protection invites retaliation by its partners, shrinking everyone's exports and possibly destroying the gains from trade altogether.
\item \textbf{Inefficiency and complacency.} Sheltered from competition, protected firms may remain high-cost and low-quality; infant industries sometimes refuse to grow up, and ``temporary'' protection becomes permanent.
\item \textbf{Reduced choice and slower growth.} Consumers enjoy fewer varieties, and the economy forgoes the efficiency gains of specialisation demonstrated by comparative advantage theory.
\item \textbf{Corruption and rent-seeking.} Quotas and licences create opportunities for favouritism and bribery in their allocation.
\item \textbf{Deadweight loss.} The net welfare loss (production and consumption inefficiency) reduces total economic surplus.
\end{enumerate}

\begin{experience}[Case study: Cameroon's poultry sector]
In the early 2000s, Cameroon faced a surge in cheap frozen chicken imports (mainly from the EU and Brazil) that threatened local poultry farmers. The government responded with a combination of measures: (1) a temporary ban on frozen chicken imports (an embargo), (2) high tariffs on poultry products under the CET, and (3) support for local feed production. Local production subsequently expanded, but consumers faced higher prices. This illustrates the trade-off: protection saved jobs and an infant industry, but at the cost of higher food prices for urban households. The WTO later ruled that some of these measures needed to be time-bound and transparent.
\end{experience}

\courssub{Protectionism in Cameroon and CEMAC}

Cameroon applies protection mainly through the \cle{Common External Tariff (CET)} of the CEMAC customs union. Under the CET, goods entering any CEMAC member state from outside the zone face a common tariff structure, historically organised in categories rising from low rates on essential goods and raw materials (Category 1: 0--5\,\%) to the highest rates on finished consumer goods (Category 4: up to 30\,\%). The logic is twofold: to protect regional industries (cement, beverages, agro-processing, textiles) from extra-regional competition, and to generate customs revenue, which is a major source of government income for CEMAC states. Within the zone, goods in principle circulate free of tariffs, although non-tariff obstacles (roadblocks, differing standards, administrative delays) often persist in practice. Cameroon has also, at times, used quantitative restrictions and temporary bans --- for example on certain imported frozen products --- to protect local producers, and it participates in trade arrangements such as the Economic Partnership Agreement (EPA) with the European Union, which progressively reshapes its tariff structure.

\begin{savaistu}[CEMAC CET categories (simplified)]
\begin{itemize}
\item Category 1 (0--5\,\%): Essential goods, raw materials, capital goods not produced locally.
\item Category 2 (10\,\%): Intermediate goods, some capital goods.
\item Category 3 (20\,\%): Finished goods with some local production.
\item Category 4 (30\,\%): Finished consumer goods competing strongly with local industry.
\end{itemize}
The exact rates and product classifications are updated periodically by CEMAC regulations.
\end{savaistu}

\begin{exoresolu}[Worked exercise: analysing a tariff]
\textbf{Statement.} The world price of sugar is 300 FCFA per kg. At this price, demand in a country is 80\,000 tonnes and domestic supply is 30\,000 tonnes. The government introduces a specific tariff of 60 FCFA per kg. As a result, demand falls to 70\,000 tonnes and domestic supply rises to 45\,000 tonnes.
\begin{enumerate}
\item Calculate imports before and after the tariff.
\item Calculate the government revenue from the tariff.
\item Identify the gainers and losers.
\end{enumerate}
\tcblower
\textbf{Solution.}
\begin{enumerate}
\item Before the tariff: imports $= 80\,000 - 30\,000 = 50\,000$ tonnes. After the tariff: imports $= 70\,000 - 45\,000 = 25\,000$ tonnes. The tariff has halved imports.
\item Revenue $=$ tariff per unit $\times$ imports after tariff $= 60 \times 25\,000\,000 = 1.5$ billion FCFA.
\item \textbf{Gainers:} domestic sugar producers, who sell 15\,000 more tonnes at a price 60 FCFA higher; and the government, which collects 1.5 billion FCFA. \textbf{Losers:} consumers, who pay 360 instead of 300 FCFA per kg and consume 10\,000 tonnes less. Because consumer losses exceed the combined gains of producers and the Treasury, the country suffers a net welfare (deadweight) loss --- the cost of protection.
\end{enumerate}
\end{exoresolu}

\begin{exoresolu}[Worked exercise: analysing a quota]
\textbf{Statement.} The world price of wheat is 150 FCFA/kg. At this price, domestic demand is 500\,000 tonnes and domestic supply is 100\,000 tonnes. The government imposes an import quota of 200\,000 tonnes. The resulting domestic price rises to 190 FCFA/kg. At this price, domestic supply becomes 150\,000 tonnes and demand becomes 350\,000 tonnes.
\begin{enumerate}
\item Verify that the quota is binding.
\item Calculate the quota rent per unit and total quota rent.
\item Who captures the quota rent if licences are given free to established importers?
\end{enumerate}
\tcblower
\textbf{Solution.}
\begin{enumerate}
\item Free-trade imports would be $500\,000 - 100\,000 = 400\,000$ tonnes. The quota of 200\,000 tonnes is lower, so it is binding.
\item Quota rent per unit $= P_q - P_w = 190 - 150 = 40$ FCFA/kg. Total quota rent $= 40 \times 200\,000\,000 = 8$ billion FCFA.
\item If licences are given free to established importers, they capture the entire 8 billion FCFA as extra profit. The government gets zero revenue unless it auctions the licences.
\end{enumerate}
\end{exoresolu}

\begin{exercice}[Practice]
Using a clearly labelled demand-and-supply diagram, analyse the effects on price, domestic production, consumption, imports and government revenue of (a) an ad valorem tariff on imported second-hand clothing, and (b) an import quota on the same product. In each case, state who gains and who loses, and give one reason why a CEMAC government might nevertheless choose protection.
\end{exercice}

\begin{corrige}[Exercise 1]
(a) \textbf{Tariff:} the domestic price rises from $P_w$ to $P_w(1 + \text{rate})$; domestic supply expands along the $S$ curve; demand contracts along the $D$ curve; imports shrink; government revenue equals $(P_t - P_w) \times$ new import volume (shade the rectangle). Gainers: domestic producers and the Treasury. Losers: consumers; net deadweight loss for society.
(b) \textbf{Quota:} total supply becomes $S + \text{quota}$; the price rises until demand minus domestic supply equals the quota; production rises and consumption falls as with a tariff, but there is \emph{no} government revenue --- the quota rent goes to licence holders.
Reason for protecting: e.g.\ to defend an infant industry, preserve employment, reduce the import bill (balance of payments), or raise customs revenue.
\end{corrige}

\begin{aretenir}[Key points]
\begin{itemize}
\item Free trade means no government barriers; protectionism means deliberate restriction of trade to shield domestic producers.
\item A tariff is a tax on imports (specific or ad valorem): it raises price, reduces consumption and imports, expands domestic output, and yields government revenue.
\item A quota is a physical limit on imports: similar price and quantity effects, but \emph{no} revenue unless licences are sold; the quota rent accrues to licence holders.
\item Other instruments: embargoes, exchange controls, subsidies, VERs, administrative barriers, local content requirements, anti-dumping duties.
\item Protection is defended by the infant industry, employment, balance of payments, anti-dumping, security, revenue and cultural arguments; it is criticised for raising prices, provoking retaliation, breeding inefficiency, limiting choice, corruption and deadweight loss.
\item Cameroon protects through the CEMAC Common External Tariff and occasional quantitative restrictions.
\end{itemize}
\end{aretenir}

\coursec{Terms of Trade}

In the previous section we examined how governments restrict international trade through tariffs, quotas and other instruments. Those policies matter because countries care deeply about \emph{what they receive in exchange for what they sell}. This naturally raises a question: is trade becoming ``cheaper'' or ``dearer'' for a country over time? When Cameroon sells a tonne of cocoa, how many tonnes of imported machinery can that tonne buy this year compared with last year? The concept that answers this question is the \cle{terms of trade}, one of the most frequently examined notions in international economics.

\courssub{Definition and Meaning of the Terms of Trade}

\begin{definition}[Terms of Trade (Commodity or Net Barter Terms of Trade)]
The \cle{terms of trade} (TOT) of a country measure the relationship between the prices of its \textbf{exports} and the prices of its \textbf{imports}. More precisely, the terms of trade are the ratio of the index of export prices to the index of import prices, expressed as a percentage. They indicate the \textbf{volume of imports that can be purchased with a given volume of exports}.
\end{definition}

The intuition is simple. Imagine a cocoa farmer in the Centre Region. His ``terms of trade'' with the town are favourable when the price of a bag of cocoa rises while the price of the fertiliser, soap and clothes he buys stays constant: each bag now buys more goods. If instead cocoa prices fall while the prices of manufactured goods rise, each bag buys less --- his exchange position has deteriorated, even if he works exactly as hard as before. A country is in the same situation: Cameroon exports cocoa, coffee, timber, cotton and crude oil, and imports machinery, vehicles, pharmaceuticals and refined petroleum products. The terms of trade tell us whether the ``exchange rate'' between what Cameroon sells and what Cameroon buys is moving in its favour or against it.

\begin{propriete}[Formula of the Terms of Trade]
\[
\text{TOT} = \frac{\text{Index of export prices}}{\text{Index of import prices}} \times 100 = \frac{P_x}{P_m} \times 100
\]
where $P_x$ is the index of export prices and $P_m$ is the index of import prices, both calculated with the same base year (base $= 100$).
\end{propriete}

\begin{attention}[Use indices, not absolute prices]
A very common mistake is to divide the \emph{price of one exported good} by the \emph{price of one imported good}. The terms of trade are built from \textbf{price indices} covering the whole basket of exports and imports, with a common base year. Also note the direction: export prices on top, import prices at the bottom. Inverting them reverses the interpretation entirely.
\end{attention}

\courssub{Calculating the Terms of Trade: A Worked Example}

Suppose we have the following data for a hypothetical economy, with year 1 as the base year.

\begin{center}
\begin{tabularx}{\linewidth}{|l|X|X|X|}
\hline
\rowcolor{popblueL} \textbf{Year} & \textbf{Export price index} $P_x$ & \textbf{Import price index} $P_m$ & \textbf{TOT} $= (P_x/P_m)\times 100$ \\
\hline
1 & 100 & 100 & 100.0 \\
\hline
2 & 108 & 102 & 105.9 \\
\hline
3 & 112 & 110 & 101.8 \\
\hline
4 & 105 & 115 & 91.3 \\
\hline
5 & 98 & 118 & 83.1 \\
\hline
\end{tabularx}
\end{center}

\begin{exemplebox}[Step-by-step calculation for Year 4]
\[
\text{TOT}_4 = \frac{105}{115} \times 100 = 91.3
\]
\textbf{Interpretation:} in year 4, a given volume of exports purchases only about $91.3\%$ of the volume of imports it could purchase in the base year. The country's exchange position has \emph{deteriorated} by roughly $8.7\%$ relative to the base year.
\end{exemplebox}

\begin{exoresolu}[Calculating and interpreting the TOT]
In 2020 (base year), the export price index of a country is 100 and its import price index is 100. In 2023, the export price index rises to 120 while the import price index rises to 150. Calculate the terms of trade in 2023 and comment.
\tcblower
\textbf{Solution.}
\[
\text{TOT}_{2023} = \frac{P_x}{P_m} \times 100 = \frac{120}{150} \times 100 = 80
\]
\textbf{Comment.} The terms of trade have deteriorated from 100 to 80, a fall of $20\%$. Although export prices rose (from 100 to 120), import prices rose even faster (from 100 to 150). One unit of exports now buys only $0.8$ of the imports it bought in 2020. The country must export $25\%$ more volume than before just to purchase the same basket of imports.
\end{exoresolu}

\courssub{Interpreting the Terms of Trade Index}

\begin{aretenir}[Reading the TOT index]
\begin{itemize}
\item $\text{TOT} > 100$: \textbf{favourable movement} (improvement). Export prices have risen faster than import prices (or fallen more slowly) since the base year. A unit of exports buys \emph{more} imports.
\item $\text{TOT} = 100$: no change relative to the base year.
\item $\text{TOT} < 100$: \textbf{unfavourable movement} (deterioration). Import prices have risen faster than export prices. A unit of exports buys \emph{fewer} imports.
\end{itemize}
Always compare with the base year (100), and always state \emph{which prices moved faster}.
\end{aretenir}

\begin{methode}[How to interpret a TOT index in an exam question]
\begin{enumerate}
\item \textbf{Compute} the index for each year using $\text{TOT} = (P_x/P_m)\times 100$, showing the substitution of values.
\item \textbf{Compare} with 100: above means improvement, below means deterioration relative to the base year.
\item \textbf{Explain the direction of movement}: did export prices rise faster, or did import prices rise faster? Quote the figures.
\item \textbf{State the consequence in real terms}: the quantity of imports obtainable per unit of exports has risen or fallen.
\item \textbf{Add a critical remark}: distinguish the \emph{movement} of the TOT from its \emph{impact} on the trade balance (see below), and mention the country's export structure (e.g.\ primary commodities for Cameroon).
\end{enumerate}
\end{methode}

The figure below shows a hypothetical terms of trade index over ten years, illustrating how the index fluctuates around the base value of 100.

\begin{popfigure}
\begin{center}
\begin{tikzpicture}
\begin{axis}[
    width=0.85\linewidth,
    height=6.5cm,
    xlabel={Year},
    ylabel={TOT index (base year 1 = 100)},
    xmin=1, xmax=10,
    ymin=70, ymax=120,
    xtick={1,2,3,4,5,6,7,8,9,10},
    ytick={70,80,90,100,110,120},
    grid=both,
    grid style={popline, dashed, very thin},
    axis line style={popink},
    tick label style={font=\small},
    label style={font=\small},
]
\addplot[color=popblue, very thick, mark=*, mark size=1.8pt] coordinates {
    (1,100) (2,106) (3,110) (4,104) (5,96) (6,88) (7,82) (8,90) (9,97) (10,102)
};
\addplot[color=poporange, dashed, thick, domain=1:10] {100};
\node[anchor=west, font=\small, text=poporange] at (axis cs:6.2,101.5) {Base $=100$};
\end{axis}
\end{tikzpicture}
\end{center}
\legende{Hypothetical terms of trade index over ten years. The index improves up to year 3, deteriorates sharply between years 4 and 7 (falling below 100), then recovers.}
\end{popfigure}

\courssub{Causes of Changes in the Terms of Trade}

The terms of trade move whenever export prices and import prices move at different speeds. The main causes are the following.

\textbf{1. Changes in world demand.}\par
If world demand for a country's exports rises (for example, a boom in Asian demand for cocoa or crude oil), export prices increase faster than import prices and the TOT improve. Conversely, a world recession reduces demand for primary commodities, pushing export prices down and deteriorating the TOT of exporters like Cameroon.

\textbf{2. Changes in world supply.}\par
A good harvest everywhere increases the supply of a commodity and lowers its world price. If many cocoa-producing countries record bumper harvests in the same season, the world cocoa price falls and Cameroon's TOT deteriorate. Supply shocks on the import side matter too: a poor harvest in grain-exporting countries raises the price of imported food.

\textbf{3. Inflation differentials.}\par
If domestic inflation is higher than in trading partners, domestic export prices tend to rise faster than import prices (measured in domestic currency), which may \emph{apparently} improve the TOT --- but this improvement is often artificial, since the country's goods become less competitive internationally. Conversely, high inflation in partner countries raises import prices and worsens the TOT.

\textbf{4. Exchange rate changes.}\par
A depreciation of the domestic currency makes imports more expensive in local currency; if import prices rise faster than export prices (in domestic currency), the TOT deteriorate. An appreciation has the opposite effect. For Cameroon, whose currency (the CFA franc) is pegged to the euro, movements of the euro against the US dollar matter greatly, since oil and many commodities are priced in dollars.

\textbf{5. Structural factors (Prebisch--Singer hypothesis).}\par
Over the long run, the prices of manufactured goods and services tend to rise relative to those of primary products, because demand for manufactures grows faster with income (higher income elasticity) and technological progress reduces the relative cost of manufactures. This explains the historical tendency for the TOT of primary-commodity exporters to deteriorate --- a concern associated with the \emph{Prebisch--Singer hypothesis}.

\courssub{A Favourable Movement Is Not Necessarily a Favourable Impact}

\begin{attention}[Improvement in the TOT $\neq$ improvement in the trade balance]
This is the classic examination trap. The TOT measure \emph{prices only}; the trade balance depends on \emph{prices and quantities}. An improvement in the TOT caused by a \textbf{rise in export prices} will improve the trade balance only if the quantity of exports demanded does not fall too much --- that is, if the \textbf{demand for exports is price-inelastic}. If demand for exports is elastic, higher export prices reduce export volume sharply, export revenue may fall, and the trade balance can worsen \emph{despite} a ``favourable'' TOT movement. Similarly, a deterioration of the TOT caused by \emph{cheaper} exports can boost export volumes so much that export earnings actually rise. Always ask: \emph{what caused the price change, and how do quantities respond?}
\end{attention}

\begin{exemplebox}[Elasticity at work]
Suppose the world price of cocoa rises by $10\%$, improving Cameroon's TOT. If world demand for cocoa is price-inelastic (chocolate producers need cocoa and have few substitutes), the quantity demanded falls by less than $10\%$, so Cameroon's export \emph{revenue} rises: favourable movement, favourable impact. But if the price of a more easily substituted export rises and buyers switch to competitors, export volumes collapse and revenue falls: favourable movement, unfavourable impact.
\end{exemplebox}

\courssub{Cameroon's Terms of Trade and Primary Commodity Dependence}

\begin{savaistu}[Cocoa, coffee and oil: why Cameroon's TOT are volatile]
Cameroon's export earnings depend heavily on a small number of primary commodities: crude oil, cocoa, coffee, timber, cotton and bananas. Their prices are fixed on world markets (London, New York) and can swing sharply with weather, speculation and world demand, while Cameroon's imports --- machinery, vehicles, medicines, refined fuel --- have relatively stable or steadily rising prices. As a result, Cameroon's terms of trade tend to improve during commodity booms (for example when oil prices surge) and deteriorate when commodity prices collapse. This volatility complicates budget planning, since State revenue itself depends on export earnings. It is a major argument for \textbf{diversifying} the economy and processing raw materials locally (e.g.\ transforming cocoa into chocolate) so as to export higher-value, more price-stable products.
\end{savaistu}

This dependence illustrates a general lesson: a country exporting few primary products and importing diversified manufactures has \emph{no control} over either its export or its import prices. It is a \cle{price-taker} on both sides, and its terms of trade are at the mercy of world markets.

\courssub{Consolidation: Worked Exercise}

\begin{exercice}[Terms of trade of a hypothetical economy]
The table below gives the export and import price indices of an economy (base year 2019 $= 100$).

\begin{center}
\begin{tabularx}{\linewidth}{|l|X|X|}
\hline
\rowcolor{popblueL} \textbf{Year} & \textbf{Export price index} & \textbf{Import price index} \\
\hline
2020 & 104 & 100 \\
\hline
2021 & 110 & 108 \\
\hline
2022 & 108 & 120 \\
\hline
\end{tabularx}
\end{center}

\begin{enumerate}
\item Calculate the terms of trade for each year.
\item Comment on the evolution of the country's exchange position.
\item Give two possible causes of the movement observed in 2022.
\item Does the movement in 2021 necessarily improve the trade balance? Justify.
\end{enumerate}
\end{exercice}

\begin{corrige}[Exercise: Terms of trade]
\textbf{1. Calculations.}
\begin{align*}
\text{TOT}_{2020} &= \frac{104}{100}\times 100 = 104 \\
\text{TOT}_{2021} &= \frac{110}{108}\times 100 \approx 101.9 \\
\text{TOT}_{2022} &= \frac{108}{120}\times 100 = 90
\end{align*}
\textbf{2. Comment.} The TOT improved slightly in 2020 and 2021 (above 100: export prices rose faster than import prices), then deteriorated sharply in 2022 to 90: import prices rose much faster than export prices, so a unit of exports buys $10\%$ fewer imports than in the base year.\par
\textbf{3. Possible causes of the 2022 deterioration:} a fall in world demand for the country's exports (e.g.\ recession in partner countries) lowering export prices; a rise in import prices due to inflation abroad or a depreciation of the domestic currency; or a world supply shock raising the price of imported food or energy.\par
\textbf{4. Not necessarily.} The 2021 improvement reflects prices only. If it resulted from higher export prices and demand for those exports is price-elastic, export volumes and revenue may have fallen, worsening the trade balance. Only if export demand is inelastic (or if the improvement came from cheaper imports with stable volumes) would the trade balance clearly improve.
\end{corrige}

\begin{aretenir}[Key points on the Terms of Trade]
\begin{itemize}
\item $\text{TOT} = (P_x/P_m)\times 100$, using price indices with a common base year.
\item TOT above 100: improvement; below 100: deterioration relative to the base year.
\item Movements are caused by world demand and supply shifts, inflation differentials and exchange rate changes.
\item A favourable \emph{movement} does not guarantee a favourable \emph{impact} on the trade balance --- quantities and elasticities must be considered.
\item Primary-commodity exporters like Cameroon experience volatile TOT because commodity prices fluctuate sharply while import prices are more stable.
\end{itemize}
\end{aretenir}

Having understood how the \emph{prices} of exports and imports compare, we now turn to a broader accounting question: how do we record \emph{all} of a country's economic transactions with the rest of the world? This is the purpose of the \emph{balance of payments}, the subject of the next section.

\coursec{The Balance of Payments: Structure and Recording}

In the previous section we studied the \cle{terms of trade}, which measure the \emph{price} relationship between what a country sells abroad and what it buys from abroad. But knowing whether export prices are rising faster than import prices is not enough: a government, a firm or an examiner also wants to know the \emph{quantities and values} of all the money flowing into and out of the country. Does Cameroon earn more foreign currency from cocoa, oil and timber than it spends on imported machinery, fuel and medicines? Are foreigners investing in the country, or is capital leaving? The document that answers all these questions in a systematic way is the \cle{balance of payments}. It is one of the most examined topics in the GCE Advanced Level paper, both in structured questions and in data-response exercises, so we shall build it up carefully: first the idea, then the structure, then the recording rules, and finally a full worked account.

\courssub{Definition and Purpose of the Balance of Payments}

\textbf{Intuition first.} Imagine a household that keeps a notebook of every franc it receives from outside (salary, gifts, sales of farm produce) and every franc it pays to outsiders (food, school fees, loan repayments). At the end of the year the notebook shows whether the household is a net receiver or a net payer. A country keeps exactly such a notebook, but at the level of \emph{all its residents} dealing with \emph{all non-residents}. That notebook is the balance of payments.

\begin{definition}[Balance of Payments]
The \cle{balance of payments} (BOP) is a \textbf{systematic statistical record of all economic transactions between the residents of a country and the rest of the world during a given period of time}, usually one year (sometimes a quarter). Transactions are recorded in money values using the principle of double-entry bookkeeping, so that every transaction appears twice: once as a \cle{credit} and once as a \cle{debit}.
\end{definition}

Three points in this definition deserve emphasis:

\begin{itemize}
\item \textbf{Residents, not citizens.} What matters is residence (living and working in the country for at least about a year), not nationality. A French engineer resident in Douala counts as a Cameroonian resident for BOP purposes.
\item \textbf{All economic transactions.} This covers trade in goods, trade in services, income flows (wages, profits, interest, dividends), transfers (gifts, grants, remittances) and financial flows (investment, loans, changes in reserves).
\item \textbf{A flow, over a period.} The BOP records flows during a year; it is not a stock like the national debt.
\end{itemize}

The balance of payments serves several purposes: it informs the government and the central bank about the country's external position, it signals pressure on the foreign exchange reserves and the exchange rate, and it guides policy on trade, borrowing and investment.

\courssub{The Structure of the Balance of Payments}

The standard presentation (consistent with the IMF methodology used by BEAC for Cameroon and the CEMAC zone) divides the account into three main parts: the \cle{current account}, the \cle{capital and financial account}, and the \cle{balancing item} (net errors and omissions).

\begin{center}
\begin{tabularx}{\linewidth}{|l|X|}
\hline
\rowcolor{popblueL} \textbf{Component} & \textbf{Contents} \\
\hline
\textbf{A. Current account} & \\
\hline
\quad 1. Visible trade (goods) & Exports (credit) and imports (debit) of physical goods: cocoa, crude oil, timber, machinery, food \\
\hline
\quad 2. Invisible trade (services) & Transport, insurance, tourism, banking, telecommunications, government services \\
\hline
\quad 3. Income & Compensation of employees; investment income (profits, dividends, interest) flowing in and out \\
\hline
\quad 4. Current transfers & Workers' remittances, gifts, grants, pensions, aid transfers with no quid pro quo \\
\hline
\textbf{B. Capital account} & Capital transfers (debt forgiveness, migrants' transfers), transactions in non-produced non-financial assets \\
\hline
\textbf{C. Financial account} & \\
\hline
\quad 1. Direct investment & Long-term investment giving control or lasting influence (e.g.\ a foreign firm building a factory) \\
\hline
\quad 2. Portfolio investment & Purchase of shares and bonds without control \\
\hline
\quad 3. Other investment & Trade credits, loans, currency and deposits \\
\hline
\quad 4. Reserve assets & Changes in official foreign exchange reserves held by the central bank \\
\hline
\textbf{D. Net errors and omissions} & Balancing item making the accounts add up to zero \\
\hline
\end{tabularx}
\end{center}

\begin{popfigure}
\begin{tikzpicture}[
box/.style={draw=popblue, thick, rounded corners=2pt, fill=popblueL, text width=4.1cm, align=center, minimum height=1.1cm, font=\small},
box2/.style={draw=popgreen, thick, rounded corners=2pt, fill=popgreenL, text width=4.1cm, align=center, minimum height=1.1cm, font=\small},
box3/.style={draw=poporange, thick, rounded corners=2pt, fill=white, text width=4.1cm, align=center, minimum height=1.1cm, font=\small},
arr/.style={-{Stealth[length=2.5mm]}, thick, popmuted}]
\node[box, align=center] (ca) at (0,2.4){\textbf{Current account}\\ goods, services, income, transfers};
\node[box2, align=center] (cfa) at (0,0.6){\textbf{Capital \& financial account}\\ investment, loans, reserves};
\node[box3, align=center] (neo) at (0,-1.2){\textbf{Net errors \& omissions}\\ (balancing item)};
\node[draw=popdark, thick, rounded corners=2pt, text width=4.6cm, align=center, font=\small\bfseries] (bop) at (7,0.6) {BALANCE OF PAYMENTS\\ always balances};
\draw[arr] (ca.east) -- (bop.north west);
\draw[arr] (cfa.east) -- (bop.west);
\draw[arr] (neo.east) -- (bop.south west);
\end{tikzpicture}
\legende{The three building blocks of the balance of payments: the current account, the capital and financial account, and the balancing item.}
\end{popfigure}

\begin{definition}[Balance of Trade (Visible Balance)]
The \cle{balance of trade} (or \cle{visible balance}) is the difference between the value of \textbf{exports of goods} and the value of \textbf{imports of goods} over a period:
\[ \text{Balance of trade} = \text{Value of visible exports} - \text{Value of visible imports}. \]
If exports of goods exceed imports of goods there is a \emph{trade surplus}; if imports exceed exports there is a \emph{trade deficit}.
\end{definition}

Notice carefully the hierarchy of balances, which examiners love to test:

\begin{itemize}
\item The \textbf{balance of trade} covers \emph{goods only}.
\item The \textbf{current account balance} covers goods \emph{plus} services, income and current transfers. It is therefore wider than the balance of trade.
\item The \textbf{overall balance of payments} position is usually judged by the balance on \emph{autonomous} transactions (current plus capital and non-reserve financial flows), the counterpart of which is the change in \emph{reserve assets} (and exceptional financing).
\end{itemize}

\courssub{Double-Entry Recording: Credits and Debits}

Every international transaction has two sides. When Cameroon exports cocoa worth 100 million FCFA to the Netherlands, two things happen: goods leave Cameroon (an export, recorded as a \cle{credit} in the current account) and a payment or claim on a foreign bank enters (recorded as a \cle{debit} in the financial account, since Cameroon acquires a financial asset abroad). Because each transaction is entered twice, once as a credit and once as a debit, the total of all credits must equal the total of all debits: \textbf{in accounting terms the balance of payments always balances}.

\begin{methode}[Classifying a Transaction as Credit or Debit]
Apply these rules step by step:
\begin{enumerate}
\item Ask: does the transaction bring foreign currency \emph{into} the country or create an obligation for foreigners to pay us? If yes, it is a \textbf{credit} ($+$).
\item Ask: does the transaction send currency \emph{out of} the country or create an obligation for us to pay foreigners? If yes, it is a \textbf{debit} ($-$).
\item Typical \textbf{credits}: exports of goods and services; income received from abroad; transfers received (remittances, grants); foreign investment flowing into the country; borrowing from abroad; a \emph{decrease} in official reserves.
\item Typical \textbf{debits}: imports of goods and services; income paid abroad (profits repatriated by foreign firms); transfers sent abroad; domestic investment abroad; lending to foreigners; an \emph{increase} in official reserves.
\item Then place the item in the correct sub-account (goods, services, income, transfers, direct investment, portfolio investment, other investment, reserves).
\end{enumerate}
\end{methode}

\begin{attention}[``The BOP Always Balances'' — So What Is a Deficit?]
A very common mistake is to write that ``a country has a balance of payments deficit, so its accounts do not balance.'' This is contradictory. Because of double-entry recording, the BOP \emph{always balances in accounting terms}. When economists speak of a \textbf{BOP deficit} or \textbf{surplus}, they refer to the balance on \textbf{autonomous transactions} (transactions undertaken for their own sake: trade, income, investment). A deficit on autonomous transactions must be financed by \textbf{accommodating transactions}: drawing down reserves or borrowing officially. Never confuse ``the accounts do not balance'' with ``autonomous transactions are in deficit.''
\end{attention}

\courssub{Practical Work: Constructing and Balancing a Simplified BOP Account}

We now do exactly what Practical Work 25 requires: build a balance of payments account from raw data and find the balancing item.

\begin{exoresolu}[Constructing a Balance of Payments Account]
The following data (in billions of FCFA) refer to the hypothetical economy of ``Littoralland'' for the year 2024:
exports of goods $= 850$; imports of goods $= 1\,050$; exports of services $= 180$; imports of services $= 260$; net income received from abroad $= -40$ (i.e.\ paid out); net current transfers received $= +120$; capital account balance $= +15$; net direct investment inflow $= +90$; net portfolio investment inflow $= +25$; net other investment $= +60$; net errors and omissions $= +10$. Find the change in reserve assets and present the full account.
\tcblower
\textbf{Step 1: the current account.}
\begin{align*}
\text{Visible balance} &= 850 - 1\,050 = -200 \\
\text{Invisible (services) balance} &= 180 - 260 = -80 \\
\text{Current account} &= -200 + (-80) + (-40) + 120 = -200
\end{align*}
\textbf{Step 2: capital and financial account (excluding reserves).}
\[ \text{Capital account} + \text{DI} + \text{PI} + \text{other} = 15 + 90 + 25 + 60 = +190 \]
\textbf{Step 3: impose the accounting identity.} The sum of all items, including the change in reserves $\Delta R$ and net errors and omissions, must be zero:
\[ -200 + 190 + 10 + \Delta R = 0 \quad\Longrightarrow\quad \Delta R = 0. \]
Here the autonomous deficit of $-200 + 190 + 10 = 0$ happens to be fully covered, so reserves are unchanged. The presented account is:
\begin{center}
\begin{tabularx}{\linewidth}{|X|r|r|}
\hline
\rowcolor{popblueL} \textbf{Item} & \textbf{Credit} & \textbf{Debit / net} \\
\hline
Visible trade (goods) & 850 & $-1\,050$ \\
\hline
Invisible trade (services) & 180 & $-260$ \\
\hline
Income (net) & \multicolumn{2}{r|}{$-40$} \\
\hline
Current transfers (net) & \multicolumn{2}{r|}{$+120$} \\
\hline
\textbf{Current account balance} & \multicolumn{2}{r|}{$\mathbf{-200}$} \\
\hline
Capital account & \multicolumn{2}{r|}{$+15$} \\
\hline
Direct investment (net) & \multicolumn{2}{r|}{$+90$} \\
\hline
Portfolio investment (net) & \multicolumn{2}{r|}{$+25$} \\
\hline
Other investment (net) & \multicolumn{2}{r|}{$+60$} \\
\hline
Net errors and omissions & \multicolumn{2}{r|}{$+10$} \\
\hline
\textbf{Change in reserve assets} & \multicolumn{2}{r|}{$\mathbf{0}$} \\
\hline
\textbf{Total} & \multicolumn{2}{r|}{$\mathbf{0}$} \\
\hline
\end{tabularx}
\end{center}
The account balances exactly, as double-entry bookkeeping requires. The ``deficit'' story here is the current account deficit of 200 billion FCFA, financed by net capital inflows of 190 billion plus 10 billion of unrecorded net inflows.
\end{exoresolu}

\begin{attention}[Sign of Reserve Changes]
An \emph{increase} in official reserves is a \textbf{debit} (the central bank is acquiring foreign assets, i.e.\ money is ``going out'' into reserve holdings); a \emph{decrease} in reserves is a \textbf{credit}. Students routinely reverse this. Think of reserves as the country ``buying'' foreign currency assets: buying is always a debit.
\end{attention}

\courssub{Balance of Trade, Current Account and Overall Balance}

It is essential to keep the three concepts distinct:

\begin{itemize}
\item \textbf{Balance of trade} $=$ visible exports $-$ visible imports. In the worked example: $-200$ billion FCFA (a trade deficit).
\item \textbf{Current account balance} $=$ balance of trade $+$ services $+$ income $+$ current transfers. In the example: $-200$ billion FCFA (the services and income deficits exactly offset the transfers surplus).
\item \textbf{Overall balance} (balance on autonomous transactions) $=$ current account $+$ capital account $+$ non-reserve financial account $+$ errors and omissions. Its counterpart is the change in reserves. In the example the overall balance is zero; had autonomous transactions summed to $-50$, reserves would have fallen by $50$ (a credit of $+50$ in the reserve line).
\end{itemize}

\begin{exemplebox}[Cameroon's Typical Pattern]
Without quoting precise figures, the structure of Cameroon's BOP illustrates these ideas well: exports are dominated by a few commodities (crude oil, cocoa, timber, cotton, aluminium), while imports of manufactured goods, equipment and food are large, so the visible balance is often in deficit. The services and income balances also tend to be negative (freight and insurance on imports, profits and interest paid abroad), while current transfers — notably remittances from Cameroonians living abroad — provide a partial offset. The financial account records foreign direct investment (for example in oil, mining and infrastructure) and borrowing, and any remaining gap is reflected in the foreign exchange reserves pooled within the BEAC for the CEMAC zone.
\end{exemplebox}

\begin{aretenir}[Key Points]
\begin{itemize}
\item The BOP is a systematic, double-entry record of all transactions between residents and the rest of the world over a period.
\item Main parts: \textbf{current account} (goods, services, income, current transfers), \textbf{capital account}, \textbf{financial account} (direct, portfolio, other investment, reserve assets), plus \textbf{net errors and omissions}.
\item Credits bring currency in; debits send currency out. Reserve increases are debits.
\item The BOP always balances in accounting terms; ``deficit/surplus'' refers to autonomous transactions only.
\item Balance of trade $\subset$ current account $\subset$ overall balance: know which is which.
\end{itemize}
\end{aretenir}

\begin{exercice}[Building a BOP Account]
From the following data for ``Sahelland'' (billions of FCFA): exports of goods $= 600$; imports of goods $= 500$; services (net) $= -70$; income (net) $= -30$; current transfers (net) $= +50$; capital account $= +5$; direct investment (net) $= +40$; portfolio investment (net) $= -10$; other investment (net) $= -25$; net errors and omissions $= -5$. Compute: (a) the balance of trade; (b) the current account balance; (c) the change in reserve assets; (d) state whether the country faces an autonomous deficit or surplus.
\end{exercice}

\begin{corrige}[Exercise 1]
(a) Balance of trade $= 600 - 500 = +100$ (surplus). (b) Current account $= 100 - 70 - 30 + 50 = +50$. (c) Sum of autonomous items: $50 + 5 + 40 - 10 - 25 - 5 = +55$; hence $\Delta R = -55$ in the presentation above, i.e.\ reserves \emph{increase} by 55 (recorded as a debit of $-55$). (d) Autonomous transactions are in surplus of 55, matched by an accumulation of reserves.
\end{corrige}

\begin{savaistu}[Exam Tip: The Four Components of the Current Account]
Examiners frequently ask candidates to ``state the components of the current account.'' Memorise the quartet: \textbf{goods} (visible trade), \textbf{services} (invisible trade), \textbf{income} (compensation of employees and investment income), and \textbf{current transfers} (remittances, grants, gifts). A mnemonic: \emph{``Goods, Services, Income, Transfers''} — \textbf{GSIT}. Do not include investment \emph{flows} (direct or portfolio): those belong to the financial account, and confusing them with investment \emph{income} (a current account item) is a classic loss of marks.
\end{savaistu}

\coursec{Balance of Payments Disequilibrium: Deficit, Surplus and Correction}

In the previous section we learnt how the balance of payments (BOP) is structured and recorded: the current account, the capital account, the financial account and the balancing items. We also saw that, because of double-entry bookkeeping, the BOP always \emph{balances} in an accounting sense. Yet economists and journalists constantly speak of a country ``suffering a balance of payments deficit''. What do they mean? They mean that the \cle{autonomous transactions} --- the exports, imports, investment and capital flows undertaken for their own sake, independently of the state of the balance of payments --- do not cancel out, so that the gap must be filled by \cle{accommodating transactions}: drawing down foreign exchange reserves or borrowing from abroad. It is this underlying imbalance, called a \cle{BOP disequilibrium}, that we now study: its two forms (deficit and surplus), their causes and consequences, and the policies used to correct them.

\courssub{The Balance of Payments Deficit}

\begin{definition}[Balance of payments deficit and surplus]
A \cle{balance of payments deficit} exists when a country's total autonomous payments to foreigners (for imports of goods and services, income payments and capital outflows) exceed its total autonomous receipts from foreigners, so that the difference must be financed by running down official reserves or by borrowing abroad.

A \cle{balance of payments surplus} exists when autonomous receipts from foreigners exceed autonomous payments to foreigners, leading to an accumulation of foreign exchange reserves or net lending abroad.
\end{definition}

Intuitively, a deficit country is like a household that spends more on shopping than it earns from its salary: it must either dip into its savings (reserves) or borrow. Neither can go on forever. A surplus country, by contrast, is like a household that earns more than it spends and piles up savings --- comfortable, but not necessarily optimal if the savings lie idle.

\textbf{Causes of a balance of payments deficit.}\par

\begin{itemize}
\item \textbf{Excessive imports.} If domestic incomes rise rapidly, or if consumers develop a strong taste for foreign goods (cars, electronics, refined foods), import expenditure grows faster than export earnings. Rapid population growth and ambitious public investment programmes (which require imported machinery) have the same effect.
\item \textbf{Falling export prices.} A country like Cameroon, whose exports are concentrated in primary commodities (cocoa, coffee, timber, crude oil, cotton), suffers when world prices of these commodities fall. Even if the \emph{volume} exported is unchanged, export \emph{revenue} collapses --- a deterioration in the terms of trade.
\item \textbf{Domestic inflation.} If prices rise faster at home than abroad, domestic goods become uncompetitive: exports fall (foreigners buy cheaper goods elsewhere) and imports rise (residents switch to relatively cheaper foreign goods).
\item \textbf{An overvalued currency.} If the exchange rate is fixed at a level that makes foreign currency artificially cheap, imports are encouraged and exports discouraged, producing a persistent current account deficit.
\item \textbf{Debt servicing.} Heavy external debt obliges the country to make large interest and amortisation payments abroad every year --- a drain on the current account (investment income debits) that can itself create a deficit.
\item \textbf{Structural and political factors.} Civil unrest, poor infrastructure, low productivity or a bad harvest can reduce export supply, while capital flight (residents moving money abroad) worsens the financial account.
\end{itemize}

\textbf{Consequences of a persistent deficit.}\par

A temporary deficit is not dangerous --- it may even finance growth-promoting imports of machinery. The problem is a \emph{persistent} (fundamental) deficit:

\begin{itemize}
\item \textbf{Depletion of foreign exchange reserves.} The central bank must sell its stock of foreign currency to pay for the excess imports. Reserves are finite; once they fall below a safe level (often measured in months of import cover), the country faces a crisis.
\item \textbf{Increased external borrowing and indebtedness.} The gap may be financed by foreign loans, but these add to future debt-servicing obligations, worsening the deficit later --- a vicious circle.
\item \textbf{Pressure on the currency.} Persistent excess demand for foreign exchange creates downward pressure on the currency: under a floating rate the currency depreciates; under a fixed rate, speculation may force a devaluation.
\item \textbf{Loss of confidence and policy constraints.} Foreign investors and creditors become reluctant to lend, and the government may be forced into painful austerity measures imposed by creditors such as the IMF.
\end{itemize}

\courssub{Measures to Correct a Balance of Payments Deficit}

Corrective policies fall into two great families, plus two supplementary tools.

\begin{methode}[Distinguishing expenditure-reducing from expenditure-switching policies]
\begin{enumerate}
\item Ask: does the policy work by making the whole economy \emph{spend less} (on everything, imports included)? If yes, it is \cle{expenditure-reducing} (e.g.\ deflationary fiscal or monetary policy).
\item Ask: does the policy work by changing \emph{relative prices} so that spending is \emph{switched} between foreign and home goods? If yes, it is \cle{expenditure-switching} (e.g.\ devaluation, tariffs, quotas, export subsidies).
\item Remember: expenditure-reducing policies treat the symptom through lower income; expenditure-switching policies attack the cause through competitiveness. In examination essays, always present at least one policy of each type and evaluate both.
\end{enumerate}
\end{methode}

\textbf{1. Expenditure-reducing policies (deflation).}\par
The government reduces aggregate demand through \emph{contractionary fiscal policy} (higher taxes, lower public spending) and \emph{contractionary monetary policy} (higher interest rates, credit restrictions). Lower national income reduces the demand for imports, since imports vary positively with income; lower domestic inflation also improves competitiveness. \emph{Limitation:} deflation causes unemployment and slows growth --- a high price to pay, especially for a developing country.

\textbf{2. Expenditure-switching policies.}\par
\begin{itemize}
\item \textbf{Devaluation / depreciation} of the currency makes exports cheaper in foreign currency and imports dearer in domestic currency, switching demand towards home goods.
\item \textbf{Tariffs and quotas} raise the domestic price of imports or limit their quantity directly.
\item \textbf{Export promotion}: subsidies, tax holidays, export credit and marketing support for exporters (e.g.\ support to cocoa and timber exporters).
\end{itemize}

\textbf{3. Exchange controls.}\par
The state rations foreign exchange: importers and travellers must obtain official authorisation to buy foreign currency. This directly limits payments abroad but encourages black markets and corruption, and discourages foreign investment.

\textbf{4. Borrowing from the IMF.}\par
A member country in deficit can draw on the International Monetary Fund to finance the gap while adjustment takes place. IMF support usually comes with \emph{conditionality} (structural adjustment conditions: devaluation, budget cuts, liberalisation), which may be socially painful.

\begin{attention}[Devaluation works only if the Marshall--Lerner condition holds]
A very common examination mistake is to write ``devaluation automatically cures a deficit''. It does not. Devaluation improves the trade balance \emph{only if} the \cle{Marshall--Lerner condition} is satisfied: the sum of the price elasticities of demand for exports and imports must be greater than one (in absolute value):
\[
|PED_X| + |PED_M| > 1.
\]
If demand is inelastic (as for essential imports such as fuel, medicine or spare parts), the import bill \emph{rises} after devaluation because the same quantities now cost more in domestic currency. Always state this condition when evaluating devaluation.
\end{attention}

\textbf{The J-curve.}\par
Even when the Marshall--Lerner condition holds, the improvement is not immediate. In the short run, import and export \emph{volumes} are slow to adjust (contracts are fixed, consumers keep their habits), while import \emph{prices} in domestic currency rise at once. The trade balance therefore \emph{worsens first} before improving --- tracing a path shaped like the letter J.

\begin{popfigure}
\begin{tikzpicture}
\begin{axis}[
width=11cm, height=7cm,
axis lines=middle,
xlabel={Time after devaluation}, ylabel={Trade balance},
xlabel style={anchor=west}, ylabel style={anchor=south},
xtick=\empty, ytick={0},
yticklabels={$0$},
xmin=0, xmax=10, ymin=-3.2, ymax=2.6,
domain=0:10, samples=100,
]
\addplot[very thick, popblue] {2.2*(1-exp(-0.55*x)) - 2.6*exp(-0.9*x)};
\draw[dashed, popmuted] (axis cs:0,0) -- (axis cs:10,0);
\node[popblue, anchor=west] at (axis cs:6.6,1.9) {\small Improvement};
\node[poporange, anchor=north] at (axis cs:1.6,-2.5) {\small Initial worsening};
\node[popmuted, anchor=south west] at (axis cs:0.1,0.1) {\small Devaluation};
\end{axis}
\end{tikzpicture}
\legende{The J-curve: after a devaluation, the trade balance typically deteriorates before it improves, because quantities adjust more slowly than prices.}
\end{popfigure}

\begin{popfigure}
\begin{tikzpicture}[
box/.style={draw=popblue, thick, rounded corners, fill=popblueL, text width=3.6cm, align=center, minimum height=1cm, font=\small},
head/.style={draw=popdark, thick, rounded corners, fill=popdark, text=white, text width=4.2cm, align=center, minimum height=1cm, font=\small\bfseries},
arr/.style={->, thick, popdark}
]
\node[head] (start) at (0,0) {BOP deficit diagnosed};
\node[box] (red) at (-5.4,-2.2) {Expenditure-reducing: deflationary fiscal \& monetary policy};
\node[box] (switch) at (0,-2.2) {Expenditure-switching: devaluation, tariffs, quotas, export promotion};
\node[box] (control) at (5.4,-2.2) {Direct controls: exchange control, import licensing};
\node[box] (imf) at (0,-4.6) {Financing the gap: IMF borrowing, reserves};
\node[box, fill=popgreenL, draw=popgreen] (goal) at (0,-6.8) {Restored external equilibrium};
\draw[arr] (start) -- (red);
\draw[arr] (start) -- (switch);
\draw[arr] (start) -- (control);
\draw[arr] (red) |- (goal.west);
\draw[arr] (switch) -- (imf);
\draw[arr] (control) |- (goal.east);
\draw[arr] (imf) -- (goal);
\end{tikzpicture}
\legende{Policy options for correcting a balance of payments deficit.}
\end{popfigure}

\courssub{The Balance of Payments Surplus}

A surplus sounds like good news --- and usually it is. But a \emph{persistent, large} surplus can itself be a problem:

\begin{itemize}
\item \textbf{Inflationary pressure.} The inflow of foreign exchange, when converted into domestic currency, increases the money supply and aggregate demand, fuelling demand-pull inflation.
\item \textbf{Retaliation by trading partners.} One country's surplus is another's deficit. Deficit partners may impose tariffs and quotas in response, provoking trade conflicts.
\item \textbf{Idle reserves.} Excessive reserves held in low-yield foreign assets represent resources that could have been invested at home in roads, schools and hospitals --- an opportunity cost.
\item \textbf{Upward pressure on the currency}, which may eventually harm exporters.
\end{itemize}

\textbf{Measures to correct a surplus.}\par
\begin{itemize}
\item \textbf{Revaluation} (or appreciation) of the currency: exports become dearer, imports cheaper, reducing the surplus.
\item \textbf{Import liberalisation}: removing tariffs and quotas encourages imports.
\item \textbf{Expansionary domestic policy}: lower interest rates and higher public spending raise incomes and therefore imports, while the resulting higher prices reduce export competitiveness.
\item \textbf{Encouraging capital outflows and foreign aid}: lending or granting the surplus to poorer partners.
\end{itemize}

\begin{center}
\begin{tabularx}{\linewidth}{|l|X|X|}
\hline
\rowcolor{popblueL} & \textbf{Deficit} & \textbf{Surplus} \\
\hline
\textbf{Meaning} & Autonomous payments exceed receipts; financed by reserves or borrowing & Autonomous receipts exceed payments; reserves accumulate \\
\hline
\textbf{Main dangers} & Reserve depletion, indebtedness, currency pressure, loss of confidence & Inflation, retaliation by partners, idle reserves, overvalued currency \\
\hline
\textbf{Expenditure policies} & Deflation (reduce spending); devaluation, tariffs, export promotion (switch spending) & Reflation (raise spending); revaluation, import liberalisation \\
\hline
\textbf{Other tools} & Exchange control, IMF borrowing & Capital outflows, foreign aid \\
\hline
\end{tabularx}
\end{center}

\courssub{Further Studies: Improving Cameroon's Balance of Payments Position}

\begin{aretenir}[Cameroon and the BEAC monetary zone]
Cameroon belongs to the \cle{CEMAC} monetary zone, whose common currency, the CFA franc, is issued by the \cle{BEAC} and pegged to the euro. Consequently, Cameroon \textbf{cannot devalue its currency independently}: devaluation is a collective decision of the zone (as in January 1994, when the CFA franc was devalued by 50\%). Cameroon must therefore rely mainly on fiscal policy, structural reforms and supply-side measures rather than on the exchange rate instrument.
\end{aretenir}

Given this constraint, the strategies most often recommended for Cameroon are:

\begin{itemize}
\item \textbf{Export diversification.} Reducing dependence on a few primary commodities (oil, cocoa, timber, cotton) by developing new export products (processed foods, manufactured goods, services) so that a fall in one world price no longer wrecks the current account.
\item \textbf{Value addition.} Processing cocoa into chocolate and timber into furniture \emph{before} export raises export unit values, creates jobs and captures more of the final price, instead of exporting raw beans and logs cheaply.
\item \textbf{Import substitution.} Producing locally goods currently imported --- rice, maize, fish, pharmaceuticals, cement inputs --- reduces the import bill. This requires cheap energy, credit and infrastructure.
\item \textbf{Tourism promotion.} Earnings from tourism (``Africa in miniature'') count as invisible exports and can diversify foreign exchange receipts.
\item \textbf{Regional trade within CEMAC and CEEAC.} Expanding exports to neighbouring markets (Chad, Central African Republic, Gabon, Nigeria) and exploiting the African Continental Free Trade Area widens the market for Cameroonian goods.
\item \textbf{Prudent debt management and the fight against capital flight}, so that debt servicing and illicit outflows stop draining the balance of payments.
\end{itemize}

\begin{savaistu}[Why can't Cameroon simply print its way out of a deficit?]
Because the CFA franc is issued by a regional central bank (the BEAC) serving six CEMAC states, no single member can finance its external deficit by creating money at will. Statutory rules limit how much credit the BEAC may extend to each government. This discipline protects the currency's stability, but it also means adjustment must come from real policies --- competitiveness, diversification and fiscal prudence --- rather than from monetary shortcuts.
\end{savaistu}

\courssub{Worked Exercise and Evaluation}

\begin{exoresolu}[Analysing a deficit and choosing policies]
The Republic of Ndongué exports cocoa and imports machinery and refined fuel. World cocoa prices fall by 30\% while domestic inflation reaches 8\%, double the rate of its trading partners. (a) Explain two reasons why Ndongué's balance of payments moves into deficit. (b) Distinguish between an expenditure-reducing and an expenditure-switching policy the government could use. (c) State one condition necessary for a devaluation to succeed.
\tcblower
\textbf{(a)} First, the fall in world cocoa prices reduces export revenue even if the volume exported is unchanged --- a deterioration of the terms of trade that widens the trade gap. Second, domestic inflation above that of trading partners makes Ndongué's goods uncompetitive: exports fall and residents switch to relatively cheaper imports, worsening the current account on both sides.

\textbf{(b)} An \emph{expenditure-reducing} policy lowers total spending in the economy so that imports fall with income: for example, cutting government expenditure and raising taxes (deflationary fiscal policy) or raising interest rates. An \emph{expenditure-switching} policy changes relative prices to divert demand from foreign to home goods: for example, devaluing the currency or imposing tariffs on imported machinery and fuel.

\textbf{(c)} Devaluation improves the trade balance only if the Marshall--Lerner condition holds, i.e.\ the sum of the price elasticities of demand for exports and imports exceeds unity ($|PED_X| + |PED_M| > 1$). Even then, the J-curve warns that the balance may worsen before it improves.
\end{exoresolu}

\begin{exercice}[Essay practice]
``A balance of payments deficit is always a sign of economic failure, and devaluation is always the best cure.'' Discuss this statement with reference to the causes of deficits, the Marshall--Lerner condition, the J-curve, and the constraints faced by a CEMAC member country such as Cameroon.
\end{exercice}

\begin{corrige}[Exercise 1]
\textbf{Plan and key points.} \emph{Introduction:} define a BOP deficit as an excess of autonomous payments over receipts financed by reserves or borrowing. \emph{Against the statement:} a deficit may finance growth-enhancing capital imports and may be temporary; it is only a \emph{persistent, fundamental} deficit that signals trouble (reserve depletion, indebtedness, currency pressure). \emph{On devaluation:} it is an expenditure-switching policy that works only if the Marshall--Lerner condition holds; the J-curve shows an initial worsening; it is inflationary and raises the domestic-currency burden of foreign debt; and Cameroon, inside the BEAC/CFA franc zone, cannot devalue unilaterally. \emph{Alternatives:} expenditure-reducing policies (costly in unemployment), tariffs and quotas, export promotion, and structural measures (diversification, value addition, import substitution, tourism, regional trade). \emph{Conclusion:} a deficit is a warning signal, not proof of failure, and the appropriate cure depends on the cause and on the country's institutional constraints --- there is no single ``best'' remedy.
\end{corrige}

\coursec{International Economic Institutions and Regional Groupings}

In the previous section we saw that a country facing a persistent balance of payments deficit must find ways to correct it: expenditure-reducing policies, expenditure-switching policies such as devaluation, and direct controls. But where does a country like Cameroon turn when its foreign reserves are exhausted and it cannot finance its deficit alone? It turns to the \cle{international economic institutions} -- above all the International Monetary Fund and the World Bank. And when it seeks to expand its markets and bargain collectively, it joins \cle{regional economic groupings} such as CEMAC and the AfCFTA. This section studies these two pillars of international economic cooperation.

\courssub{The International Monetary Fund (IMF)}

\textbf{Origin.} The IMF was created at the Bretton Woods Conference (USA) in 1944 and began operations in 1945. Its headquarters are in Washington, D.C. It was born from the chaos of the 1930s, when competitive devaluations and trade wars deepened the Great Depression. The founders wanted an institution that would guarantee stability in international payments.

\begin{definition}[The International Monetary Fund]
The \cle{IMF} is an international organisation whose purpose is to promote international monetary cooperation, ensure the stability of exchange rates, facilitate the balanced growth of international trade, and provide temporary financial assistance to member countries experiencing balance of payments difficulties.
\end{definition}

\textbf{Objectives and functions.} The IMF's work can be grouped into three main functions:

\begin{enumerate}
\item \cle{Surveillance}: the IMF monitors the economic and financial policies of its member countries (annual ``Article IV consultations'') and gives policy advice to prevent crises.
\item \cle{Lending}: it provides short- and medium-term loans to members facing balance of payments deficits, financed from a pool of \cle{quotas} -- subscriptions paid by each member in proportion to the size of its economy. A country's quota determines how much it can borrow and its voting power.
\item \cle{Capacity development}: it offers technical assistance and training (for example in statistics, tax administration and central banking).
\end{enumerate}

\textbf{Conditionalities.} IMF loans are not free gifts. They come with \cle{conditionalities}: policy conditions the borrowing country must implement, typically fiscal discipline (reducing budget deficits), tight monetary policy, removal of subsidies, trade liberalisation and privatisation. In the 1980s and 1990s these were packaged as \cle{Structural Adjustment Programmes} (SAPs). Cameroon implemented SAPs from the late 1980s, involving civil service salary cuts and devaluation of the CFA franc in 1994.

\begin{attention}[Evaluation of conditionalities]
Supporters argue conditionalities restore macroeconomic discipline and confidence. Critics argue they are socially costly (cuts in health and education spending, lay-offs), are imposed uniformly without regard to local conditions, and can deepen recession in the short run. A good GCE essay \emph{evaluates both sides}.
\end{attention}

\courssub{The World Bank (IBRD and IDA)}

The World Bank was also created at Bretton Woods in 1944 and is headquartered in Washington, D.C. It is a family of institutions; the two most important for our syllabus are:

\begin{itemize}
\item The \cle{IBRD} (International Bank for Reconstruction and Development, 1944): originally created to rebuild post-war Europe, it now lends to \emph{middle-income} countries at near-market rates, raising funds by selling bonds on world capital markets.
\item The \cle{IDA} (International Development Association, 1960): provides \emph{concessional} loans (very low interest, long repayment periods) and grants to the \emph{poorest} countries, including most of Sub-Saharan Africa.
\end{itemize}

\begin{definition}[The World Bank]
The \cle{World Bank} is an international financial institution whose central goal is the \cle{reduction of poverty} and the promotion of long-term \cle{economic development}, through financing investment projects (roads, dams, schools, health, agriculture) and supporting policy and institutional reforms in developing countries.
\end{definition}

\begin{aretenir}[IMF versus World Bank -- do not confuse them]
\begin{itemize}
\item \textbf{IMF}: short/medium-term help for \emph{balance of payments} problems; macroeconomic stability; lends to governments (treasuries/central banks); money from quotas.
\item \textbf{World Bank}: long-term finance for \emph{development projects} and poverty reduction; lends for specific projects (infrastructure, education); money from world capital markets and donor contributions.
\item Memory aid: \textbf{IMF = fire brigade} (puts out payments crises); \textbf{World Bank = builder} (finances long-run development).
\end{itemize}
\end{aretenir}

\courssub{The World Trade Organization (WTO)}

The \cle{WTO}, created in 1995 as the successor to the GATT (General Agreement on Tariffs and Trade, 1947), is headquartered in Geneva, Switzerland. It is the only global organisation dealing with the rules of trade between nations.

\textbf{Key principles:}
\begin{enumerate}
\item \cle{Non-discrimination}: the \emph{Most-Favoured-Nation} (MFN) rule -- a trade concession granted to one member must be extended to all; and \emph{national treatment} -- imported goods must be treated like domestic goods once inside the market.
\item \cle{Freer trade}: progressive reduction of tariffs and other barriers through negotiating ``rounds''.
\item \cle{Predictability}: members ``bind'' their tariffs and commit not to raise them arbitrarily.
\item \cle{Fair competition} and \cle{special treatment for developing countries} (longer transition periods).
\end{enumerate}

The WTO also operates a \cle{dispute settlement mechanism}, acting as a referee when members accuse each other of breaking the rules. \emph{Evaluation:} the WTO has expanded world trade enormously, but critics argue that rich countries protect their agriculture while demanding openness from poor countries, and that small states like Cameroon have limited influence in negotiations.

\courssub{Other International Institutions}

\begin{itemize}
\item \cle{UNCTAD} (United Nations Conference on Trade and Development, Geneva, 1964): the UN body defending the trade and development interests of developing countries; it campaigns for fairer commodity prices and better terms of trade -- directly relevant to cocoa and coffee exporters like Cameroon.
\item \cle{AfDB} (African Development Bank, headquartered in Abidjan, C\^ote d'Ivoire, founded 1964): finances development projects and programmes across Africa (transport corridors, energy, agriculture), including many projects in Cameroon.
\item \cle{BEAC} (Bank of Central African States, headquartered in Yaound\'e, Cameroon): the common central bank of the six CEMAC states; it issues the \cle{CFA franc (XAF)}, conducts monetary policy and manages the region's foreign exchange reserves.
\end{itemize}

\begin{savaistu}[Exam tip: dates and headquarters]
Examiners frequently ask for the headquarters and founding year of each institution. Learn this table:
\begin{center}
\begin{tabularx}{\linewidth}{|l|X|l|}
\hline
\rowcolor{popblueL} \textbf{Institution} & \textbf{Headquarters} & \textbf{Founded} \\
\hline
IMF & Washington, D.C. (USA) & 1944/45 \\
\hline
World Bank (IBRD) & Washington, D.C. (USA) & 1944 \\
\hline
WTO & Geneva (Switzerland) & 1995 (GATT 1947) \\
\hline
UNCTAD & Geneva (Switzerland) & 1964 \\
\hline
AfDB & Abidjan (C\^ote d'Ivoire) & 1964 \\
\hline
BEAC & Yaound\'e (Cameroon) & 1972 (as BEAC) \\
\hline
\end{tabularx}
\end{center}
\end{savaistu}

\courssub{Regional Economic Integration: Meaning and Stages}

\begin{definition}[Economic integration]
\cle{Economic integration} is the process by which two or more countries agree to remove barriers to trade and factor movements among themselves and to coordinate their economic policies, in order to create a larger, more competitive economic space.
\end{definition}

\textbf{Objectives} of integration include: enlarging the market for member states' products, achieving economies of scale, attracting investment, strengthening bargaining power in world negotiations, and promoting peace and political stability.

Integration proceeds through \cle{stages}, each deeper than the previous one, as shown in the staircase below.

\begin{popfigure}
\begin{tikzpicture}[scale=1.0, every node/.style={font=\small}]
% staircase
\fill[popblueL] (0,0) rectangle (2.4,1);
\fill[popblue!70!white] (2.4,0) rectangle (4.8,2);
\fill[popgreen!70!white] (4.8,0) rectangle (7.2,3);
\fill[poporange!80!white] (7.2,0) rectangle (9.6,4);
\draw[thick, popdark] (0,0) -- (0,1) -- (2.4,1) -- (2.4,2) -- (4.8,2) -- (4.8,3) -- (7.2,3) -- (7.2,4) -- (9.6,4);
\node[align=center, text width=2.1cm] at (1.2,0.5) {\textbf{Free trade area}};
\node[align=center, text width=2.1cm] at (3.6,1.0) {\textbf{Customs union}};
\node[align=center, text width=2.1cm] at (6.0,1.5) {\textbf{Common market}};
\node[align=center, text width=2.1cm] at (8.4,2.0) {\textbf{Economic \& monetary union}};
% annotations
\node[align=center, text width=2.6cm, font=\scriptsize] at (1.2,1.45) {Internal tariffs removed};
\node[align=center, text width=2.6cm, font=\scriptsize] at (3.6,2.45) {+ Common external tariff};
\node[align=center, text width=2.6cm, font=\scriptsize] at (6.0,3.45) {+ Free movement of labour \& capital};
\node[align=center, text width=2.6cm, font=\scriptsize] at (8.4,4.45) {+ Common currency \& policies};
% example labels
\node[font=\scriptsize, popmuted] at (1.2,-0.4) {e.g. AfCFTA (goal)};
\node[font=\scriptsize, popmuted] at (3.6,-0.4) {e.g. CEMAC};
\node[font=\scriptsize, popmuted] at (6.0,-0.4) {e.g. EU (single market)};
\node[font=\scriptsize, popmuted] at (8.4,-0.4) {e.g. Eurozone, CEMAC};
\end{tikzpicture}
\legende{The staircase of economic integration: each stage adds a deeper level of cooperation.}
\end{popfigure}

\begin{enumerate}
\item \cle{Free trade area}: members abolish tariffs and quotas on trade \emph{among themselves}, but each keeps its own external tariffs against non-members (e.g. the AfCFTA's ambition).
\item \cle{Customs union}: a free trade area \emph{plus} a \cle{common external tariff} (CET) applied by all members to imports from the rest of the world (e.g. CEMAC).
\item \cle{Common market}: a customs union \emph{plus} free movement of factors of production -- labour and capital -- between members (e.g. the EU single market).
\item \cle{Economic and monetary union}: a common market \emph{plus} harmonised economic policies and a \cle{common currency} managed by a common central bank (e.g. the Eurozone; CEMAC with the CFA franc and the BEAC).
\end{enumerate}

\courssub{Regional Groupings Relevant to Cameroon}

\begin{popfigure}
\begin{tikzpicture}[scale=0.85, every node/.style={font=\small}]
% schematic map of Central Africa
\draw[thick, popdark, fill=popblueL] (0,0) -- (0,4.6) -- (2.2,5.2) -- (4.6,4.8) -- (5.4,3.4) -- (5.0,1.4) -- (3.4,0) -- cycle;
% Cameroon highlighted
\draw[thick, poporange, fill=poporange!40!white] (0.3,2.6) -- (0.3,4.4) -- (2.0,4.9) -- (3.0,4.0) -- (2.4,2.8) -- cycle;
\node[font=\small\bfseries, popdark] at (1.55,3.8) {Cameroon};
\node[font=\scriptsize, popmuted] at (1.55,3.45) {(CEMAC + CEEAC)};
% neighbours
\node[align=center, font=\scriptsize] at (0.9,1.6) {Gabon\\ (C+E)};
\node[align=center, font=\scriptsize] at (2.6,1.0) {Congo\\ (C+E)};
\node[align=center, font=\scriptsize] at (3.6,2.2) {CAR\\ (C+E)};
\node[align=center, font=\scriptsize] at (4.4,3.6) {Chad\\ (C+E)};
\node[align=center, font=\scriptsize] at (0.9,0.5) {Eq. Guinea\\ (C+E)};
% outside
\node[align=center, font=\scriptsize, popmuted] at (6.6,4.2) {Nigeria\\ (ECOWAS only)};
\node[align=center, font=\scriptsize, popmuted] at (6.6,1.2) {DRC, Angola,\\ Rwanda, Burundi,\\ S\~ao Tom\'e (CEEAC)};
\draw[->, popmuted] (5.5,3.9) -- (4.8,4.0);
% legend
\node[align=left, font=\scriptsize, anchor=west] at (0,-0.7) {C = CEMAC member \quad E = CEEAC/ECCAS member};
\end{tikzpicture}
\legende{Schematic map: Cameroon belongs to both CEMAC (6 states) and CEEAC/ECCAS (11 states).}
\end{popfigure}

\begin{itemize}
\item \cle{CEMAC} (Economic and Monetary Community of Central Africa, created 1994, successor of UDEAC): Cameroon, Chad, Central African Republic, Equatorial Guinea, Gabon and Congo-Brazzaville. It is both a \emph{customs union} and a \emph{monetary union}: members share the \cle{CFA franc (XAF)} and the common central bank, the BEAC, based in Yaound\'e.
\item \cle{CEEAC / ECCAS} (Economic Community of Central African States, created 1983): a wider grouping of 11 Central African states (the six CEMAC countries plus Angola, Burundi, DRC, Rwanda and S\~ao Tom\'e and Pr\'incipe), aiming at a Central African common market and regional peace.
\item \cle{African Union (AU)} and the \cle{AfCFTA} (African Continental Free Trade Area, launched operationally in 2021, headquartered in Accra, Ghana): the largest free trade area in the world by number of members, aiming to create a single continental market for goods and services. Cameroon has ratified it.
\item \cle{ECOWAS} (Economic Community of West African States, 1975): a West African grouping (Nigeria, Ghana, C\^ote d'Ivoire, etc.). Cameroon is \emph{not} a member, but trades with it -- Nigeria is one of Cameroon's main trading partners.
\item The \cle{European Union (EU)}: the deepest integration scheme in the world (economic and monetary union for the euro area); linked to Cameroon through trade partnership agreements.
\item Cameroon also belongs to the \cle{Commonwealth} and the \cle{Francophonie}, reflecting its bilingual heritage; these foster cultural, educational and economic cooperation rather than formal trade integration.
\end{itemize}

\begin{attention}[Common mistake: CEMAC versus ECOWAS]
Do not write that ``Cameroon is in ECOWAS'' or that ``ECOWAS is Cameroon's monetary union''. \textbf{CEMAC is a monetary union} (common currency, the CFA franc, and a common central bank, the BEAC). \textbf{ECOWAS is not a monetary union} -- its members use different currencies (naira, cedi, CFA franc BCEAO\ldots) and Cameroon is not even a member. Confusing the two costs marks.
\end{attention}

\courssub{Benefits and Problems of Regional Integration for Cameroon}

\textbf{Benefits:}
\begin{enumerate}
\item \cle{Trade creation}: cheaper imports from partners and new export markets replace more expensive domestic or outside production.
\item A \cle{larger market} allows Cameroonian firms (e.g. in cement, food processing) to exploit \cle{economies of scale}.
\item The common currency (CFA franc) eliminates \cle{exchange rate risk} and transaction costs within CEMAC and anchors price stability.
\item Stronger \cle{collective bargaining power} in world trade negotiations (e.g. through the AfCFTA).
\item Free movement of labour and capital encourages \cle{investment} and skills transfer; regional infrastructure (roads, pipelines) becomes viable.
\end{enumerate}

\textbf{Problems:}
\begin{enumerate}
\item \cle{Trade diversion}: members may switch from the cheapest world supplier to a more expensive partner because of the common external tariff, reducing welfare.
\item \cle{Loss of policy sovereignty}: Cameroon cannot conduct its own monetary policy or devalue the CFA franc alone -- the BEAC decides for all six states.
\item \cle{Unequal gains}: the more industrialised member (often Cameroon itself in CEMAC) may dominate, while fiscal revenues fall as tariffs are removed -- Cameroon depends significantly on customs duties.
\item \cle{Similar economic structures}: CEMAC members all export raw materials (oil, cocoa, timber) and import manufactures, so there is little to trade \emph{with each other}; intra-regional trade remains very low.
\item Poor infrastructure, non-tariff barriers (roadblocks, corruption) and political instability slow down real integration.
\end{enumerate}

\begin{exoresolu}[Worked exercise: identifying stages of integration]
\textbf{Statement.} (a) Six Central African states abolish tariffs among themselves, adopt a common external tariff, share one currency issued by a single central bank in Yaound\'e, but labour mobility remains limited in practice. At which stage(s) of integration is this grouping? Justify. (b) Explain one benefit and one cost of the common currency for Cameroon.
\tcblower
\textbf{Solution.} (a) The grouping described is \textbf{CEMAC}. The abolition of internal tariffs plus a common external tariff makes it a \textbf{customs union}; the shared currency (CFA franc) and single central bank (BEAC) make it a \textbf{monetary union}. Because labour and capital do not yet move perfectly freely, it is not yet a complete common market. (b) \emph{Benefit}: the common currency removes exchange rate risk and conversion costs on trade between Cameroon and, say, Gabon, and the BEAC's discipline has kept inflation relatively low. \emph{Cost}: Cameroon loses an independent monetary and exchange rate policy -- it cannot devalue to correct its own balance of payments deficit (as seen in the previous section), since any devaluation (like that of 1994) must be decided collectively for the whole zone.
\end{exoresolu}

\begin{exercice}[Practice]
\begin{enumerate}
\item Distinguish clearly between the roles of the IMF and the World Bank, giving one Cameroonian example of the intervention of each.
\item List the four stages of economic integration in ascending order of depth, and state the additional feature introduced at each stage.
\item ``Regional integration is beneficial to Cameroon.'' Discuss, using CEMAC and the AfCFTA as illustrations.
\end{enumerate}
\end{exercice}

\begin{corrige}[Exercise 1]
\textbf{1.} IMF: short/medium-term balance of payments support with conditionalities -- e.g. the Structural Adjustment Programmes and the 1994 CFA franc devaluation context. World Bank: long-term project finance for development -- e.g. financing of roads, energy and agricultural projects in Cameroon. \textbf{2.} Free trade area (internal tariffs removed) $\rightarrow$ customs union (+ common external tariff) $\rightarrow$ common market (+ free movement of labour and capital) $\rightarrow$ economic and monetary union (+ common currency and harmonised policies). \textbf{3.} Essay plan: \emph{benefits} -- trade creation, larger market, monetary stability via the CFA franc, bargaining power via AfCFTA; \emph{costs} -- trade diversion, loss of monetary sovereignty, low intra-CEMAC trade due to similar export structures, loss of customs revenue. \emph{Conclusion}: beneficial in principle, but gains depend on diversification of Cameroon's economy and effective implementation.
\end{corrige}

\coursec{Foreign Exchange Market, Exchange Rates and Adjustments}

Having examined the role of international economic institutions and regional groupings in shaping Cameroon's trade environment, we now turn to the mechanism that makes cross-border transactions possible: the \cle{foreign exchange market}. Every import of French wine, every export of cocoa beans, every repatriation of profits by a multinational, and every tourist's spending in Yaound\'e ultimately requires a conversion of one currency into another. This section explains how that market works, how exchange rates are quoted and determined, why they move, and what the consequences are for the Cameroonian economy.

\courssub{The Foreign Exchange Market and its Participants}

\begin{definition}[Foreign Exchange Market]
The \textbf{foreign exchange market (forex market)} is the global decentralised marketplace where currencies are bought and sold. It operates 24 hours a day, five days a week, through a network of banks, brokers, electronic platforms and official institutions. Its primary function is to enable the conversion of one currency into another so that international trade, investment and tourism can take place.
\end{definition}

The main participants are:
\begin{itemize}
\item \textbf{Commercial banks} -- the core dealers who quote buying (bid) and selling (ask) prices and execute the bulk of transactions.
\item \textbf{Firms (importers/exporters)} -- importers demand foreign currency to pay for goods bought abroad; exporters supply the foreign currency they have earned from their sales.
\item \textbf{Central banks} (e.g. the \emph{BEAC} for the CEMAC zone) -- they manage foreign reserves, implement monetary policy and, under a fixed regime, guarantee the official parity.
\item \textbf{Speculators and investment funds} -- they buy and sell currencies hoping to profit from short-term exchange rate movements.
\item \textbf{Bureaux de change} -- retail outlets that serve tourists and small businesses.
\end{itemize}

\begin{attention}[Who Demands and Who Supplies Foreign Currency?]
This is a frequent source of confusion. \textbf{Demand} for foreign currency comes from \emph{domestic residents}: importers paying foreign suppliers, tourists travelling abroad, firms investing overseas. \textbf{Supply} of foreign currency comes from \emph{foreigners}: buyers of our exports, foreign tourists spending in Cameroon, foreign investors bringing capital into the country. Always state this explicitly in essays; examiners award marks for identifying the correct agents.
\end{attention}

\courssub{Exchange Rate: Definition, Quotation and Conversion}

\begin{definition}[Exchange Rate]
An \textbf{exchange rate} is the price of one currency expressed in terms of another. If the exchange rate of the euro against the FCFA is 655.957, then $1\ \text{EUR} = 655.957\ \text{FCFA}$.
\end{definition}

Two quotation methods exist:
\begin{itemize}
\item \textbf{Direct quotation} (price quotation): the number of units of \emph{domestic} currency per unit of \emph{foreign} currency. In Cameroon the direct quote is $\text{FCFA}/\text{EUR} = 655.957$.
\item \textbf{Indirect quotation} (volume quotation): the number of units of \emph{foreign} currency per unit of \emph{domestic} currency. The indirect quote is $\text{EUR}/\text{FCFA} = 1/655.957 \approx 0.001524$.
\end{itemize}

\begin{methode}[Converting between Currencies]
\begin{enumerate}
\item Identify the quotation method used (direct or indirect).
\item If the rate is given as \emph{domestic per foreign} (direct), then:
\begin{itemize}
\item To convert \emph{foreign $\to$ domestic}: multiply the foreign amount by the rate.
\item To convert \emph{domestic $\to$ foreign}: divide the domestic amount by the rate.
\end{itemize}
\item If the rate is given as \emph{foreign per domestic} (indirect), reverse the operations.
\item Always state the currency of the result and round sensibly.
\end{enumerate}
\end{methode}

\begin{exoresolu}[Currency Conversion]
\textbf{Statement:} The BEAC publishes a direct quote of $1\ \text{EUR} = 655.957\ \text{FCFA}$. (a) A Cameroonian importer must pay \texteuro 12,500 for a shipment of machinery. How many FCFA are needed? (b) A cocoa exporter receives 45,000,000 FCFA. How many euros does this represent?

\tcblower
\textbf{Solution:}
\begin{enumerate}
\item[(a)] Foreign currency needed = \texteuro 12,500. Using the direct quote (foreign $\to$ domestic: multiply):
\[ \text{FCFA} = 12,500 \times 655.957 = 8,199,462.5\ \text{FCFA}. \]
The importer needs \textbf{8,199,462.5 FCFA} (about 8.2 million FCFA).
\item[(b)] Domestic currency received = 45,000,000 FCFA. Convert to euros (domestic $\to$ foreign: divide):
\[ \text{EUR} = 45,000,000 \div 655.957 \approx 68,602.1\ \text{EUR}. \]
The exporter's revenue represents approximately \textbf{\texteuro 68,602}.
\end{enumerate}
\end{exoresolu}

\courssub{Determination of Exchange Rates under a Floating System}

In a \textbf{floating (flexible) exchange rate system}, the equilibrium rate is determined by the interaction of the \cle{demand for} and the \cle{supply of} foreign currency, exactly as the price of any good is determined in a product market. The demand for foreign currency comes from agents who need it to buy foreign goods, services or assets; the supply comes from foreigners who need domestic currency to buy our exports or domestic assets.

\begin{popfigure}
\begin{tikzpicture}[scale=0.9, every node/.style={font=\small}]
\draw[->] (0,0) -- (8,0) node[right] {$Q$ of foreign currency};
\draw[->] (0,0) -- (0,5) node[above] {$E$ (FCFA per unit of foreign currency)};
\draw[popblue, thick] (0.5,4.5) .. controls (2,3.5) and (4,2.5) .. (7,1) node[right, popblue] {$D$};
\draw[popgreen, thick] (0.5,0.5) .. controls (2,1.5) and (4,2.5) .. (7,4) node[right, popgreen] {$S$};
\draw[dashed, popdark] (4,0) -- (4,2.5) -- (0,2.5);
\node[popdark] at (4,-0.3) {$Q^*$};
\node[popdark, left] at (0,2.5) {$E^*$};
\node[align=center] at (4,4.8) {Equilibrium in the foreign exchange market};
\end{tikzpicture}
\legende{Demand and supply in the foreign exchange market. The equilibrium exchange rate $E^*$ clears the market at quantity $Q^*$.}
\end{popfigure}

An increase in the demand for foreign currency (e.g. a surge in imports) shifts the demand curve to the right, raising the equilibrium exchange rate. Because the exchange rate is expressed as domestic currency per unit of foreign currency, a higher $E$ means that more FCFA are needed to buy one unit of foreign currency: the domestic currency has \cle{depreciated}.

\begin{popfigure}
\begin{tikzpicture}[scale=0.9, every node/.style={font=\small}]
\draw[->] (0,0) -- (8.5,0) node[right] {$Q$};
\draw[->] (0,0) -- (0,5) node[above] {$E$};
\draw[popblue, thick] (0.5,4.5) .. controls (2,3.5) and (4,2.5) .. (7,1) node[right, popblue] {$D_1$};
\draw[poporange, thick, dashed] (1.5,4.5) .. controls (3,3.5) and (5,2.5) .. (8,1) node[right, poporange] {$D_2$};
\draw[popgreen, thick] (0.5,0.5) .. controls (2,1.5) and (4,2.5) .. (7,4) node[right, popgreen] {$S$};
\draw[dashed, popdark] (4,0) -- (4,2.5) -- (0,2.5);
\draw[dashed, poporange] (5.5,0) -- (5.5,3.2) -- (0,3.2);
\node[popdark] at (4,-0.3) {$Q_1^*$};
\node[poporange] at (5.5,-0.3) {$Q_2^*$};
\node[popdark, left] at (0,2.5) {$E_1^*$};
\node[poporange, left] at (0,3.2) {$E_2^*$};
\draw[->, poporange, thick] (0.6,2.5) -- (0.6,3.2) node[midway, left, align=center] {rise in $E$};
\node[align=center] at (4,4.8) {Rightward shift of demand: higher $E$ (depreciation)};
\end{tikzpicture}
\legende{Effect of an increase in demand for foreign currency: the exchange rate rises from $E_1^*$ to $E_2^*$, a depreciation of the domestic currency.}
\end{popfigure}

Symmetrically, an increase in the \emph{supply} of foreign currency (e.g. a boom in cocoa exports) shifts the supply curve to the right, lowering $E$: the domestic currency \cle{appreciates}.

\courssub{Causes of Exchange Rate Changes}

The main drivers of exchange rate movements in a floating regime are:
\begin{itemize}
\item \textbf{Trade flows:} a rise in exports increases the supply of foreign currency (appreciation); a rise in imports increases the demand for it (depreciation).
\item \textbf{Inflation differentials:} higher domestic inflation than abroad reduces the purchasing power of the domestic currency, leading to depreciation (the \emph{purchasing power parity} approach).
\item \textbf{Interest rate differentials:} higher domestic interest rates attract foreign capital, increasing the demand for the domestic currency (appreciation) -- the \cle{interest rate parity} condition.
\item \textbf{Speculation:} expectations of future appreciation or depreciation can trigger self-fulfilling short-run movements.
\end{itemize}

\courssub{Exchange Rate Systems}

Three broad regimes exist:
\begin{enumerate}
\item \textbf{Fixed (pegged) exchange rate:} the monetary authority commits to a specific parity and stands ready to buy or sell foreign currency to maintain it. The CFA franc is a classic example.
\item \textbf{Floating exchange rate:} the rate is determined solely by market forces, without official intervention.
\item \textbf{Managed float (dirty float):} the rate is primarily market-determined, but the central bank occasionally intervenes to smooth excessive volatility.
\end{enumerate}

\begin{aretenir}[The CFA Franc Peg]
The \textbf{FCFA} (franc de la coop\'eration financi\`ere en Afrique centrale, issued by the BEAC) is pegged to the euro at a fixed rate of \textbf{655.957 FCFA = 1 EUR}. The peg is guaranteed by the French Treasury, which ensures unlimited convertibility of the CFA franc into euros; in return, the BEAC holds a substantial part of its foreign reserves (mainly in euros) in an operations account with the French Treasury. This regime eliminates exchange rate risk vis-\`a-vis the euro and anchors price stability, but it means the CEMAC countries cannot use the exchange rate as an independent instrument of adjustment.
\end{aretenir}

\courssub{Exchange Rate Adjustments}

When the value of a currency changes, the terminology depends on the regime:
\begin{itemize}
\item \textbf{Fixed regime:} \cle{devaluation} = official reduction in the value of the domestic currency decided by the authorities (e.g. the CFA franc devaluation of January 1994, when the parity moved from 50 to 100 FCFA per French franc). \cle{Revaluation} = official increase in value.
\item \textbf{Floating regime:} \cle{depreciation} = market-driven fall in value. \cle{Appreciation} = market-driven rise in value.
\end{itemize}

\begin{attention}[Devaluation vs Depreciation -- Common Mistake]
\textbf{Devaluation} and \textbf{depreciation} both mean the domestic currency becomes cheaper, but they are \emph{not synonyms}. Devaluation is a deliberate policy act under a fixed peg; depreciation is the outcome of market forces under a float. Similarly, \textbf{revaluation} is an official upward adjustment, while \textbf{appreciation} is market-driven. In exams, always specify the regime when using these terms.
\end{attention}

\textbf{Effects on trade and inflation:}
\begin{itemize}
\item A \textbf{depreciation/devaluation} makes exports cheaper for foreigners (export volume tends to rise) and imports more expensive for domestic residents (import volume tends to fall). This improves the trade balance \emph{if} the \cle{Marshall--Lerner condition} holds, i.e. if the sum of the price elasticities of demand for exports and imports is greater than 1. However, it also raises the domestic price of imported inputs, feeding \cle{cost-push inflation}.
\item An \textbf{appreciation/revaluation} has the opposite effects: exports become dearer, imports cheaper, potentially worsening the trade balance but dampening inflation.
\end{itemize}

\courssub{Effects on the Balance of Payments and the Cameroonian Economy}

Cameroon operates under the CFA franc fixed peg, so the exchange rate does not adjust automatically to balance of payments shocks. Instead, adjustment occurs through:
\begin{itemize}
\item \textbf{Reserve flows:} a current account deficit leads to a loss of foreign reserves at the BEAC; a surplus builds reserves up.
\item \textbf{Domestic credit and monetary policy:} the BEAC can tighten or loosen credit to influence domestic demand and thus the demand for imports.
\item \textbf{Structural policies:} diversification of exports (beyond cocoa, oil and timber), improvement of competitiveness, and attraction of foreign direct investment are the long-run solutions.
\end{itemize}

A hypothetical \textbf{devaluation of the CFA franc} (if the parity were changed) would:
\begin{itemize}
\item increase the FCFA price of imported fuel, machinery and food $\rightarrow$ higher inflation, hurting low-income households;
\item make Cameroonian cocoa, cotton and timber cheaper on world markets $\rightarrow$ higher export volumes and possibly higher export earnings in FCFA terms;
\item increase the FCFA value of external debt service (since the debt is denominated in foreign currencies) $\rightarrow$ a heavier fiscal burden.
\end{itemize}

\begin{experience}[Case Study: Impact of a 10\% Depreciation of the FCFA (Hypothetical)]
Assume the parity is adjusted from 655.957 to 721.55 FCFA/EUR, i.e. a 10\% depreciation ($655.957 \times 1.10 \approx 721.55$). Consider a cocoa exporter who sells one tonne at \texteuro 2,000.
\begin{itemize}
\item \textbf{Before:} revenue $= 2,000 \times 655.957 = 1,311,914$ FCFA.
\item \textbf{After:} revenue $= 2,000 \times 721.55 = 1,443,100$ FCFA.
\end{itemize}
The exporter gains $1,443,100 - 1,311,914 = 131,186$ FCFA per tonne in domestic currency. However, the same exporter pays more for imported fertiliser (priced in euros). If fertiliser costs \texteuro 500 per tonne, its cost rises from $500 \times 655.957 = 327,978.5$ FCFA to $500 \times 721.55 = 360,775$ FCFA. The net effect depends on the share of imported inputs in production costs. This illustrates why the BEAC and the Cameroonian government carefully weigh any parity change.
\end{experience}

\begin{savaistu}[Did You Know? The 1994 Devaluation]
In January 1994, the CFA franc was devalued by 50\% against the French franc (from 50 to 100 FCFA per French franc), the first parity change since 1948. The aim was to restore the competitiveness of the CFA-zone economies, including Cameroon, whose exports had become overvalued. The devaluation boosted cash-crop exports such as cocoa and cotton but also caused a sharp, temporary rise in the prices of imported goods.
\end{savaistu}

\courssub{Summary of Key Points}

\begin{aretenir}[Key Points for Revision]
\begin{itemize}
\item The forex market converts currencies; participants include banks, firms, central banks, speculators and bureaux de change.
\item Exchange rate = price of one currency in terms of another; direct quote = domestic per foreign, indirect quote = foreign per domestic.
\item Under a float, equilibrium is where demand for foreign currency equals supply; shifts in demand or supply cause depreciation or appreciation.
\item Main causes of rate changes: trade flows, inflation differentials, interest rate differentials, speculation.
\item Regimes: fixed (the CFA franc is pegged to the euro at 655.957), floating, managed float.
\item Adjustments: devaluation/revaluation (fixed regime) vs depreciation/appreciation (floating regime) -- they are not synonyms.
\item Depreciation improves export competitiveness but raises import prices and inflation; appreciation does the opposite. The trade balance improves only if the Marshall--Lerner condition holds.
\item For Cameroon, the fixed peg means external adjustment occurs via reserves and domestic policy, not through the exchange rate itself.
\end{itemize}
\end{aretenir}

\newpage

\coursec{Integration activity}

\begin{exercice}[Exercice d'application]
Énoncé de l'exercice : analyser la situation économique proposée en mobilisant les notions du cours.
\tcblower
Solution détaillée : 
\begin{enumerate}
    \item Identifier les agents et les marchés concernés.
    \item Appliquer le modèle d'offre et de demande ou le modèle IS-LM selon le contexte.
    \item Calculer les équilibres et interpréter les résultats économiques.
    \item Discuter les limites de l'analyse (hypothèses, chocs exogènes).
\end{enumerate}
\end{exercice}

\newpage

\coursec{Methods and worked examples}

\courssub{Methods and Worked Examples}

This section presents the essential techniques you must master to solve GCE Advanced Level International Economics problems. Each method is broken down into clear, examinable steps, followed by a fully worked examination-style question. Pay close attention to the \textbf{layout of calculations}, the \textbf{terminology} (e.g., "opportunity cost", "terms of trade", "current account"), and the \textbf{evaluation points} which earn the higher marks.

\begin{methode}[Calculating Comparative Advantage and Determining Specialisation]
\textbf{Context:} Practical Work 23/24 (TU 47/48). This is the core numerical skill for the Theories of Trade.
\textbf{Steps:}
\begin{enumerate}
    \item \textbf{Set up the Production Possibility Table:} Organise data clearly: Country (or Region) vs. Goods. Units = Output per unit of resource (e.g., per worker, per hectare, per day) \emph{or} Input required per unit of output. \textbf{Check which one is given!}
    \item \textbf{Calculate Opportunity Costs (OC):}
        \begin{itemize}
            \item If \textbf{Output method} (Output per resource): $OC(\text{Good A}) = \frac{\text{Output of Good B given up}}{\text{Output of Good A gained}}$.
            \item If \textbf{Input method} (Resource per unit output): $OC(\text{Good A}) = \frac{\text{Input for Good A}}{\text{Input for Good B}}$.
        \end{itemize}
    \item \textbf{Compare Opportunity Costs:} The country with the \textbf{lower} OC for a good has \textbf{Comparative Advantage} (CA) in that good.
    \item \textbf{Determine Specialisation:} Each country specialises \textbf{completely} in the good where it has CA (assuming constant opportunity costs / linear PPF).
    \item \textbf{Calculate Total World Output Before \& After:} Show that total output of at least one good rises (or both) without using more resources $\rightarrow$ \textbf{Gains from Trade}.
    \item \textbf{Determine Terms of Trade (Exchange Rate):} The trading ratio (Price of Good A / Price of Good B) must lie \textbf{between the two domestic opportunity cost ratios} for both countries to gain.
\end{enumerate}
\textbf{Examiner's Reflex:} Always state "Country X has a comparative advantage in Good Y because its opportunity cost (Z) is lower than Country W's (K)." Never just say "it produces more."
\end{methode}

\begin{exoresolu}[Comparative Advantage: Cameroon and Nigeria — Cocoa and Petroleum]
\textbf{Statement:} Cameroon and Nigeria produce only Cocoa and Crude Petroleum using labour as the only input. The output per worker per day is given below:
\begin{center}
\begin{tabular}{|c|c|c|}\hline
\textbf{Country} & \textbf{Cocoa (tonnes)} & \textbf{Petroleum (barrels)} \\ \hline
Cameroon & 10 & 5 \\ \hline
Nigeria & 4 & 20 \\ \hline
\end{tabular}
\end{center}
\begin{enumerate}
    \item Calculate the opportunity cost of producing 1 tonne of Cocoa in each country. \textbf{(2 marks)}
    \item Determine which country has a comparative advantage in Cocoa and which in Petroleum. \textbf{(2 marks)}
    \item If each country has 100 workers and initially splits labour equally (50 each), calculate total world output before specialisation. \textbf{(2 marks)}
    \item If they specialise completely according to comparative advantage, calculate the new total world output. \textbf{(2 marks)}
    \item Identify a mutually beneficial Terms of Trade (exchange ratio) for 1 tonne of Cocoa in barrels of Petroleum. \textbf{(2 marks)}
    \item Explain \textbf{one} reason why complete specialisation might not occur in reality. \textbf{(2 marks)}
\end{enumerate}

\tcblower
\textbf{Detailed Solution:}

\textbf{1. Opportunity Cost Calculations (Output Method)}
Formula: $OC(\text{Cocoa}) = \frac{\text{Petroleum given up}}{\text{Cocoa gained}}$ per worker shift.

\begin{itemize}
    \item \textbf{Cameroon:} 1 worker produces 10 tonnes Cocoa \textbf{OR} 5 barrels Petroleum.
    $OC_{\text{Cam}}(\text{1 tonne Cocoa}) = \frac{5 \text{ barrels}}{10 \text{ tonnes}} = \mathbf{0.5 \text{ barrels of Petroleum}}$.
    \item \textbf{Nigeria:} 1 worker produces 4 tonnes Cocoa \textbf{OR} 20 barrels Petroleum.
    $OC_{\text{Nig}}(\text{1 tonne Cocoa}) = \frac{20 \text{ barrels}}{4 \text{ tonnes}} = \mathbf{5 \text{ barrels of Petroleum}}$.
\end{itemize}

\textbf{2. Comparative Advantage}
\begin{itemize}
    \item \textbf{Cocoa:} Cameroon OC = 0.5 barrels; Nigeria OC = 5 barrels. \textbf{Cameroon has lower OC $\rightarrow$ CA in Cocoa.}
    \item \textbf{Petroleum:} Calculate OC of 1 barrel Petroleum (inverse).
    $OC_{\text{Cam}}(\text{1 barrel Petrol}) = \frac{10}{5} = 2 \text{ tonnes Cocoa}$.
    $OC_{\text{Nig}}(\text{1 barrel Petrol}) = \frac{4}{20} = 0.2 \text{ tonnes Cocoa}$.
    Nigeria has lower OC (0.2 < 2) $\rightarrow$ \textbf{Nigeria has CA in Petroleum}.
\end{itemize}

\textbf{3. World Output Before Specialisation (50 workers each)}
\begin{center}
\begin{tabular}{|c|c|c|c|c|}\hline
 & \multicolumn{2}{c|}{\textbf{Cameroon (50W each)}} & \multicolumn{2}{c|}{\textbf{Nigeria (50W each)}} \\ \hline
\textbf{Good} & \textbf{Output/Worker} & \textbf{Total} & \textbf{Output/Worker} & \textbf{Total} \\ \hline
Cocoa & 10 & $50 \times 10 = 500$ & 4 & $50 \times 4 = 200$ \\ \hline
Petroleum & 5 & $50 \times 5 = 250$ & 20 & $50 \times 20 = 1000$ \\ \hline
\end{tabular}
\end{center}
\textbf{Total World Output:} Cocoa = \textbf{700 tonnes}; Petroleum = \textbf{1250 barrels}.

\textbf{4. World Output After Complete Specialisation}
\begin{itemize}
    \item Cameroon (100 workers) $\rightarrow$ Cocoa only: $100 \times 10 = \mathbf{1000 \text{ tonnes Cocoa}}$.
    \item Nigeria (100 workers) $\rightarrow$ Petroleum only: $100 \times 20 = \mathbf{2000 \text{ barrels Petroleum}}$.
\end{itemize}
\textbf{Total World Output:} Cocoa = \textbf{1000 tonnes} (+300); Petroleum = \textbf{2000 barrels} (+750).
\textbf{Gain:} World output of \textbf{both} goods has increased. This is the gain from trade.

\textbf{5. Mutually Beneficial Terms of Trade (TOT)}
Domestic Exchange Ratios (OC of Cocoa):
\begin{itemize}
    \item Cameroon: 1 Cocoa = 0.5 Petroleum (Domestic cost of buying petrol).
    \item Nigeria: 1 Cocoa = 5 Petroleum (Domestic cost of buying cocoa).
\end{itemize}
For trade to benefit both, the International Exchange Ratio (TOT) must lie \textbf{between} them:
\[ \mathbf{0.5 \text{ barrels} < 1 \text{ tonne Cocoa} < 5 \text{ barrels}} \]
\textbf{Example acceptable TOT:} \textbf{1 tonne Cocoa = 2 barrels Petroleum} (or 1, 3, 4).
\textit{Check:} Cameroon sells cocoa at 2 barrels (gains 1.5 barrels vs domestic 0.5). Nigeria buys cocoa at 2 barrels (saves 3 barrels vs domestic 5). Both gain.

\textbf{6. Limits to Complete Specialisation (Evaluation)}
\begin{itemize}
    \item \textbf{Increasing Opportunity Costs (Concave PPF):} Resources are not perfectly substitutable (land specific to cocoa, oil rigs specific to petroleum). As specialisation expands, OC rises, eventually equalling the TOT, stopping further specialisation.
    \item \textbf{Transport Costs:} Shipping cocoa from Cameroon to Nigeria and oil back adds cost, narrowing the "gain zone".
    \item \textbf{Non-traded Goods / Services:} Not all output is tradable (e.g., haircuts, local construction).
    \item \textbf{Strategic / Food Security:} Cameroon may retain some oil refining capacity; Nigeria may retain some cocoa production for domestic stability.
\end{itemize}
\end{exoresolu}

\begin{methode}[Analysing the Effects of a Tariff (Diagram \& Welfare)]
\textbf{Context:} Lesson 102 (TU 49). Standard GCE essay/short answer: "Using a diagram, analyse the effects of a specific tariff on domestic price, quantity, consumer surplus, producer surplus, government revenue, and deadweight loss."
\textbf{Steps:}
\begin{enumerate}
    \item \textbf{Draw Axes:} Price (P) vertical, Quantity (Q) horizontal.
    \item \textbf{Draw Curves:} Downward sloping Domestic Demand ($D_d$), Upward sloping Domestic Supply ($S_d$). World Supply ($S_w$) is \textbf{perfectly elastic (horizontal)} at World Price $P_w$ (Small Country Assumption).
    \item \textbf{Free Trade Equilibrium:} $P_w$. Domestic Qs = $Q_1$ (intersection $S_d$ \& $P_w$). Domestic Qd = $Q_2$ (intersection $D_d$ \& $P_w$). Imports = $Q_2 - Q_1$.
    \item \textbf{Impose Specific Tariff ($t$):} Shift World Supply curve \textbf{upwards by $t$} to $S_w + t$. New Domestic Price $P_t = P_w + t$.
    \item \textbf{New Equilibrium:} New Domestic Qs = $Q_3$ ($S_d$ meets $P_t$). New Domestic Qd = $Q_4$ ($D_d$ meets $P_t$). New Imports = $Q_4 - Q_3$ (falls).
    \item \textbf{Identify Welfare Areas (Label clearly!):}
        \begin{itemize}
            \item \textbf{Consumer Surplus Loss:} Area between $P_w$ and $P_t$ under $D_d$ (Trapezoid $P_w P_t Q_4 Q_2$).
            \item \textbf{Producer Surplus Gain:} Area between $P_w$ and $P_t$ above $S_d$ (Trapezoid $P_w P_t Q_3 Q_1$).
            \item \textbf{Govt Revenue:} $t \times (Q_4 - Q_3)$ (Rectangle).
            \item \textbf{Deadweight Loss (Net Welfare Loss):} Two triangles:
                \begin{itemize}
                    \item \textbf{Production Distortion (Efficiency Loss):} Triangle between $S_d$, $P_w$, $P_t$ (left of $S_d$). Resources shifted from efficient world producers to inefficient domestic producers.
                    \item \textbf{Consumption Distortion:} Triangle between $D_d$, $P_w$, $P_t$ (right of $D_d$). Consumers priced out of market despite willingness to pay $> P_w$.
                \end{itemize}
        \end{itemize}
    \item \textbf{Conclusion:} Net welfare loss = Deadweight Loss Triangles. Tariff protects domestic jobs/profits but at a net cost to society.
\end{enumerate}
\textbf{Key Diagram Labels:} $P_w$, $P_t$, $Q_1, Q_2, Q_3, Q_4$, $S_d$, $D_d$, $S_w$, $S_w+t$, Areas (a,b,c,d or letters).
\end{methode}

\begin{exoresolu}[Tariff Analysis: Cameroon Imports of Frozen Chicken]
\textbf{Statement:} Cameroon imports frozen chicken at a world price of \SI{1000}{FCFA/kg}. Domestic demand and supply are linear:
$Q_d = 1000 - 0.5P$ ; $Q_s = -200 + 0.5P$ (Quantities in tonnes, P in FCFA/kg).
The government imposes a specific tariff of \SI{200}{FCFA/kg}.
\begin{enumerate}
    \item Calculate the free trade equilibrium price, domestic production, consumption, and import volume. \textbf{(3 marks)}
    \item Calculate the new equilibrium after the tariff: domestic price, production, consumption, import volume. \textbf{(3 marks)}
    \item Calculate the change in Consumer Surplus, Producer Surplus, Government Revenue, and Net Welfare Loss (Deadweight Loss). \textbf{(6 marks)}
    \item Briefly explain the "Infant Industry" argument for protecting the local poultry sector. \textbf{(2 marks)}
\end{enumerate}

\tcblower
\textbf{Detailed Solution:}

\textbf{1. Free Trade Equilibrium ($P_w = 1000$)}
\begin{itemize}
    \item \textbf{Price:} $P_{FT} = \mathbf{1000 \text{ FCFA/kg}}$.
    \item \textbf{Domestic Supply ($Q_s$):} $Q_s = -200 + 0.5(1000) = -200 + 500 = \mathbf{300 \text{ tonnes}}$.
    \item \textbf{Domestic Demand ($Q_d$):} $Q_d = 1000 - 0.5(1000) = 1000 - 500 = \mathbf{500 \text{ tonnes}}$.
    \item \textbf{Imports:} $M = Q_d - Q_s = 500 - 300 = \mathbf{200 \text{ tonnes}}$.
\end{itemize}

\textbf{2. Post-Tariff Equilibrium ($t = 200 \rightarrow P_t = 1200$)}
\begin{itemize}
    \item \textbf{Domestic Price:} $P_t = P_w + t = 1000 + 200 = \mathbf{1200 \text{ FCFA/kg}}$.
    \item \textbf{Domestic Supply ($Q_s'$):} $Q_s' = -200 + 0.5(1200) = -200 + 600 = \mathbf{400 \text{ tonnes}}$.
    \item \textbf{Domestic Demand ($Q_d'$):} $Q_d' = 1000 - 0.5(1200) = 1000 - 600 = \mathbf{400 \text{ tonnes}}$.
    \item \textbf{Imports:} $M' = Q_d' - Q_s' = 400 - 400 = \mathbf{0 \text{ tonnes}}$ (Prohibitive Tariff!).
\end{itemize}
\textit{Note: The tariff is high enough to eliminate imports completely. This simplifies welfare calc: Govt Revenue = 0.}

\textbf{3. Welfare Analysis (Areas under linear curves)}
We need inverse demand/supply for height of triangles.
$P_d = 2000 - 2Q_d$ (Max price = 2000 when Q=0).
$P_s = 400 + 2Q_s$ (Min price = 400 when Q=0).

\begin{itemize}
    \item \textbf{Consumer Surplus (CS) Loss:} Area under $D_d$ between $P=1000$ and $P=1200$, from $Q=400$ to $500$.
    Trapezoid area = $\frac{1}{2}(P_{max} - P_t)Q_d' - \frac{1}{2}(P_{max} - P_w)Q_d$? No, standard: $\Delta CS = -\int_{P_w}^{P_t} Q_d dP$.
    Easier: Area of Trapezoid between prices.
    Height = $200$. Parallel sides = $Q_d(P_w)=500$ and $Q_d(P_t)=400$.
    $\Delta CS = -\frac{1}{2}(500 + 400) \times 200 = -\frac{900}{2} \times 200 = \mathbf{-90,000 \text{ FCFA}}$ (Loss).
    
    \item \textbf{Producer Surplus (PS) Gain:} Area above $S_d$ between $P=1000$ and $P=1200$, from $Q=300$ to $400$.
    Trapezoid: Height = $200$. Sides = $Q_s(P_w)=300$, $Q_s(P_t)=400$.
    $\Delta PS = +\frac{1}{2}(300 + 400) \times 200 = +\frac{700}{2} \times 200 = \mathbf{+70,000 \text{ FCFA}}$ (Gain).
    
    \item \textbf{Government Revenue (GR):} $t \times M' = 200 \times 0 = \mathbf{0 \text{ FCFA}}$.
    
    \item \textbf{Deadweight Loss (Net Welfare Loss):} $\Delta CS + \Delta PS + GR$.
    $NWL = -90,000 + 70,000 + 0 = \mathbf{-20,000 \text{ FCFA}}$.
    \textit{Check via Triangles:}
    Production Loss Triangle: Base $\Delta Q_s = 100$, Height $t=200$. Area = $0.5 \times 100 \times 200 = 10,000$.
    Consumption Loss Triangle: Base $\Delta Q_d = 100$, Height $t=200$. Area = $0.5 \times 100 \times 200 = 10,000$.
    Total DWL = \textbf{20,000 FCFA}. Matches.
\end{itemize}

\textbf{4. Infant Industry Argument}
New domestic industries (e.g., Cameroonian poultry farming) face high initial average costs due to lack of economies of scale, learning-by-doing, and undeveloped supply chains. They cannot compete immediately with established, efficient foreign producers (e.g., EU/US/Brazil subsidised exports). Temporary protection (tariff/quota) allows them to expand output, move down the LRAC curve, acquire skills, and eventually compete internationally without protection. \textbf{Key condition:} Protection must be \textbf{temporary} and \textbf{performance-based} (sunset clauses), otherwise it creates permanent inefficiency and rent-seeking.
\end{exoresolu}

\begin{methode}[Calculating and Interpreting Terms of Trade (TOT)]
\textbf{Context:} Lesson 103 (TU 50). TOT = Index of Export Prices / Index of Import Prices $\times 100$.
\textbf{Steps:}
\begin{enumerate}
    \item \textbf{Identify Base Year:} TOT = 100 in base year.
    \item \textbf{Calculate Indices:}
        $I_X = \frac{P_{X,\text{current}}}{P_{X,\text{base}}} \times 100$ (Export Price Index).
        $I_M = \frac{P_{M,\text{current}}}{P_{M,\text{base}}} \times 100$ (Import Price Index).
        \textit{If given value/volume data:} $P_X = \frac{\text{Export Value}}{\text{Export Volume}}$ (Unit Value Index).
    \item \textbf{Compute TOT:} $TOT = \frac{I_X}{I_M} \times 100$.
    \item \textbf{Interpret Change:}
        \begin{itemize}
            \item TOT $> 100$ (Rising): \textbf{Improving/Favourable}. 1 unit of exports buys \textbf{more} imports. Real income $\uparrow$.
            \item TOT $< 100$ (Falling): \textbf{Deteriorating/Unfavourable}. 1 unit of exports buys \textbf{fewer} imports. Real income $\downarrow$.
        \end{itemize}
    \item \textbf{Link to Elasticities (Marshall-Lerner / Prebisch-Singer):}
        \begin{itemize}
            \item If TOT improves, does Export Revenue rise? Depends on PED$_X$. If PED$_X$ elastic ($>1$), volume falls $\% >$ price rises $\% \rightarrow$ Revenue $\downarrow$.
            \item Prebisch-Singer Hypothesis: Primary commodity exporters (Cameroon: cocoa, timber, oil, cotton) face secularly declining TOT vs manufactured goods (low YED, low PED, synthetic substitutes).
        \end{itemize}
\end{enumerate}
\textbf{Exam Trap:} "Improving TOT" $\neq$ "Improving Trade Balance". Trade Balance = $P_X Q_X - P_M Q_M$. Volume effects matter!
\end{methode}

\begin{exoresolu}[Terms of Trade Calculation and Prebisch-Singer Application]
\textbf{Statement:} Cameroon exports Cocoa and Petroleum. It imports Machinery and Food. Base Year = 2020.
\begin{center}
\begin{tabular}{|c|c|c|c|c|}\hline
\textbf{Year} & \textbf{Export Price Index ($I_X$)} & \textbf{Import Price Index ($I_M$)} & \textbf{Export Volume Index} & \textbf{Import Volume Index} \\ \hline
2020 (Base) & 100 & 100 & 100 & 100 \\ \hline
2023 & 130 & 115 & 90 & 110 \\ \hline
\end{tabular}
\end{center}
\begin{enumerate}
    \item Calculate the Terms of Trade (TOT) for 2023. \textbf{(2 marks)}
    \item State whether the TOT has improved or deteriorated and explain the meaning for Cameroon's real income. \textbf{(2 marks)}
    \item Calculate the Index of Export Value and Import Value for 2023. \textbf{(2 marks)}
    \item Determine whether the Trade Balance (Value) has improved or worsened. \textbf{(2 marks)}
    \item Explain how the \textbf{Prebisch-Singer Hypothesis} relates to Cameroon's long-term TOT prospects. \textbf{(4 marks)}
\end{enumerate}

\tcblower
\textbf{Detailed Solution:}

\textbf{1. Terms of Trade 2023}
\[ TOT_{2023} = \frac{I_X}{I_M} \times 100 = \frac{130}{115} \times 100 = \mathbf{113.04} \]

\textbf{2. Interpretation}
TOT has risen from 100 to 113.04. This is an \textbf{improvement (favourable movement)} of approx 13\%.
\textbf{Meaning:} In 2023, Cameroon needs to export \textbf{fewer} units of cocoa/oil to purchase the same quantity of machinery/food. The purchasing power of exports has increased. Real national income (command over goods) has risen, \emph{ceteris paribus}.

\textbf{3. Export Value Index \& Import Value Index}
Value Index = Price Index $\times$ Volume Index / 100.
\begin{itemize}
    \item \textbf{Export Value Index (2023)} = $\frac{130 \times 90}{100} = \mathbf{117}$.
    \item \textbf{Import Value Index (2023)} = $\frac{115 \times 110}{100} = \mathbf{126.5}$.
\end{itemize}

\textbf{4. Trade Balance Change}
\begin{itemize}
    \item Base Year (2020): Export Value = 100, Import Value = 100. Trade Balance = 0 (assumed balanced for index comparison).
    \item 2023: Export Value = 117, Import Value = 126.5.
    \item Trade Deficit = $126.5 - 117 = 9.5$ index points.
\end{itemize}
\textbf{Conclusion:} Despite \textbf{improving TOT}, the \textbf{Trade Balance has worsened (moved into deficit)}. Why? Import volume grew strongly (10\% $\uparrow$) while export volume fell (10\% $\downarrow$). The volume effect dominated the price effect. This highlights the difference between \textbf{TOT (price ratio)} and \textbf{Trade Balance (value gap)}.

\textbf{5. Prebisch-Singer Hypothesis \& Cameroon}
The Prebisch-Singer Hypothesis argues that the \textbf{secular (long-run) terms of trade of primary commodity exporters tend to deteriorate} relative to manufactured goods exporters.
\textbf{Application to Cameroon:}
\begin{itemize}
    \item \textbf{Structure of Exports:} Cameroon relies heavily on primary commodities: Crude Oil ($\sim$40-50\%), Cocoa beans, Timber (logs/sawn wood), Cotton, Aluminium, Bananas. Low income elasticity of demand (YED $< 1$) and low price elasticity of demand (PED $< 1$).
    \item \textbf{Structure of Imports:} Capital goods (machinery, transport equipment), manufactured consumer goods, refined petroleum products. High YED, high technology content.
    \item \textbf{Mechanisms causing decline:}
        \begin{enumerate}
            \item \textbf{Engel's Law / Low YED:} As world income rises, demand for primary food/raw materials rises slower than demand for manufactures/services.
            \item \textbf{Technological Progress / Synthetics:} Synthetic rubber replaces natural rubber; synthetic fibres replace cotton; lab-grown cocoa substitutes emerging. Reduces demand for natural primary products.
            \item \textbf{Market Power Asymmetry:} Primary markets often competitive (price takers); Manufactured markets often oligopolistic (price makers $\rightarrow$ markup pricing, wages rise with productivity). Productivity gains in manufacturing $\rightarrow$ higher wages/prices; Productivity gains in primary $\rightarrow$ lower prices (competitive markets pass gains to consumers).
            \item \textbf{Volatility:} Primary prices volatile; long-term trend downward.
        \end{enumerate}
    \item \textbf{Policy Implication for Cameroon:} Diversification (structural transformation) into agro-processing (cocoa $\rightarrow$ chocolate/butter), petrochemicals, manufacturing, and services (ICT, tourism) is essential to escape the "Primary Commodity Trap" and secure improving TOT.
\end{itemize}
\end{exoresolu}

\begin{methode}[Recording Balance of Payments Entries (Double Entry)]
\textbf{Context:} Lesson 104/105 (TU 51/52). The fundamental accounting identity: \textbf{Current Account + Capital \& Financial Account + Net Errors \& Omissions = 0}. Every transaction has a Credit (+) and Debit (-) entry.
\textbf{Steps:}
\begin{enumerate}
    \item \textbf{Identify the Transaction:} What flows? (Good, Service, Income, Transfer, Asset, Liability).
    \item \textbf{Determine Direction:}
        \begin{itemize}
            \item \textbf{Credit (+):} Inflow of foreign exchange / Reduction in foreign assets / Increase in foreign liabilities. \textit{Exports, Income received, Transfers received, Foreign investment in Cameroon, Loans received, Sale of reserves.}
            \item \textbf{Debit (-):} Outflow of foreign exchange / Increase in foreign assets / Reduction in foreign liabilities. \textit{Imports, Income paid, Transfers sent, Cameroonian investment abroad, Loan repayment, Purchase of reserves.}
        \end{itemize}
    \item \textbf{Classify into Accounts:}
        \begin{itemize}
            \item \textbf{Current Account:} Goods (Visible), Services (Invisible), Primary Income (Investment income, Compensation of employees), Secondary Income (Current Transfers: Remittances, Aid grants).
            \item \textbf{Capital Account:} Capital Transfers (Debt forgiveness, Migrants' transfers), Acquisition/disposal of non-produced non-financial assets.
            \item \textbf{Financial Account:} Direct Investment (FDI >10\%), Portfolio Investment (<10\%), Other Investment (Loans, Currency/Deposits, Trade Credit), Reserve Assets (BEAC reserves: Gold, SDRs, Foreign Currency).
        \end{itemize}
    \item \textbf{Write the Two Entries:} One Credit, One Debit. Sum = 0.
    \item \textbf{Calculate Balances:} CA Balance = Credits - Debits. FA Balance = Net Acquisition of Assets - Net Incurrence of Liabilities (Sign convention: + = Net Outflow/Lending; - = Net Inflow/Borrowing). \textit{Note: GCE often uses older "Capital Account" = Financial Account + Capital Account. Check syllabus convention. Modern BPM6: Current + Capital + Financial = 0.}
\end{enumerate}
\textbf{Golden Rule:} "Exports/Credits = Money Coming In (+). Imports/Debits = Money Going Out (-)."
\end{methode}

\begin{exoresolu}[Balance of Payments Recording: Cameroon Transactions]
\textbf{Statement:} Record the following transactions in Cameroon's Balance of Payments (BPM6 format). Indicate the Account (Current / Capital / Financial), the Category, and whether it is a Credit (+) or Debit (-) entry. Assume all values in Billion FCFA.
\begin{enumerate}
    \item Cameroon exports 50bn FCFA of cocoa beans to the Netherlands. Payment received in Euros deposited in a Douala bank.
    \item Cameroon imports 30bn FCFA of machinery from France. Payment made by transfer from the Douala bank.
    \item Cameroonian workers in Equatorial Guinea send 10bn FCFA remittances home via Mobile Money.
    \item A French company (TotalEnergies) reinvests 20bn FCFA profits in its Cameroonian subsidiary (FDI).
    \item The Cameroonian government receives a 15bn FCFA project grant (budget support) from the EU.
    \item A Cameroonian commercial bank buys 5bn FCFA of US Treasury Bills (Portfolio Investment).
    \item BEAC (Central Bank) sells 25bn FCFA of foreign reserves (USD) to commercial banks to defend the CFA Franc parity.
    \item China forgives 40bn FCFA of Cameroon's bilateral debt.
\end{enumerate}

\tcblower
\textbf{Detailed Solution:}
\textit{Format: Transaction | Account | Category | Credit (+) / Debit (-) | Amount}

\begin{center}
\begin{tabularx}{\linewidth}{|c|X|X|c|c|}\hline
\textbf{Trans.} & \textbf{Account} & \textbf{Category (BPM6)} & \textbf{Entry} & \textbf{Amt (Bn FCFA)} \\ \hline
1 & Current & Goods (Exports) & \textbf{Credit (+)} & 50 \\
  & Financial & Other Investment: Currency \& Deposits (Assets: Bank deposits abroad $\uparrow$? No, deposit in Douala bank = Liability to non-resident? \textbf{Correction:} Payment received in Euros deposited in Douala bank. If the account is held by a non-resident (Dutch importer) in a Cameroonian bank $\rightarrow$ Liability $\uparrow$ (Credit). If held by Cameroonian exporter $\rightarrow$ Asset (Foreign currency) $\uparrow$ (Debit). \textbf{Standard assumption:} Exporter receives payment $\rightarrow$ Claim on non-resident (Asset) increases $\rightarrow$ \textbf{Debit (-) Financial Account: Other Investment / Currency \& Deposits}. \textit{Wait, standard textbook: Export = Credit Current. Counterpart = Increase in Foreign Assets (Debit Financial) OR Decrease in Foreign Liabilities (Debit Financial). Let's use standard: Export Credit Current; Debit Financial (Asset $\uparrow$).}) & \textbf{Debit (-)} & 50 \\ \hline
2 & Current & Goods (Imports) & \textbf{Debit (-)} & 30 \\
  & Financial & Other Investment: Currency \& Deposits (Assets $\downarrow$ / Liability $\downarrow$) & \textbf{Credit (+)} & 30 \\ \hline
3 & Current & Secondary Income (Personal Transfers) & \textbf{Credit (+)} & 10 \\
  & Financial & Other Investment: Currency \& Deposits (Assets $\uparrow$) & \textbf{Debit (-)} & 10 \\ \hline
4 & Current & Primary Income (Investment Income: Reinvested Earnings) & \textbf{Debit (-)} & 20 \\
  & Financial & Direct Investment (Liabilities: Reinvestment of earnings) & \textbf{Credit (+)} & 20 \\ \hline
5 & Current & Secondary Income (Current Transfers: General Govt Grants) & \textbf{Credit (+)} & 15 \\
  & Financial & Other Investment: Currency \& Deposits (Assets $\uparrow$) & \textbf{Debit (-)} & 15 \\ \hline
6 & Financial & Portfolio Investment (Assets: Debt Securities) & \textbf{Debit (-)} & 5 \\
  & Financial & Other Investment: Currency \& Deposits (Assets $\downarrow$) & \textbf{Credit (+)} & 5 \\ \hline
7 & Financial & Reserve Assets (Foreign Exchange) & \textbf{Credit (+)} & 25 \\
  & Financial & Other Investment: Currency \& Deposits (Liabilities $\downarrow$ / Assets $\uparrow$ for banks) & \textbf{Debit (-)} & 25 \\ \hline
8 & Capital & Capital Transfers (Debt Forgiveness) & \textbf{Credit (+)} & 40 \\
  & Financial & Other Investment: Loans (Liabilities $\downarrow$) & \textbf{Debit (-)} & 40 \\ \hline
\end{tabularx}
\end{center}

\textbf{Key Explanations:}
\begin{itemize}
    \item \textbf{T1:} Export = Credit Current. Counterpart: Exporter now holds a claim (Euro deposit) $\rightarrow$ Financial Asset increases $\rightarrow$ Debit Financial (Other Investment: Currency/Deposits).
    \item \textbf{T2:} Import = Debit Current. Payment: Bank deposits (Asset) decrease $\rightarrow$ Credit Financial.
    \item \textbf{T4:} Reinvested earnings are imputed: Income payable to foreign parent (Debit Current Primary Income) matched by increase in FDI Liability (Credit Financial). \textit{Crucial GCE concept.}
    \item \textbf{T7:} BEAC sells reserves $\rightarrow$ Reserve Assets decrease $\rightarrow$ Credit (+) Reserve Assets. Commercial banks gain foreign currency (Asset $\uparrow$ $\rightarrow$ Debit) or Liabilities to BEAC fall (Debit).
    \item \textbf{T8:} Debt forgiveness = Capital Transfer (Credit Capital Account) + Reduction in Loan Liability (Debit Financial Account: Other Investment/Loans).
\end{itemize}
\textbf{Check:} Sum of all Credits = Sum of all Debits = 195. Identity holds.
\end{exoresolu}

\begin{methode}[Correcting a Balance of Payments Deficit: Expenditure Switching vs Reducing]
\textbf{Context:} Lessons 105/106/111 (TU 53/54/60). Policy responses to CA Deficit.
\textbf{Steps for Essay/Analysis:}
\begin{enumerate}
    \item \textbf{Define Deficit:} Current Account Deficit (Imports > Exports + Net Income/Transfers). Financed by Capital/Financial Inflows (Borrowing/FDI) or Running down Reserves.
    \item \textbf{Categorise Policies:}
        \begin{itemize}
            \item \textbf{Expenditure Reducing (Demand Management):} Contractionary Fiscal (Tax $\uparrow$, G $\downarrow$) / Monetary (Interest Rate $\uparrow$). Reduces $Y \rightarrow$ Reduces $M$ (Import function $M = f(Y)$). \textbf{Side effect:} Recession, Unemployment $\uparrow$. \textit{IMF Standard Conditionality.}
            \item \textbf{Expenditure Switching (Relative Prices):} \textbf{Devaluation/Depreciation} (CFA Franc fixed $\rightarrow$ requires BEAC/France agreement; difficult). Makes Exports cheaper (X $\uparrow$), Imports dearer (M $\downarrow$). \textbf{Condition:} Marshall-Lerner Condition ($PED_X + PED_M > 1$). \textbf{J-Curve:} Short run worsening (volumes fixed, values move against) $\rightarrow$ Long run improvement.
            \item \textbf{Supply-Side / Structural:} Improve competitiveness (Productivity $\uparrow$, Quality $\uparrow$, Diversification). Long term.
            \item \textbf{Direct Controls (Protectionism):} Tariffs, Quotas, Import Licensing. \textbf{Violates WTO/CEMAC rules.} Retaliation risk. Welfare loss (see Tariff method).
        \end{itemize}
    \item \textbf{Evaluate for Cameroon (CEMAC Context):}
        \begin{itemize}
            \item \textbf{Fixed Exchange Rate (CFA Franc):} Cannot devalue unilaterally. Requires regional consensus (BEAC) + French Treasury guarantee. "Internal Devaluation" (Wage cuts, austerity) is the only expenditure-switching tool nationally.
            \item \textbf{Import Structure:} Essential imports (Fuel, Medicine, Capital goods, Food) $\rightarrow$ Low PED$_M$. Devaluation $\rightarrow$ Import bill $\uparrow$ $\rightarrow$ Inflation $\uparrow$ (Cost-push) $\rightarrow$ Real wage $\downarrow$ $\rightarrow$ Social unrest.
            \item \textbf{Export Structure:} Primary commodities $\rightarrow$ Low PED$_X$ (Price takers). Devaluation $\rightarrow$ Volume $\uparrow$ little, Revenue $\uparrow$ in FCFA but terms of trade may worsen.
            \item \textbf{Conclusion:} Expenditure Reducing (Fiscal consolidation) + Structural Diversification (AfCFTA, Industrialisation) are the viable sustainable paths. Devaluation is a "nuclear option" for the region.
        \end{itemize}
\end{enumerate}
\end{methode}

\begin{exoresolu}[Policy Evaluation: Correcting Cameroon's Current Account Deficit]
\textbf{Statement:} "Cameroon has experienced persistent Current Account deficits. Evaluate the effectiveness of \textbf{Expenditure Reducing} policies versus \textbf{Expenditure Switching} (Devaluation) policies in correcting this deficit, within the context of the CEMAC zone and the CFA Franc regime." \textbf{(20 marks essay style)}

\tcblower
\textbf{Detailed Solution (Essay Plan \& Key Arguments):}

\textbf{Introduction}
\begin{itemize}
    \item Define Current Account Deficit (CAD): $X + NY + NCT < M$. Cameroon's CAD driven by structural imports (fuel, capital goods, food) vs volatile primary exports (oil, cocoa).
    \item Define Policies: Expenditure Reducing (Fiscal/Monetary contraction to lower $Y \rightarrow M$). Expenditure Switching (Exchange rate depreciation to raise $P_M/P_X \rightarrow$ Switch demand).
    \item Thesis: In the CEMAC fixed exchange rate regime, Expenditure Reducing is the \textbf{only immediately available} macro tool, but it carries high social costs. Expenditure Switching (Devaluation) is \textbf{constrained institutionally} and \textbf{structurally ineffective} due to low elasticities. Sustainable correction requires \textbf{Supply-Side Structural Reforms}.
\end{itemize}

\textbf{Body 1: Expenditure Reducing Policies (Fiscal Consolidation / IMF Programs)}
\begin{itemize}
    \item \textbf{Mechanism:} $\uparrow$ Taxes (VAT, Corporate), $\downarrow$ Govt Spending (Subsidies, Wage bill) $\rightarrow$ $\downarrow$ Aggregate Demand (AD) $\rightarrow$ $\downarrow$ National Income (Y) $\rightarrow$ $\downarrow$ Import Demand ($M = mY$).
    \item \textbf{Effectiveness:} Directly targets the "Absorption" side ($Y > Absorption$). Works if deficit is cyclical/demand-driven.
    \item \textbf{Limitations (Cameroon Context):}
        \begin{itemize}
            \item \textbf{Recession/Unemployment:} Cameroon needs growth ($\sim$4-5\%) for job creation (youth bulge). Austerity $\rightarrow$ Social tension (e.g., 2008 riots, Anglophone crisis exacerbation).
            \item \textbf{Low Tax Base:} Large informal sector ($\sim$90\% employment). Hard to raise revenue without widening base.
            \item \textbf{Import Inelasticity:} Essential imports (Fuel $\sim$20\% import bill, Medicines, Capital goods for investment) don't fall much with Y. Deficit persists.
            \item \textbf{Fiscal Multiplier:} High in developing economies $\rightarrow$ Large output loss per unit of deficit reduction.
        \end{itemize}
\end{itemize}

\textbf{Body 2: Expenditure Switching: Devaluation of CFA Franc}
\begin{itemize}
    \item \textbf{Mechanism:} $E \uparrow$ (FCFA per EUR/USD) $\rightarrow$ $P_X^{\text{foreign}} \downarrow$, $P_M^{\text{domestic}} \uparrow$ $\rightarrow$ $X \uparrow$, $M \downarrow$.
    \item \textbf{Institutional Constraint:} CFA Franc fixed to Euro (1 EUR = 655.957 FCFA). Guaranteed by French Treasury. Devaluation requires \textbf{regional consensus (BEAC/CEMAC Heads of State)} + \textbf{French approval}. Politically extremely difficult (1994 devaluation was unique, traumatic).
    \item \textbf{Structural Ineffectiveness (Elasticities):}
        \begin{itemize}
            \item \textbf{Marshall-Lerner Condition:} $PED_X + PED_M > 1$?
            \item \textbf{Exports (Oil, Cocoa, Wood):} Price Takers. $PED_X \approx 0$ (World demand insensitive to Cameroon's price). Devaluation $\rightarrow$ FCFA revenue $\uparrow$ per barrel, but Volume $\approx$ constant. \textbf{Terms of Trade may worsen} if priced in USD/EUR.
            \item \textbf{Imports (Fuel, Food, Capital):} Necessities. $PED_M < 1$. Devaluation $\rightarrow$ Import Value ($P_M \times Q_M$) $\uparrow$ $\rightarrow$ \textbf{Worsens CA initially (J-Curve)} and fuels \textbf{Imported Inflation}.
            \item \textbf{Pass-through:} High in CEMAC (shallow markets, dollarized invoicing). Inflation $\uparrow$ $\rightarrow$ Real Exchange Rate appreciates back $\rightarrow$ Competitiveness gain eroded.
        \end{itemize}
    \item \textbf{1994 Precedent:} 50\% devaluation (100 $\rightarrow$ 50 FCFA/FRF). Short term: Import compression, Export revenue boost (FCFA). Medium term: Inflation $\uparrow$, Real wages $\downarrow$, Poverty $\uparrow$, Competitiveness gains lost by 1998. Not repeatable easily.
\end{itemize}

\textbf{Body 3: The Sustainable Alternative — Supply Side \& Structural Transformation}
\begin{itemize}
    \item \textbf{Diversification:} Move up value chain. Cocoa $\rightarrow$ Cocoa Butter/Powder/Chocolate (Local processing: SIC Cacaos, Chococam). Wood $\rightarrow$ Furniture. Cotton $\rightarrow$ Textiles. Oil $\rightarrow$ Petrochemicals (SONARA upgrade).
    \item \textbf{AfCFTA:} Access to 1.3B market. Reduce intra-African tariffs $\rightarrow$ Export manufactured goods to Nigeria, Chad, Gabon (higher PED, higher value added).
    \item \textbf{Business Climate:} Reduce red tape, electricity costs (ENEO reliability), port efficiency (Douala/Kribi), access to finance (SMEs).
    \item \textbf{Import Substitution (Strategic):} Local rice, maize, poultry, pharmaceuticals (Cinpharm). Reduces structural import dependency.
\end{itemize}

\textbf{Conclusion}
Expenditure Reducing is the \textbf{default adjustment mechanism} under the fixed peg (Internal Devaluation), but it is socially painful and treats symptoms, not causes. Expenditure Switching (Nominal Devaluation) is \textbf{institutionally blocked} and \textbf{economically blunt} due to perverse elasticities (Primary exports, Essential imports). The \textbf{only durable solution} is \textbf{Supply-Side Structural Transformation}: Diversifying exports (raising $PED_X$), substituting strategic imports (lowering $PED_M$), and leveraging AfCFTA. This shifts the CA constraint outward, allowing growth \textbf{with} external balance.
\end{exoresolu}

\begin{methode}[Foreign Exchange Market Equilibrium \& Exchange Rate Determination]
\textbf{Context:} Lessons 109/110 (TU 58/59). Supply and Demand for Foreign Exchange (Forex).
\textbf{Steps:}
\begin{enumerate}
    \item \textbf{Axes:} Exchange Rate ($E$ = Domestic Currency / Foreign Currency, e.g., FCFA/USD) on Y-axis. Quantity of Foreign Currency (Q\$) on X-axis. \textit{Note: GCE sometimes uses $E$ = Foreign/Domestic. Stick to \textbf{Direct Quote (FCFA per \$)} for Cameroon context.}
    \item \textbf{Demand Curve ($D_\$$):} \textbf{Downward Sloping.} Derived from Demand for Imports ($M$), Capital Outflow, Speculation. $E \downarrow$ (FCFA appreciates) $\rightarrow$ Foreign goods cheaper $\rightarrow$ $Q_\$ \uparrow$.
    \item \textbf{Supply Curve ($S_\$$):} \textbf{Upward Sloping.} Derived from Export Earnings ($X$), Capital Inflow (FDI, Aid, Loans), Remittances, Speculation. $E \uparrow$ (FCFA depreciates) $\rightarrow$ Cameroonian goods cheaper for foreigners $\rightarrow$ $Q_\$ \uparrow$.
    \item \textbf{Equilibrium:} $D_\$ = S_\$$ $\rightarrow$ $E^*$.
    \item \textbf{Shifts (Causes of Appreciation/Depreciation):}
        \begin{itemize}
            \item $D_\$ \uparrow$ (Right): Import boom, Capital flight, Speculative attack $\rightarrow$ \textbf{Depreciation ($E \uparrow$)}.
            \item $D_\$ \downarrow$ (Left): Import compression, Capital repatriation $\rightarrow$ \textbf{Appreciation ($E \downarrow$)}.
            \item $S_\$ \uparrow$ (Right): Export boom (Oil price $\uparrow$), FDI inflow, Aid/Remittances $\uparrow$ $\rightarrow$ \textbf{Appreciation ($E \downarrow$)}.
            \item $S_\$ \downarrow$ (Left): Export collapse, Capital flight (residents buying foreign assets) $\rightarrow$ \textbf{Depreciation ($E \uparrow$)}.
        \end{itemize}
    \item \textbf{Fixed Rate (CFA Franc):} $E$ fixed at $E_{\text{peg}}$. BEAC intervenes.
        \begin{itemize}
            \item Excess Demand ($D_\$ > S_\$$ at $E_{\text{peg}}$): BEAC sells Reserves ($\downarrow$ Reserves) to meet gap.
            \item Excess Supply ($S_\$ > D_\$$): BEAC buys Reserves ($\uparrow$ Reserves).
        \end{itemize}
\end{enumerate}
\end{methode}

\begin{exoresolu}[Forex Market Dynamics: Oil Price Shock in CEMAC]
\textbf{Statement:} Cameroon is a net oil exporter. The world price of crude oil falls sharply by 40\%. Using a Supply and Demand diagram for the Foreign Exchange Market (FCFA per USD), analyse the impact on:
\begin{enumerate}
    \item The Supply Curve of USD ($S_\$$). \textbf{(3 marks)}
    \item The Equilibrium Exchange Rate ($E^*$) under a \textbf{Floating Rate} regime. \textbf{(3 marks)}
    \item BEAC's Foreign Reserves and Money Supply under the \textbf{Fixed Rate (CFA Franc)} regime. \textbf{(4 marks)}
    \item Explain the concept of "Dutch Disease" and how an oil boom (reverse shock) could harm the non-oil tradable sector (e.g., Cocoa, Textiles). \textbf{(4 marks)}
\end{enumerate}

\tcblower
\textbf{Detailed Solution:}

\textbf{1. Impact on Supply Curve ($S_\$$)}
\begin{itemize}
    \item Oil exports are a major component of \textbf{Supply of Foreign Exchange} (Export Earnings $\rightarrow$ Supply of USD/EUR).
    \item Oil Price $\downarrow$ 40\% $\rightarrow$ Value of Oil Exports ($P_{oil} \times Q_{oil}$) $\downarrow$ significantly (assuming low short-run PED).
    \item Fewer USD earned per barrel $\rightarrow$ Total Supply of USD shifts \textbf{LEFT (Inward)} at every exchange rate.
    \item \textbf{Diagram:} $S_{\$1}$ shifts to $S_{\$2}$ (leftwards).
\end{itemize}

\textbf{2. Floating Rate Equilibrium}
\begin{itemize}
    \item Initial Equilibrium: $E_1$ (intersection $D_\$$, $S_{\$1}$).
    \item New Equilibrium: Intersection $D_\$$, $S_{\$2}$.
    \item $S_\$$ shifts Left $\rightarrow$ Excess Demand for USD at $E_1$.
    \item Exchange Rate $E$ (FCFA/USD) \textbf{RISES} to $E_2$.
    \item \textbf{Result:} \textbf{Depreciation of the FCFA}. It takes more FCFA to buy 1 USD.
    \item Quantity of USD traded falls ($Q_2 < Q_1$).
\end{itemize}

\textbf{3. Fixed Rate (CFA Franc) Regime — BEAC Intervention}
\begin{itemize}
    \item Rate fixed at $E_{\text{peg}}$ (e.g., 600 FCFA/USD).
    \item At $E_{\text{peg}}$, $S_{\$2} < D_\$$ $\rightarrow$ \textbf{Excess Demand for USD (Shortage)}.
    \item BEAC \textbf{must} supply the gap ($Q_D - Q_S$) from its \textbf{Foreign Reserves}.
    \item \textbf{Reserves:} \textbf{FALL} (Credit entry in Financial Account: Reserve Assets $\downarrow$).
    \item \textbf{Money Supply (Monetary Base):} BEAC sells USD, receives FCFA from commercial banks. FCFA is withdrawn from circulation. \textbf{Money Supply (M0/M2) CONTRACTS}.
    \item \textbf{Transmission:} $\downarrow$ Reserves $\rightarrow$ $\downarrow$ Monetary Base $\rightarrow$ $\uparrow$ Interest Rates (Interbank) $\rightarrow$ $\downarrow$ Credit $\rightarrow$ $\downarrow$ Aggregate Demand $\rightarrow$ $\downarrow$ Imports (Expenditure Reducing). \textit{This is the automatic adjustment mechanism of a Currency Board / Fixed Peg.}
    \item \textbf{Risk:} If reserves fall below statutory floor (20\% of sight liabilities), credibility of peg threatened $\rightarrow$ Speculative Attack $\rightarrow$ Forced Devaluation.
\end{itemize}

\textbf{4. Dutch Disease (The Reverse Shock: Oil Boom)}
\begin{itemize}
    \item \textbf{Definition:} A resource boom (Oil price $\uparrow$) causes Real Exchange Rate Appreciation, damaging the competitiveness of the \textbf{Non-Resource Tradable Sector} (Manufacturing, Agriculture/Cocoa).
    \item \textbf{Mechanism (Two Channels):}
        \begin{enumerate}
            \item \textbf{Spending Effect (Demand Side):} Oil Revenue $\uparrow$ $\rightarrow$ Govt/Private Income $\uparrow$ $\rightarrow$ Demand for \textbf{Non-Tradables} (Construction, Services, Retail) $\uparrow$ $\rightarrow$ Prices of Non-Tradables ($P_N$) $\uparrow$ relative to Tradables ($P_T$ fixed by world market). $\rightarrow$ \textbf{Real Appreciation ($P_N/P_T \uparrow$)}. Resources (Labour, Capital) pulled from Cocoa/Textiles $\rightarrow$ Non-Tradables.
            \item \textbf{Resource Movement Effect (Supply Side):} High wages in Oil sector $\rightarrow$ Labour/Capital moves from Cocoa/Textiles $\rightarrow$ Oil/Services. Output of Non-Oil Tradables $\downarrow$.
        \end{enumerate}
    \item \textbf{Cameroon Context:} 1970s-80s Oil Boom $\rightarrow$ Real appreciation $\rightarrow$ Neglect of Agriculture (Cocoa/Coffee output stagnated), Manufacturing base remained weak. "Resource Curse".
    \item \textbf{Policy:} Stabilisation Fund (save oil revenue abroad), Invest in Infrastructure/Human Capital (raise productivity in Tradables), Avoid excessive non-tradable spending.
\end{itemize}
\end{exoresolu}

\begin{methode}[Exchange Rate Adjustments: Marshall-Lerner Condition \& J-Curve]
\textbf{Context:} Lesson 111 (TU 60). The condition for a devaluation/depreciation to improve the Current Account (CA).
\textbf{Steps:}
\begin{enumerate}
    \item \textbf{Define Elasticities:} $PED_X$ (Price Elasticity of Demand for Exports), $PED_M$ (Price Elasticity of Demand for Imports). Both defined as \textbf{absolute values} (positive numbers).
    \item \textbf{Marshall-Lerner Condition (MLC):} $\mathbf{PED_X + PED_M > 1}$.
        \begin{itemize}
            \item If True: Depreciation $\rightarrow$ CA Improves (Value of X $\uparrow$ more than Value of M $\uparrow$, or X $\uparrow$ + M $\downarrow$).
            \item If False ($<1$): Depreciation $\rightarrow$ CA Worsens (Inelastic: Price effect dominates Volume effect).
        \end{itemize}
    \item \textbf{Derivation (Intuition):}
        $CA = P_X Q_X - P_M Q_M$ (in domestic currency).
        Depreciation: $P_X^{\text{for}} \downarrow$, $P_M^{\text{dom}} \uparrow$.
        $\Delta CA \approx \text{Volume Effects} - \text{Price Effects}$.
        Requires volume response strong enough to offset adverse price movement on imports.
    \item \textbf{J-Curve Effect (Time Path):}
        \begin{enumerate}
            \item \textbf{Short Run (t0-t1):} Contracts fixed, recognition lags, delivery lags $\rightarrow$ $PED_X \approx 0, PED_M \approx 0$. MLC fails. Import Bill $\uparrow$ (Price $\uparrow$, Q same), Export Revenue $\downarrow$ or flat (Foreign price $\downarrow$, Q same). \textbf{CA Worsens.}
            \item \textbf{Medium Run (t1-t2):} Volumes adjust. $PED$ rises. MLC holds. CA improves.
            \item \textbf{Long Run (t2+):} Full adjustment. New equilibrium.
        \end{enumerate}
    \item \textbf{Cameroon/CEMAC Specifics:} Low $PED_X$ (Primary commodities, price takers), Low $PED_M$ (Essentials: Fuel, Medicine, Capital goods). \textbf{MLC likely FAILS or is very weak.} Devaluation risky.
\end{enumerate}
\end{methode}

\begin{exoresolu}[Marshall-Lerner \& J-Curve: Numerical Application]
\textbf{Statement:} Cameroon's Current Account is initially balanced ($X = M = 1000$ bn FCFA). The CFA Franc depreciates by 20\% against the USD (major invoicing currency).
Elasticities (Absolute Values):
Short Run: $PED_X = 0.2$, $PED_M = 0.1$.
Long Run: $PED_X = 0.8$, $PED_M = 0.6$.
Assume export prices in foreign currency adjust fully (Pass-through = 1), import prices in domestic currency adjust fully.
\begin{enumerate}
    \item State the Marshall-Lerner Condition. Does it hold in the Short Run? Long Run? \textbf{(3 marks)}
    \item Calculate the approximate percentage change in Export Value ($X$) and Import Value ($M$) in the \textbf{Short Run}. Determine the new CA Balance. \textbf{(4 marks)}
    \item Calculate the approximate percentage change in Export Value ($X$) and Import Value ($M$) in the \textbf{Long Run}. Determine the new CA Balance. \textbf{(4 marks)}
    \item Sketch the J-Curve diagram labelling the phases. \textbf{(3 marks)}
\end{enumerate}

\tcblower
\textbf{Detailed Solution:}

\textbf{1. Marshall-Lerner Condition}
Condition: $PED_X + PED_M > 1$.
\begin{itemize}
    \item \textbf{Short Run:} $0.2 + 0.1 = 0.3 < 1$. \textbf{Condition FAILS.} Depreciation worsens CA.
    \item \textbf{Long Run:} $0.8 + 0.6 = 1.4 > 1$. \textbf{Condition HOLDS.} Depreciation improves CA.
\end{itemize}

\textbf{2. Short Run Impact (Approximation Formula: $\% \Delta Value \approx \% \Delta P + \% \Delta Q$)}
\textit{Depreciation 20\% $\rightarrow$ $P_X^{\text{FCFA}} \uparrow 20\%$ (Foreign price constant), $P_M^{\text{FCFA}} \uparrow 20\%$.}
\textit{Volume change: $\% \Delta Q_X = -PED_X \times \% \Delta P_X^{\text{for}}$. Foreign price falls 20\%? Wait.}
\textit{Standard Setup: Domestic Currency Depreciates 20\% ($E \uparrow 20\%$).}
\begin{itemize}
    \item \textbf{Exports:} Foreign Price $P_X^* = P_X / E$. If $P_X$ (FCFA) sticky, $P_X^* \downarrow 20\%$. $\% \Delta Q_X = PED_X \times 20\% = 0.2 \times 20\% = +4\%$.
    Export Value (FCFA) $V_X = P_X \times Q_X$. $\% \Delta V_X \approx \% \Delta P_X + \% \Delta Q_X$.
    If local currency price $P_X$ unchanged (sticky): $\% \Delta V_X \approx 0\% + 4\% = \mathbf{+4\%}$.
    (New $X = 1040$).
    \item \textbf{Imports:} Domestic Price $P_M = P_M^* \times E$. Foreign price $P_M^*$ constant. $P_M \uparrow 20\%$.
    $\% \Delta Q_M = -PED_M \times 20\% = -0.1 \times 20\% = -2\%$.
    Import Value (FCFA) $V_M = P_M \times Q_M$. $\% \Delta V_M \approx +20\% - 2\% = \mathbf{+18\%}$.
    (New $M = 1180$).
\end{itemize}
\textbf{Short Run CA Balance:} $X - M = 1040 - 1180 = \mathbf{-140 \text{ bn FCFA (Deficit)}}$.
\textit{CA Worsens from 0 to -140.}

\textbf{3. Long Run Impact}
\begin{itemize}
    \item \textbf{Exports:} $\% \Delta Q_X = 0.8 \times 20\% = +16\%$. $\% \Delta V_X \approx 0\% + 16\% = \mathbf{+16\%}$. (New $X = 1160$).
    \item \textbf{Imports:} $\% \Delta Q_M = -0.6 \times 20\% = -12\%$. $\% \Delta V_M \approx +20\% - 12\% = \mathbf{+8\%}$. (New $M = 1080$).
\end{itemize}
\textbf{Long Run CA Balance:} $X - M = 1160 - 1080 = \mathbf{+80 \text{ bn FCFA (Surplus)}}$.
\textit{CA Improves from -140 to +80.}

\textbf{4. J-Curve Diagram Description}
\begin{popfigure}
\begin{tikzpicture}[scale=0.8]
\begin{axis}[
    width=12cm, height=7cm,
    xlabel={Time}, ylabel={Current Account Balance},
    xmin=0, xmax=5, ymin=-200, ymax=150,
    xtick={1,2,3,4}, xticklabels={$t_0$ (Deprec.), $t_1$, $t_2$, $t_3$},
    ytick={-140,0,80}, yticklabels={-140, 0, +80},
    axis lines=middle,
    clip=false
]
% J-Curve path
\addplot[popblue, very thick, smooth, samples=50, domain=0:4] 
    ({x}, {80 - 220*exp(-x) + 80*exp(-2*x)}) ; % Rough J shape
% Horizontal line at 0
\addplot[dashed, gray] coordinates {(0,0) (4,0)};
% Annotations
\node[popblue, anchor=south west] at (axis cs:0.5, -100) {\textbf{Phase 1: Deterioration}};
\node[popblue, anchor=south west] at (axis cs:2.5, 50) {\textbf{Phase 2: Improvement}};
\node[popgreen, anchor=south] at (axis cs:4, 80) {\textbf{Long Run Surplus}};
\node[poporange, anchor=north] at (axis cs:1, -140) {\textbf{Trough}};
\end{axis}
\end{tikzpicture}
\legende{The J-Curve Effect: Initial deterioration due to low elasticities (Price effect $>$ Volume effect), followed by improvement as volumes adjust (MLC satisfied).}
\end{popfigure}
\textbf{Key Labels for Exam:} Time (X-axis), CA Balance (Y-axis). Initial equilibrium (0). Depreciation at $t_0$. Immediate drop (Inelastic). Trough. Crossing zero (MLC holds). Long run surplus.
\end{exoresolu}

\begin{methode}[Analysing Regional Integration: CEMAC / AfCFTA Effects]
\textbf{Context:} Lesson 108 (TU 57). Trade Creation vs Trade Diversion (Viner).
\textbf{Steps:}
\begin{enumerate}
    \item \textbf{Define the Bloc:} Identify members (CEMAC: Cameroon, CAR, Chad, Congo, Eq. Guinea, Gabon). Common External Tariff (CET).
    \item \textbf{Trade Creation (Welfare Gain):} Removal of tariffs \textbf{within} bloc $\rightarrow$ High cost domestic producer replaced by Low cost partner producer.
        \begin{itemize}
            \item Condition: Partner's Cost $<$ Domestic Cost $<$ Partner's Cost + Tariff (pre-union).
            \item Effect: Consumption shifts to efficient partner $\rightarrow$ Resource allocation improves $\rightarrow$ Welfare $\uparrow$.
        \end{itemize}
    \item \textbf{Trade Diversion (Welfare Loss):} Removal of tariffs \textbf{within} bloc + CET \textbf{against} Rest of World (RoW) $\rightarrow$ Low cost RoW producer replaced by Higher cost Partner producer.
        \begin{itemize}
            \item Condition: RoW Cost $<$ Partner Cost $<$ RoW Cost + CET.
            \item Effect: Imports switch from efficient RoW to inefficient Partner $\rightarrow$ Welfare $\downarrow$ (Tariff revenue lost, higher resource cost).
        \end{itemize}
    \item \textbf{Net Welfare Effect:} Creation $-$ Diversion. Dynamic effects (Economies of scale, Competition, Investment) usually positive long run.
    \item \textbf{AfCFTA Specifics:} Phased tariff liberalisation (90\% lines $\rightarrow$ 0\%, 7\% sensitive, 3\% excluded). Rules of Origin critical (prevent transshipment). Non-Tariff Barriers (NTBs) removal key (corridors, standards).
    \item \textbf{Cameroon's Position:} Largest CEMAC economy. Hub for transit (Chad, CAR). Potential manufacturing hub for AfCFTA. Risks: Deindustrialisation if flooded by cheaper imports (Nigeria, Egypt, SA), Revenue loss from tariffs (fiscal gap).
\end{enumerate}
\end{methode}

\begin{exoresolu}[Trade Creation vs Diversion: Cameroon in CEMAC \& AfCFTA]
\textbf{Statement:} Cameroon produces wheat flour domestically at a cost of \SI{250}{FCFA/kg}. It imports from two sources:
\begin{itemize}
    \item \textbf{France (RoW):} Cost \SI{200}{FCFA/kg} (Efficient).
    \item \textbf{Chad (CEMAC Partner):} Cost \SI{220}{FCFA/kg} (Less efficient than France, more efficient than Cameroon).
\end{itemize}
CEMAC Common External Tariff (CET) on flour = 20\%. Intra-CEMAC tariff = 0\%.
\begin{enumerate}
    \item Calculate the domestic price of flour from each source \textbf{before} CEMAC integration (assuming MFN tariff 20\% on all). \textbf{(2 marks)}
    \item Determine the source of imports before integration. \textbf{(1 mark)}
    \item Calculate the domestic price of flour from each source \textbf{after} CEMAC integration (CET on France, 0\% on Chad). \textbf{(2 marks)}
    \item Determine the new source of imports. Identify if this is Trade Creation or Trade Diversion. Calculate the welfare change (Net gain/loss per kg). \textbf{(5 marks)}
    \item Discuss \textbf{one} dynamic benefit and \textbf{one} risk for Cameroon from the AfCFTA. \textbf{(4 marks)}
\end{enumerate}

\tcblower
\textbf{Detailed Solution:}

\textbf{1. Prices Before Integration (MFN 20\% on all)}
\begin{itemize}
    \item Domestic Production Cost: \textbf{250 FCFA/kg}.
    \item France (RoW): Landed Cost = $200 \times 1.20 = \mathbf{240 \text{ FCFA/kg}}$.
    \item Chad (Partner): Landed Cost = $220 \times 1.20 = \mathbf{264 \text{ FCFA/kg}}$.
\end{itemize}

\textbf{2. Source Before Integration}
Cheapest price = \textbf{France at 240 FCFA/kg}. (Domestic 250, Chad 264).
Cameroon imports from France. Tariff Revenue = $40 \text{ FCFA/kg} \times Q$.

\textbf{3. Prices After CEMAC Integration}
\begin{itemize}
    \item Domestic: \textbf{250 FCFA/kg}.
    \item France (RoW): CET 20\% applies. Price = $200 \times 1.20 = \mathbf{240 \text{ FCFA/kg}}$.
    \item Chad (Partner): Intra-CEMAC Tariff 0\%. Price = $220 \times 1.00 = \mathbf{220 \text{ FCFA/kg}}$.
\end{itemize}

\textbf{4. New Source \& Welfare Analysis}
Cheapest price = \textbf{Chad at 220 FCFA/kg}.
\textbf{Import Source Switches:} From France (Efficient, 200) $\rightarrow$ Chad (Less efficient, 220).
\textbf{Classification:} \textbf{TRADE DIVERSION.}
\textit{Why?} Partner cost (220) $>$ RoW cost (200). The union diverts trade from a more efficient non-member to a less efficient member due to the tariff preference.

\textbf{Welfare Calculation (per kg):}
\begin{itemize}
    \item \textbf{Consumer Gain:} Price falls from 240 (France+Tariff) to 220 (Chad). Gain = \textbf{+20 FCFA/kg}.
    \item \textbf{Govt Revenue Loss:} Tariff on France (40 FCFA/kg) is lost. Loss = \textbf{-40 FCFA/kg}.
    \item \textbf{Resource Cost Change:} Real resource cost of supply rises from 200 (France) to 220 (Chad). Loss = \textbf{-20 FCFA/kg}.
    \item \textbf{Net Welfare Change:} Consumer Gain (+20) + Govt Loss (-40) = \textbf{-20 FCFA/kg}. (Equals Resource Cost Loss).
\end{itemize}
\textbf{Result:} \textbf{Net Welfare Loss of 20 FCFA/kg.} Trade Diversion dominates.

\textbf{5. AfCFTA: Dynamic Benefit \& Risk for Cameroon}
\begin{itemize}
    \item \textbf{Dynamic Benefit: Economies of Scale \& Industrialisation.}
    Cameroon's domestic market ($\sim$28M) is small for efficient manufacturing (cement, textiles, agro-processing, pharmaceuticals). AfCFTA (1.3B consumers) allows Cameroonian firms to reach \textbf{Minimum Efficient Scale (MES)}, lowering LRAC. Example: Cimencam/Dangote Cement exporting to Chad/CAR/Nigeria; Chococam processing cocoa for West Africa. Attracts FDI targeting regional market (not just national). "Market size $\rightarrow$ Investment $\rightarrow$ Productivity $\rightarrow$ Growth".
    \item \textbf{Risk: Deindustrialisation / Fiscal Shock.}
    Cameroon's industrial base is fragile. AfCFTA eliminates tariffs on 90\% of lines. Cheaper, more efficient imports from Nigeria (cement, textiles, processed food), Egypt (pharma, ceramics), South Africa (machinery, vehicles) could flood the Douala market. Local SMEs (uncompetitive due to high energy costs, poor infrastructure, informality) may collapse $\rightarrow$ Job losses. Simultaneously, \textbf{Customs Revenue} (major part of tax base, $\sim$30-40\% of revenue) declines sharply $\rightarrow$ Fiscal deficit pressure $\rightarrow$ Cuts in public investment (roads, energy) $\rightarrow$ Further hurts competitiveness. \textbf{Mitigation:} AfCFTA Adjustment Fund, sensitive/exclusion lists (7\%+3\%), aggressive supply-side reforms (energy, logistics, skills).
\end{itemize}
\end{exoresolu}

\begin{aretenir}[Summary of Key Methods for Revision]
\begin{itemize}
    \item \textbf{Comparative Advantage:} Always use \textbf{Opportunity Cost}. Lower OC = CA. TOT must lie between domestic OCs.
    \item \textbf{Tariff Diagram:} $S_w$ horizontal $\rightarrow$ Shift up by $t$. Identify 5 areas: CS Loss, PS Gain, Govt Rev, 2x DWL Triangles.
    \item \textbf{Terms of Trade:} $TOT = \frac{I_X}{I_M} \times 100$. Rising = Good (usually). Distinguish from Trade Balance (Value = Price $\times$ Volume).
    \item \textbf{BoP Recording:} Double Entry. Credit = Inflow/FX earning. Debit = Outflow/FX spending. Current vs Capital vs Financial vs Reserves.
    \item \textbf{Deficit Correction:} Expenditure Reducing (Austerity) vs Switching (Devaluation). \textbf{CEMAC Constraint:} Fixed Peg $\rightarrow$ Internal Devaluation only. MLC usually fails for primary exporters.
    \item \textbf{Forex Market:} $D_\$ \downarrow$, $S_\$ \uparrow$. Fixed Rate $\rightarrow$ Reserves adjust, Money Supply adjusts automatically.
    \item \textbf{J-Curve / MLC:} Short run Inelastic $\rightarrow$ Worsening. Long run Elastic $\rightarrow$ Improvement. Condition: $PED_X + PED_M > 1$.
    \item \textbf{Regional Integration:} Trade Creation (Good: Low cost partner replaces High cost domestic) vs Trade Diversion (Bad: High cost partner replaces Low cost RoW). AfCFTA = Scale benefits vs Fiscal/Industrial risks.
\end{itemize}
\end{aretenir}

\newpage

\coursec{Exercises}

\courssub{I check my knowledge}

\begin{exercice}[MCQ: Basis of Trade]
Which of the following statements best describes the principle of comparative advantage?
\begin{enumerate}
\item A country should produce only goods for which it has an absolute advantage.
\item A country should specialise in goods for which its opportunity cost is lower than that of other countries.
\item Trade is beneficial only if both countries have identical production possibilities.
\item Comparative advantage requires that a country has fewer natural resources than its trading partner.
\end{enumerate}
\end{exercice}

\begin{exercice}[True/False with Justification]
State whether each statement is true or false and give a brief justification.
\begin{enumerate}
\item A tariff always increases the domestic price of the imported good by the full amount of the tariff.
\item The current account of the balance of payments records only trade in goods.
\item An improvement in the terms of trade means that a country can buy more imports for a given quantity of exports.
\item Under a fixed exchange rate regime, the central bank must hold foreign exchange reserves to maintain the peg.
\end{enumerate}
\end{exercice}

\begin{exercice}[Definitions]
Define the following terms precisely:
\begin{enumerate}
\item Terms of Trade (TOT)
\item Balance of Payments (BOP) overall balance
\item Devaluation (as distinct from depreciation)
\item Infant industry argument for protection
\end{enumerate}
\end{exercice}

\begin{exercice}[Direct Calculations]
Cameroon exports cocoa beans and imports machinery. In 2023 the export price index for cocoa was 120 (base 2015=100) and the import price index for machinery was 150 (base 2015=100).
\begin{enumerate}
\item Compute the Terms of Trade index for 2023.
\item If the 2022 TOT index was 95, calculate the percentage change in the terms of trade between 2022 and 2023.
\item Interpret the result for Cameroon's purchasing power of imports.
\end{enumerate}
\end{exercice}

\courssub{I practise}

\begin{exercice}[Absolute and Comparative Advantage: Cameroon and Guinea]
The table shows output per worker per day in Cameroon and Guinea for two goods, cocoa and aluminium.

\begin{center}
\begin{tabular}{|l|c|c|}
\hline
 & Cocoa (tonnes) & Aluminium (tonnes) \\
\hline
Cameroon & 4 & 2 \\
Guinea & 1 & 3 \\
\hline
\end{tabular}
\end{center}

\begin{enumerate}
\item Determine which country has an absolute advantage in each good.
\item Calculate the opportunity cost of producing one tonne of cocoa in each country.
\item Identify the comparative advantage of each country.
\item Propose a mutually beneficial terms of trade (exchange ratio) for cocoa and aluminium and show the gains from specialisation and trade for both countries.
\end{enumerate}
\end{exercice}

\begin{exercice}[Effects of a Tariff]
Cameroon imposes a specific tariff of 500 FCFA per tonne on imported wheat. The world price is 10,000 FCFA per tonne. Domestic demand and supply equations (in tonnes) are:
\[
Q_d = 5000 - 0.2P,\quad Q_s = 1000 + 0.1P
\]
where \(P\) is the domestic price in FCFA.
\begin{enumerate}
\item Find the free‑trade equilibrium quantity of imports.
\item Determine the new domestic price, domestic consumption, domestic production and quantity of imports after the tariff.
\item Calculate the tariff revenue, the consumer surplus loss, the producer surplus gain and the deadweight loss.
\item Discuss briefly the welfare implications for Cameroon.
\end{enumerate}
\end{exercice}

\begin{exercice}[Balance of Payments Classification]
Classify each of the following transactions into the appropriate Balance of Payments account (Current Account: Goods, Services, Primary Income, Secondary Income; Capital Account; Financial Account: Direct Investment, Portfolio Investment, Other Investment, Reserve Assets). Indicate whether each entry is a credit (+) or debit (–) for Cameroon.

\begin{enumerate}
\item Export of 10,000 tonnes of cocoa beans worth 15 billion FCFA.
\item Import of petroleum products worth 20 billion FCFA.
\item Remittances from Cameroonian diaspora: 5 billion FCFA.
\item Interest payments on external debt: 2 billion FCFA.
\item A French firm builds a new factory in Douala (FDI inflow): 30 billion FCFA.
\item The Central Bank sells 10 billion FCFA worth of foreign reserves to support the exchange rate.
\item Tourism receipts from foreign visitors: 3 billion FCFA.
\item Debt forgiveness by a bilateral creditor: 8 billion FCFA.
\end{enumerate}
\end{exercice}

\begin{exercice}[Exchange Rate Calculations]
Assume the following exchange rates:
\begin{itemize}
\item 1 EUR = 655.957 FCFA (fixed peg)
\item 1 USD = 0.90 EUR
\item 1 USD = 600 FCFA (market rate)
\end{itemize}

\begin{enumerate}
\item Calculate the cross rate USD/FCFA implied by the EUR peg.
\item A Cameroonian exporter sells cocoa for 2,000 USD per tonne. How many FCFA does the exporter receive at the official peg? How many at the market rate?
\item If the EUR depreciates by 10\% against the USD while the FCFA/EUR peg remains unchanged, what happens to the FCFA/USD cross rate? Explain the impact on Cameroon's export competitiveness.
\end{enumerate}
\end{exercice}

\courssub{I prepare for the GCE}

\begin{exercice}[Data Response: Terms of Trade and Protectionism — 20 marks]
The table below shows export and import price indices (base 2015=100) for Cameroon for three years.

\begin{center}
\begin{tabular}{|c|c|c|}
\hline
Year & Export Price Index & Import Price Index \\
\hline
2020 & 110 & 130 \\
2021 & 115 & 140 \\
2022 & 118 & 155 \\
\hline
\end{tabular}
\end{center}

\begin{enumerate}
\item Compute the Terms of Trade index for each year. [4 marks]
\item Calculate the percentage change in the Terms of Trade between 2020 and 2022. [2 marks]
\item Explain two possible reasons for the deterioration in Cameroon's Terms of Trade over this period. [6 marks]
\item Discuss whether imposing tariffs on imported manufactured goods could improve Cameroon's Terms of Trade. Use a diagram to support your answer. [8 marks]
\end{enumerate}
\end{exercice}

\begin{exercice}[Essay: Balance of Payments Deficit in Cameroon — 25 marks]
"Persistent balance of payments deficits in Cameroon are mainly structural and cannot be corrected by expenditure‑switching policies alone."

\begin{enumerate}
\item Identify and explain four structural causes of a persistent current account deficit in Cameroon. [10 marks]
\item Evaluate the effectiveness of the following measures in correcting the deficit, given that the FCFA is pegged to the Euro:
\begin{itemize}
\item Expenditure‑reducing fiscal and monetary policies.
\item Expenditure‑switching policies (devaluation is not an option).
\item Supply‑side policies to boost export competitiveness.
\item Seeking assistance from the IMF and World Bank.
\end{itemize}
[15 marks]
\end{enumerate}
\end{exercice}

\begin{exercice}[Essay: International Institutions and Regional Groupings — 20 marks]
\begin{enumerate}
\item Describe the main functions of the International Monetary Fund (IMF) and the World Bank. [8 marks]
\item Assess the role of the IMF and the World Bank in Cameroon's economic adjustment programmes since the 1980s. [12 marks]
\end{enumerate}
\end{exercice}

\newpage

\coursec{Solutions to the exercises}

\begin{corrige}[Exercise 1 — MCQ: Basis of Trade]
The correct answer is \textbf{b)}.

\textbf{Justification:} The principle of \cle{comparative advantage}, developed by David Ricardo (1817), states that a country gains from specialising in and exporting the goods for which its \cle{opportunity cost} of production is lower than that of its trading partners, even if it has an absolute advantage (or disadvantage) in \emph{all} goods.

Why the other options are wrong:
\begin{itemize}
\item \textbf{a)} is false: absolute advantage is neither necessary nor sufficient for beneficial trade; a country with no absolute advantage anywhere still gains by specialising where its \emph{relative} (opportunity) cost is lowest.
\item \textbf{c)} is false: trade is beneficial precisely because production possibilities (and therefore opportunity costs) \emph{differ} between countries.
\item \textbf{d)} is false: comparative advantage concerns relative efficiency (opportunity cost), not the absolute endowment of natural resources.
\end{itemize}
\end{corrige}

\begin{corrige}[Exercise 2 — True/False with Justification]
\begin{enumerate}
\item \textbf{False.} The domestic price rises by the full tariff only if the country is a \emph{price taker} facing a perfectly elastic world supply (small-country case). If the importing country is large, part of the tariff is absorbed by foreign exporters through a lower world price, so the domestic price rises by \emph{less} than the tariff. Even in the small-country case, the price rises by the full amount only if the good continues to be imported.
\item \textbf{False.} The current account records trade in \textbf{goods} \emph{and} \textbf{services}, plus \textbf{primary income} (investment income, compensation of employees) and \textbf{secondary income} (current transfers such as remittances and foreign aid grants). Trade in goods alone is only the ``visible balance''.
\item \textbf{True.} The terms of trade are $TOT = \dfrac{P_X}{P_M}\times 100$. An improvement (rise) means export prices have risen relative to import prices, so a given volume of exports now purchases a larger volume of imports — the country's import purchasing power has increased.
\item \textbf{True.} To defend a fixed (pegged) rate, the central bank must stand ready to buy its own currency (using reserves) when there is downward pressure and to sell it (accumulating reserves) when there is upward pressure. Holding adequate \cle{foreign exchange reserves} is therefore essential to maintain the peg — as the BEAC does for the FCFA within the CEMAC zone.
\end{enumerate}
\end{corrige}

\begin{corrige}[Exercise 3 — Definitions]
\begin{enumerate}
\item \textbf{Terms of Trade (TOT):} the ratio of a country's export price index to its import price index, expressed as a percentage:
\[ TOT = \frac{\text{Index of export prices}}{\text{Index of import prices}}\times 100 \]
It measures the volume of imports obtainable per unit of exports. A rise is an ``improvement'', a fall a ``deterioration''.
\item \textbf{Balance of Payments overall balance:} the net sum of the current account, the capital account and the financial account (excluding reserve assets), together with net errors and omissions. It equals the change in official reserves: a positive overall balance means reserves are accumulating; a negative one means the country is financing a gap by drawing down reserves or borrowing.
\item \textbf{Devaluation:} a deliberate \emph{official} reduction in the value of the domestic currency against a foreign currency (or gold/SDR) decided by the monetary authorities under a \emph{fixed} exchange rate regime (e.g.\ the 1994 devaluation of the FCFA from 50 to 100 FCFA per French franc). \textbf{Depreciation}, by contrast, is a fall in the currency's value brought about by \emph{market forces} (demand and supply) under a \emph{floating} regime.
\item \textbf{Infant industry argument:} the claim that new domestic industries, which initially face higher costs than established foreign competitors, should be granted \emph{temporary} protection (tariffs, quotas, subsidies) until they grow, exploit economies of scale and ``learning by doing'', and become internationally competitive — after which protection should be removed.
\end{enumerate}
\end{corrige}

\begin{corrige}[Exercise 4 — Direct Calculations (Terms of Trade)]
\begin{enumerate}
\item Terms of Trade index for 2023:
\[ TOT_{2023} = \frac{P_X}{P_M}\times 100 = \frac{120}{150}\times 100 = \boxed{80} \]
\item Percentage change between 2022 and 2023:
\[ \%\Delta TOT = \frac{80 - 95}{95}\times 100 = \frac{-15}{95}\times 100 \approx \boxed{-15.8\%} \]
\item \textbf{Interpretation:} the terms of trade have \emph{deteriorated} sharply. In 2023, one unit of Cameroon's cocoa exports buys only 80\% of the imports it could buy in the base year (2015). Compared with 2022, the import-purchasing power of exports has fallen by about 15.8\%: Cameroon must export roughly 15.8\% more cocoa than in 2022 to finance the same volume of machinery imports. This squeezes the trade balance and reduces real national income available for investment goods.
\end{enumerate}
\end{corrige}

\begin{corrige}[Exercise 5 — Absolute and Comparative Advantage: Cameroon and Guinea]
\textbf{Data (output per worker per day):} Cameroon: 4 t cocoa or 2 t aluminium; Guinea: 1 t cocoa or 3 t aluminium.

\begin{enumerate}
\item \textbf{Absolute advantage:} Cameroon produces more cocoa per worker ($4>1$): absolute advantage in \textbf{cocoa}. Guinea produces more aluminium ($3>2$): absolute advantage in \textbf{aluminium}.

\item \textbf{Opportunity cost of 1 tonne of cocoa:}
\begin{itemize}
\item Cameroon: one worker-day yields 4 t cocoa \emph{or} 2 t aluminium, so 1 t cocoa costs $\dfrac{2}{4} = \textbf{0.5 t of aluminium}$.
\item Guinea: one worker-day yields 1 t cocoa \emph{or} 3 t aluminium, so 1 t cocoa costs $\dfrac{3}{1} = \textbf{3 t of aluminium}$.
\end{itemize}

\item \textbf{Comparative advantage:} cocoa is relatively cheaper in Cameroon ($0.5 < 3$ tonnes of aluminium forgone), so \textbf{Cameroon has the comparative advantage in cocoa}. Symmetrically, 1 t of aluminium costs 2 t of cocoa in Cameroon but only $\tfrac{1}{3}$ t of cocoa in Guinea: \textbf{Guinea has the comparative advantage in aluminium}.

\item \textbf{Mutually beneficial exchange ratio:} any ratio between the two opportunity costs benefits both:
\[ 0.5 \text{ t Al} < 1 \text{ t cocoa} < 3 \text{ t Al} \]
Take, for example, \textbf{1 tonne of cocoa = 1 tonne of aluminium}.

\textbf{Gains from trade (per worker-day of resources shifted):}
\begin{itemize}
\item \emph{Cameroon:} specialising, it produces 4 t cocoa. It exports 1 t cocoa and receives 1 t aluminium. Producing that 1 t of aluminium itself would have required giving up 2 t of cocoa (its opportunity cost of aluminium is 2 t cocoa). Through trade it obtains 1 t aluminium for only 1 t cocoa instead of 2 t cocoa — a gain of 1 t of cocoa equivalent.
\item \emph{Guinea:} it exports 1 t aluminium (which costs it only $\tfrac{1}{3}$ t cocoa to produce) and receives 1 t cocoa, which would have cost it 3 t of aluminium to produce domestically. It gains $3 - 1 = 2$ t of aluminium equivalent.
\end{itemize}
Both countries consume beyond their own production possibility frontiers: trade is mutually beneficial.
\end{enumerate}
\end{corrige}

\begin{corrige}[Exercise 6 — Effects of a Tariff]
\textbf{Given:} $Q_d = 5000 - 0.2P$, $Q_s = 1000 + 0.1P$, world price $P_w = 10\,000$ FCFA, specific tariff $t = 500$ FCFA/tonne.

\begin{enumerate}
\item \textbf{Free-trade imports} ($P = P_w = 10\,000$):
\begin{align*}
Q_d &= 5000 - 0.2(10\,000) = 5000 - 2000 = 3000 \text{ t}\\
Q_s &= 1000 + 0.1(10\,000) = 1000 + 1000 = 2000 \text{ t}\\
M_0 &= Q_d - Q_s = \boxed{1000 \text{ tonnes}}
\end{align*}

\item \textbf{After the tariff} the domestic price rises to $P' = 10\,000 + 500 = \boxed{10\,500 \text{ FCFA}}$:
\begin{align*}
Q_d' &= 5000 - 0.2(10\,500) = 5000 - 2100 = \boxed{2900 \text{ t}}\\
Q_s' &= 1000 + 0.1(10\,500) = 1000 + 1050 = \boxed{2050 \text{ t}}\\
M_1 &= 2900 - 2050 = \boxed{850 \text{ tonnes}}
\end{align*}

\item \textbf{Welfare calculations:}
\begin{itemize}
\item \emph{Tariff revenue:} $T = t \times M_1 = 500 \times 850 = \boxed{425\,000 \text{ FCFA}}$.
\item \emph{Consumer surplus loss} (trapezium between old and new price under the demand curve):
\[ \Delta CS = -(P'-P_w)\times\frac{Q_d + Q_d'}{2} = -500\times\frac{3000+2900}{2} = \boxed{-1\,475\,000 \text{ FCFA}} \]
\item \emph{Producer surplus gain} (trapezium above the supply curve):
\[ \Delta PS = +(P'-P_w)\times\frac{Q_s + Q_s'}{2} = 500\times\frac{2000+2050}{2} = \boxed{+1\,012\,500 \text{ FCFA}} \]
\item \emph{Deadweight loss} (net national loss):
\[ DWL = |\Delta CS| - \Delta PS - T = 1\,475\,000 - 1\,012\,500 - 425\,000 = \boxed{37\,500 \text{ FCFA}} \]
(Equivalently: production distortion $\tfrac{1}{2}\times 500 \times 50 = 12\,500$ plus consumption distortion $\tfrac{1}{2}\times 500 \times 100 = 25\,000$.)
\end{itemize}

\item \textbf{Welfare discussion:} consumers lose (higher price, lower consumption), domestic producers gain (higher price, expanded output), and the government gains revenue. But the gains to producers and government ($1\,012\,500 + 425\,000 = 1\,437\,500$) fall short of the consumer loss ($1\,475\,000$): the nation as a whole suffers a net welfare (deadweight) loss of 37\,500 FCFA, because resources are drawn into higher-cost domestic production and mutually beneficial consumption is forgone. The tariff may still be defended on non-efficiency grounds (food security, infant industry, revenue), but it is economically inefficient.
\end{enumerate}
\end{corrige}

\begin{corrige}[Exercise 7 — Balance of Payments Classification]
\begin{center}
\begin{tabularx}{\linewidth}{|c|X|l|c|}
\hline
\rowcolor{popblueL} N\textdegree & Transaction & BOP account & Sign \\
\hline
1 & Cocoa exports, 15 bn FCFA & Current Account — Goods & Credit ($+$) \\
\hline
2 & Petroleum imports, 20 bn FCFA & Current Account — Goods & Debit ($-$) \\
\hline
3 & Diaspora remittances, 5 bn FCFA & Current Account — Secondary Income & Credit ($+$) \\
\hline
4 & Interest on external debt, 2 bn FCFA & Current Account — Primary Income & Debit ($-$) \\
\hline
5 & French firm's factory in Douala, 30 bn FCFA & Financial Account — Direct Investment & Credit ($+$) \\
\hline
6 & Central Bank sells reserves, 10 bn FCFA & Financial Account — Reserve Assets & Credit ($+$)\textsuperscript{*} \\
\hline
7 & Tourism receipts, 3 bn FCFA & Current Account — Services & Credit ($+$) \\
\hline
8 & Debt forgiveness, 8 bn FCFA & Capital Account (debt forgiveness transfer) & Credit ($+$) \\
\hline
\end{tabularx}
\end{center}
\textsuperscript{*}A \emph{decrease} in reserve assets is recorded as a credit (an inflow of foreign exchange into the domestic economy); an accumulation of reserves would be a debit.

\textbf{Net effect check:} Current account $= 15 - 20 + 5 - 2 + 3 = +1$ bn; capital account $= +8$ bn; financial account (excl.\ reserves) $= +30$ bn; reserves $+10$ bn (drawdown). The accounts balance overall once the counterpart entries (e.g.\ payment flows) are included.
\end{corrige}

\begin{corrige}[Exercise 8 — Exchange Rate Calculations]
\begin{enumerate}
\item \textbf{Cross rate implied by the EUR peg:} since 1 USD $= 0.90$ EUR and 1 EUR $= 655.957$ FCFA:
\[ 1 \text{ USD} = 0.90 \times 655.957 = \boxed{590.36 \text{ FCFA (approx.)}} \]

\item \textbf{Exporter's receipts for 2,000 USD:}
\begin{itemize}
\item At the official peg: $2000 \times 590.36 = \boxed{1\,180\,720 \text{ FCFA}}$ (approx.).
\item At the market rate: $2000 \times 600 = \boxed{1\,200\,000 \text{ FCFA}}$.
\end{itemize}
The market rate yields about 19\,280 FCFA more per tonne; the gap reflects the fact that the market implicitly values the USD slightly higher against the FCFA than the EUR-peg cross rate.

\item \textbf{EUR depreciates 10\% against the USD, peg unchanged:} the new rate is 1 USD $= 0.90 \times 1.10 = 0.99$ EUR (the USD buys 10\% more euros). Hence:
\[ 1 \text{ USD} = 0.99 \times 655.957 \approx \boxed{649.4 \text{ FCFA}} \]
The FCFA \emph{depreciates} against the USD (more FCFA per dollar). \textbf{Impact on competitiveness:} Cameroon's exports priced in USD (cocoa, oil, timber, cotton) become cheaper for dollar-area buyers and earn more FCFA per unit sold, improving price competitiveness and exporters' FCFA revenues. Conversely, imports invoiced in USD become dearer in FCFA, which raises import costs — a mixed effect typical of a pegged currency whose anchor currency moves against third currencies.
\end{enumerate}
\end{corrige}

\begin{corrige}[Exercise 9 — Data Response: Terms of Trade and Protectionism (20 marks)]
\textbf{(a) TOT indices} $TOT = \dfrac{P_X}{P_M}\times 100$ \textbf{[4 marks — 1 per correct value, 1 for formula]}
\begin{itemize}
\item 2020: $\dfrac{110}{130}\times 100 \approx \mathbf{84.6}$
\item 2021: $\dfrac{115}{140}\times 100 \approx \mathbf{82.1}$
\item 2022: $\dfrac{118}{155}\times 100 \approx \mathbf{76.1}$
\end{itemize}

\textbf{(b) Percentage change 2020--2022 [2 marks]}
\[ \frac{76.1 - 84.6}{84.6}\times 100 \approx \mathbf{-10.0\%} \]
(1 mark method, 1 mark accuracy.)

\textbf{(c) Two reasons for the deterioration [6 marks — 3 each: identification 1, explanation 2]}
\begin{itemize}
\item \emph{Fall in world prices of primary exports / rise in export supply:} Cameroon's exports are concentrated in primary commodities (cocoa, coffee, oil, timber) whose world prices are volatile and trend weakly because of low income elasticity of demand, inelastic supply responses and competition from other producers; export prices rose only modestly (110 to 118).
\item \emph{Rise in import prices of manufactures and energy:} manufactured goods, machinery and fuel embody higher value added and their prices rose faster (130 to 155); global inflation, shipping costs and currency effects raise the import bill of a country dependent on imported capital and intermediate goods.
\end{itemize}
(Also acceptable: declining productivity/quality of export crops; Prebisch--Singer hypothesis of secular deterioration for primary exporters; exchange-rate pass-through.)

\textbf{(d) Can tariffs on manufactured imports improve the TOT? [8 marks]}

\emph{Analysis:} a tariff improves the terms of trade only if the country is \textbf{large enough} to influence world prices: by reducing import demand, it forces foreign suppliers to lower their prices, so $P_M$ falls relative to $P_X$.

\begin{popfigure}
\begin{center}
\begin{tikzpicture}[scale=0.85]
\draw[->, thick] (0,0) -- (9,0) node[below left]{\footnotesize Q};
\draw[->, thick] (0,0) -- (0,6) node[above left]{\footnotesize P};
\draw[popblue, very thick, domain=1:8] plot (\x, {5.5-0.55*\x}) node[right]{\footnotesize $D_M$};
\draw[poporange, very thick, domain=1:8] plot (\x, {1.2+0.45*\x}) node[right]{\footnotesize $S_M$};
\draw[poporange, dashed, thick] (1,2.6) -- (8,2.6) node[right]{\footnotesize $S_M$ elastic};
\draw[dashed, gray] (4.05,0) node[below]{\footnotesize $M_0$} -- (4.05,3.27);
\fill[popink] (4.05,3.27) circle (2pt);
\end{tikzpicture}
\end{center}
\legende{Import market: with upward-sloping foreign supply, a tariff lowers the world price (TOT gain); with perfectly elastic supply (small country), no TOT gain.}
\end{popfigure}

\emph{Evaluation:}
\begin{itemize}
\item Cameroon is a \textbf{small open economy}: it is a price taker for manufactures, so foreign supply is effectively perfectly elastic and the tariff raises the domestic price without lowering world prices — \textbf{no TOT gain}.
\item Tariffs create deadweight losses, invite \textbf{retaliation}, may breach CEMAC/WTO/AfCFTA commitments, and protect inefficiency.
\item However, a TOT improvement is theoretically possible for a large buyer, and tariffs may be justified on infant-industry or revenue grounds even without a TOT gain.
\end{itemize}
\emph{Mark scheme:} definition/mechanism of TOT effect of tariff (2), correct diagram (2), small-country application to Cameroon (2), evaluation/judgement (2).
\end{corrige}

\begin{corrige}[Exercise 10 — Essay: Balance of Payments Deficit in Cameroon (25 marks)]
\textbf{(a) Four structural causes of a persistent current account deficit [10 marks — 2.5 each]}
\begin{enumerate}
\item \textbf{Narrow, commodity-concentrated export base:} exports depend on a few primary products (crude oil, cocoa, coffee, timber, cotton) with volatile world prices and low value added; export earnings stagnate when prices fall.
\item \textbf{High import dependence:} the industrial base is shallow, so capital goods, intermediate inputs, refined fuel (despite crude exports, given limited refining capacity after SONARA's difficulties) and even food must be imported — imports are income-elastic and price-inelastic.
\item \textbf{Low productivity and weak competitiveness:} poor infrastructure (energy, roads, port congestion), high factor costs and skills gaps raise production costs, limiting diversification into manufactured exports.
\item \textbf{Structural debt-service and income outflows:} interest payments on external debt and profit/dividend repatriation by foreign firms create persistent primary-income debits; population growth raises import demand faster than export capacity.
\end{enumerate}

\textbf{(b) Evaluation of corrective measures under the Euro peg [15 marks]}

\emph{Expenditure-reducing policies (fiscal consolidation, tight monetary policy via BEAC):} lower domestic income and inflation reduce import demand and can make exports relatively cheaper (internal devaluation). \emph{But:} they contract output and employment, are socially costly, and work slowly; within CEMAC, monetary policy is regional, limiting national discretion. Effective as a short-run stabiliser, not a structural cure.

\emph{Expenditure-switching policies with no devaluation option:} since the FCFA is pegged to the euro, nominal devaluation requires a CEMAC-wide decision (as in 1994). Switching must therefore rely on tariffs, quotas, import-substitution incentives and subsidies. These can redirect demand to home goods, \emph{but} risk retaliation, WTO/CEMAC/AfCFTA incompatibility, smuggling, and protection of inefficiency; they do not raise export \emph{supply} capacity.

\emph{Supply-side policies:} investment in energy, transport and port infrastructure; education and vocational training; support to agro-processing and local value addition (e.g.\ processing cocoa locally); improving the business climate. These raise productivity, diversify exports and reduce import dependence — they address the \textbf{structural} causes directly. \emph{But:} they are slow (long gestation), require financing and institutional quality, and their payoff lies beyond the short run.

\emph{IMF/World Bank assistance:} provides foreign exchange, policy credibility and structural reform frameworks (Extended Credit Facility-type programmes), easing the financing constraint and supporting reforms. \emph{But:} conditionality (austerity, privatisation) can be socially painful and politically sensitive; assistance treats symptoms unless domestic reforms follow through.

\textbf{Conclusion:} the statement is largely correct. Because the deficit is structural (concentrated exports, import dependence, low competitiveness), expenditure-switching alone — especially without the devaluation instrument — cannot restore equilibrium. A durable correction requires supply-side transformation, supported by prudent demand management and, where necessary, external assistance.

\emph{Examiner expectations:} clear structure (introduction defining BOP deficit; developed paragraphs; reasoned conclusion); accurate use of concepts (expenditure-reducing vs switching); Cameroon/CEMAC context (peg, BEAC, commodity exports); balanced evaluation in each part. Indicative split: (a) 10, (b) 15 with roughly 3--4 marks per measure including evaluation.
\end{corrige}

\begin{corrige}[Exercise 11 — Essay: International Institutions and Regional Groupings (20 marks)]
\textbf{(a) Main functions of the IMF and the World Bank [8 marks — 4 each]}

\emph{International Monetary Fund (IMF):}
\begin{itemize}
\item Surveillance of members' economic and exchange-rate policies.
\item Providing short- and medium-term \textbf{balance-of-payments financing} (stand-by and extended arrangements, concessional facilities for low-income countries).
\item Promoting international monetary cooperation and orderly exchange arrangements.
\item Technical assistance and capacity building (statistics, central banking, public finance).
\end{itemize}

\emph{World Bank (IBRD/IDA):}
\begin{itemize}
\item Long-term \textbf{development financing} for projects and programmes (infrastructure, education, health, agriculture).
\item Poverty reduction and support for structural reforms and institution building.
\item Concessional lending and grants to low-income countries (IDA).
\item Technical expertise, research and policy advice on development.
\end{itemize}
Key distinction: the IMF deals with \emph{macroeconomic stabilisation and BOP support}; the World Bank with \emph{long-run development and poverty reduction}.

\textbf{(b) Role of the IMF and World Bank in Cameroon's adjustment since the 1980s [12 marks]}

\emph{Context:} from the mid-1980s Cameroon suffered a severe crisis (fall in oil, cocoa and coffee prices, overvalued real exchange rate, fiscal deficits). From the late 1980s it adopted successive \textbf{Structural Adjustment Programmes (SAPs)} designed and financed with the IMF and World Bank.

\emph{Positive contributions:}
\begin{itemize}
\item Restored macroeconomic stability: fiscal consolidation, monetary discipline within CEMAC, and the 1994 FCFA devaluation (supported by the institutions) restored competitiveness and growth in the late 1990s.
\item Debt relief: Cameroon reached the completion point of the \textbf{HIPC initiative} (mid-2000s), obtaining substantial debt cancellation that freed budgetary resources for social spending.
\item Financing and reforms in public financial management, privatisation of inefficient parastatals, and sectoral projects (roads, health, education).
\item Technical assistance and policy credibility that helped mobilise other donors.
\end{itemize}

\emph{Criticisms and limits:}
\begin{itemize}
\item Social costs: cuts in public wages and employment, reduced subsidies and cost-recovery in health/education hurt the poor and fuelled emigration of skilled workers in the 1990s.
\item ``One-size-fits-all'' conditionality sometimes ignored local institutional realities; implementation was uneven and ownership weak.
\item Retrenchment of the state weakened public services; privatisation results were mixed.
\item Despite adjustment, diversification remained limited: the economy stayed dependent on a few commodity exports, so vulnerability persisted.
\end{itemize}

\emph{Balanced judgement:} IMF/World Bank programmes helped Cameroon avert collapse, stabilise the macroeconomy and obtain debt relief, but their social costs and limited success in transforming the productive structure explain the persistent debate about their design. Later programmes emphasise growth, governance and social protection more explicitly.

\emph{Mark scheme (12):} context of the crisis (2); stabilisation/devaluation/debt-relief achievements (4); criticisms with Cameroonian evidence (4); balanced conclusion (2). Examiner expects accurate chronology in broad terms (no invented dates/figures), correct naming of instruments (SAPs, HIPC), and genuine two-sided assessment.
\end{corrige}

\newpage

\coursec{Summary sheet}

\begin{popfigure}
\begin{tikzpicture}[
    mindmap,
    concept color=popblue,
    level 1 concept/.append style={level distance=4.5cm, sibling angle=360/7},
    level 2 concept/.append style={level distance=3cm},
    every node/.style={font=\small, text width=2.5cm, align=center},
    grow cyclic,
    popblue/.style={concept color=popblue},
    poporange/.style={concept color=poporange},
    popgreen/.style={concept color=popgreen},
    poppurple/.style={concept color=poppurple},
    poppink/.style={concept color=poppink},
    popgold/.style={concept color=popgold},
    popdark/.style={concept color=popdark}
]
\node[concept, align=center]{International\\ Economics}
    child[popblue] { node[concept, align=center]{Meaning \&\\ Theories}
        child { node {Absolute Advantage (Smith)} }
        child { node {Comparative Advantage (Ricardo)} }
        child { node {Heckscher-Ohlin (Factor Endowments)} }
    }
    child[poporange] { node[concept, align=center]{Gains from Trade\\ \& Specialisation}
        child { node {Opportunity Cost} }
        child { node {Production Possibility Frontier} }
        child { node {Terms of Trade (TOT)} }
    }
    child[popgreen] { node[concept, align=center]{Trade\\ Restrictions}
        child { node {Tariffs (Specific/Ad valorem)} }
        child { node {Quotas \& VERs} }
        child { node {Non-tariff Barriers} }
        child { node {Arguments For/Against} }
    }
    child[poppurple] { node[concept, align=center]{Balance of\\ Payments}
        child { node {Current Account (Goods, Services, Income, Transfers)} }
        child { node {Capital \& Financial Account} }
        child { node {Double Entry Recording} }
        child { node {Disequilibrium: Deficit/Surplus} }
    }
    child[poppink] { node[concept, align=center]{Correction\\ Policies}
        child { node {Expenditure Switching (Deval/Deprec)} }
        child { node {Expenditure Reducing (Fiscal/Monetary)} }
        child { node {Direct Controls} }
        child { node {Marshall-Lerner Condition} }
    }
    child[popgold] { node[concept, align=center]{Institutions\\ \& Groupings}
        child { node {IMF, World Bank, WTO} }
        child { node {CEMAC, ECCAS, AfCFTA} }
        child { node {CFA Franc Zone} }
    }
    child[popdark] { node[concept, align=center]{Forex Market\\ \& Exchange Rates}
        child { node {Demand/Supply of FX} }
        child { node {Regimes: Fixed, Floating, Managed} }
        child { node {Appreciation/Depreciation} }
        child { node {Purchasing Power Parity} }
    };
\end{tikzpicture}
\legende{Chapter Mind Map: International Economics (ECO06)}
\end{popfigure}

\begin{bilan}

\textbf{Key Definitions}
\begin{itemize}
    \item \cle{International Trade}: Exchange of goods, services and capital across national borders.
    \item \cle{Comparative Advantage}: A country produces a good at lower \cle{opportunity cost} than another (Ricardo).
    \item \cle{Terms of Trade (TOT)}: $TOT = \frac{\text{Index of Export Prices}}{\text{Index of Import Prices}} \times 100$. Improvement = rising index.
    \item \cle{Balance of Payments (BoP)}: Systematic record of all economic transactions between residents of a country and the rest of the world in a given period (Double-entry: Credit +, Debit -).
    \item \cle{Current Account}: Trade balance + Net primary income + Net secondary income (transfers).
    \item \cle{Capital/Financial Account}: Capital transfers + Acquisition/disposal of non-produced assets + Financial assets/liabilities (FDI, Portfolio, Reserves).
    \item \cle{Exchange Rate}: Price of one currency in terms of another. \cle{Appreciation} (rise in value), \cle{Depreciation} (fall in value, floating); \cle{Revaluation}/\cle{Devaluation} (fixed regime).
    \item \cle{Marshall-Lerner Condition}: Devaluation improves BoP only if $PED_x + PED_m > 1$ (sum of price elasticities of demand for exports and imports).
\end{itemize}

\textbf{Laws, Formulas \& Theorems to Know}
\begin{itemize}
    \item \textbf{Absolute Advantage (Adam Smith)}: Specialise if $a_{LC} < a_{LC}^*$ (lower labour input per unit).
    \item \textbf{Comparative Advantage (David Ricardo)}: Specialise based on \emph{opportunity cost ratios}. Gains from trade exist if $\frac{a_{LC}}{a_{LW}} \neq \frac{a_{LC}^*}{a_{LW}^*}$.
    \item \textbf{Heckscher-Ohlin Theorem}: Countries export goods intensive in their abundant factors (Capital-abundant $\to$ capital-intensive goods).
    \item \textbf{Stolper-Samuelson Theorem}: Trade raises real return to abundant factor, lowers return to scarce factor.
    \item \textbf{TOT Formula}: $TOT = \frac{P_x}{P_m} \times 100$. \cle{Prebisch-Singer Hypothesis}: Long-run decline in TOT for primary commodity exporters (e.g. Cameroon cocoa, coffee, oil).
    \item \textbf{BoP Identity}: $CA + KA + FA + \text{Net Errors \& Omissions} = 0$. $\Delta \text{Reserves} = -(CA + KA + FA)$.
    \item \textbf{Elasticities Approach}: $\Delta B \approx \eta_x \Delta E + \eta_m \Delta E$ (simplified).
    \item \textbf{PPP (Absolute)}: $S = \frac{P}{P^*}$. \textbf{PPP (Relative)}: $\%\Delta S \approx \pi - \pi^*$.
\end{itemize}

\textbf{Methods (Step-by-Step)}
\begin{enumerate}
    \item \textbf{Comparative Advantage Numerical}: 
    \begin{enumerate}
        \item Calculate Opportunity Cost (OC) for each good in each country (OC of Good X = Output Y foregone / Output X gained).
        \item Compare OCs. Lower OC $\to$ Comparative Advantage.
        \item Determine Specialisation pattern.
        \item Find \cle{Terms of Trade limits}: Must lie between the two domestic opportunity cost ratios.
        \item Calculate Consumption Possibilities with trade; show gains (consumption outside PPF).
    \end{enumerate}
    \item \textbf{BoP Recording (Double Entry)}:
    \begin{enumerate}
        \item Identify transaction (Export, Import, Remittance, Loan, FDI).
        \item Classify: Current vs Capital/Financial.
        \item Apply sign: Credit (+) = Inflow (Export, Loan received); Debit (-) = Outflow (Import, Loan repayment).
        \item Record two entries (e.g. Export cocoa: Credit Current Account/Goods + Credit Financial Account/Currency/Reserves).
    \end{enumerate}
    \item \textbf{Exchange Rate Diagram (Forex Market)}:
    \begin{enumerate}
        \item Axes: $Y$ = Exchange Rate (FCFA/\$), $X$ = Qty of Foreign Currency (\$).
        \item Demand ($D_{\$}$): Downward sloping (imports, capital outflow).
        \item Supply ($S_{\$}$): Upward sloping (exports, capital inflow).
        \item Equilibrium $E_0$. Shifts: $\uparrow$ Exports $\to$ $\uparrow S_{\$}$ $\to$ Appreciation (FCFA strengthens).
    \end{enumerate}
    \item \textbf{Devaluation Analysis (J-Curve)}:
    \begin{enumerate}
        \item Short run: Inelastic demand $\to$ Import bill $\uparrow$, Export revenue $\downarrow$ $\to$ BoP worsens.
        \item Long run: Elasticities rise $\to$ Volumes adjust $\to$ BoP improves (if M-L holds).
    \end{enumerate}
\end{enumerate}

\textbf{Common Mistakes \& Traps}
\begin{itemize}
    \item \textbf{Confusing Absolute vs Comparative Advantage}: Absolute = productivity; Comparative = opportunity cost. Trade based on Comparative.
    \item \textbf{TOT Direction}: Rising TOT = \textbf{Improvement} (favourable). Falling TOT = \textbf{Deterioration} (unfavourable). Do not confuse with "better/worse" loosely.
    \item \textbf{BoP Signs}: Credits (+) are sources of forex (Exports, Foreign Investment In). Debits (-) are uses of forex (Imports, Capital Outflow). \emph{Increase in Reserves is a Debit (outflow of domestic asset)}.
    \item \textbf{Current vs Capital Account}: Remittances (diaspora) = Current Transfers (Current Account). Buying shares in SONARA by French firm = FDI (Financial Account).
    \item \textbf{Devaluation vs Depreciation}: Devaluation = Policy decision (Fixed regime). Depreciation = Market force (Floating regime).
    \item \textbf{Marshall-Lerner}: Must state condition $PED_x + PED_m > 1$. If inelastic, devaluation worsens deficit.
    \item \textbf{CFA Franc}: Fixed parity with Euro (1 EUR = 655.957 FCFA). Guaranteed convertibility by French Treasury. Not a floating rate.
\end{itemize}

\textbf{GCE Advanced Level Tips}
\begin{itemize}
    \item \textbf{Essay Structure}: Define $\to$ Explain Mechanism (Diagram/Math) $\to$ Apply to Cameroon/CEMAC (Cocoa, Oil, CFA, AfCFTA) $\to$ Evaluate (Limits, Time lags, Elasticities).
    \item \textbf{Diagrams are Marks}: Draw PPF (Concave), Forex Market (S/D), Tariff Diagram (Domestic S/D + World Price), J-Curve. Label axes, curves, equilibria ($E_0, E_1$), areas (Consumer Surplus, Govt Revenue, DWL).
    \item \textbf{Calculation Questions}: Show formula $\to$ Substitution $\to$ Answer with units. TOT index, Opportunity Cost ratios, BoP balancing item.
    \item \textbf{Contextualise}: Mention \cle{AfCFTA} (intra-African trade), \cle{CEMAC} (regional integration), \cle{SONARA} (oil refining/imports), \cle{Cocoa/Coffee} (primary exports, TOT vulnerability), \cle{Diaspora Remittances} (major Current Account credit).
    \item \textbf{Policy Mix}: For BoP deficit under Fixed Rate (CFA zone): Expenditure Reducing (Fiscal austerity) + Structural Reforms (Supply side) since Expenditure Switching (Deval) is constrained by parity.
    \item \textbf{Command Words}: "Distinguish" $\to$ Table/Contrast. "Analyse" $\to$ Cause/Effect chain. "Evaluate" $\to$ Pros/Cons + Judgement. "Calculate" $\to$ Working essential.
\end{itemize}

\end{bilan}

\findecours

\end{document}
